\documentclass[11pt,a4paper]{article}
\usepackage[margin=2.4cm]{geometry}
\usepackage{amsmath,amssymb,amsthm,mathtools,bm}
\usepackage{booktabs,tabularx,array,longtable}
\usepackage{xcolor}
\usepackage[T1]{fontenc}
\usepackage{lmodern}
\usepackage{microtype}
\usepackage[colorlinks=true,linkcolor=blue!50!black,citecolor=blue!50!black,urlcolor=blue!50!black]{hyperref}

\newtheorem{proposition}{Proposition}
\theoremstyle{definition}
\newtheorem{definition}{Definition}
\theoremstyle{remark}
\newtheorem{remark}{Remark}

\newcommand{\R}{\mathbb{R}}
\newcommand{\E}{\mathbb{E}}
\newcommand{\Var}{\operatorname{Var}}
\newcommand{\Pois}{\operatorname{Pois}}
\newcommand{\Cat}{\operatorname{Cat}}
\newcommand{\Dir}{\operatorname{Dir}}
\newcommand{\DP}{\operatorname{DP}}
\newcommand{\Ga}{\operatorname{Gamma}}
\newcommand{\Mult}{\operatorname{Mult}}
\newcommand{\LKJ}{\operatorname{LKJ}}
\newcommand{\TN}{\operatorname{TN}}

\title{relsynth: Statistical Guarantees and Checks for Relational Synthetic Data\\[2pt]
\large Mixed types, star and snowflake schemas, self-referencing hierarchies, temporal facts and text}
\author{Data Synthesisation project}
\date{5 October 2026; additions up to v1.0, 19 October 2026\\[2pt]\small Builds on \texttt{algorithms.pdf} (layer 1). Code: \texttt{src/relsynth}; checks: \texttt{legacy/checks/run\_checks.py} (C00--C78) and \texttt{tests/}}

\begin{document}
\maketitle

\begin{abstract}
Layer 1 fixed the methods and their complexity. This document states, for the most common synthesis situations, what the implemented model guarantees and proves it, then names the numerical check that verifies each statement on fake data with known ground truth (checks C00--C68; Table~\ref{tab:results}). Version 0.3 (Section~\ref{sec:v03}) adds what the first real-database run required: grains (unique keys made of foreign keys), structural zeros, computed columns (lookups, derives, aggregates), nullable keys, a rank-likelihood copula and posterior predictive draws. The situations are: one table with mixed column types and missing values; star schemas; snowflake chains; self-referencing hierarchies; time-stamped facts; free text; business rules; and an enterprise schema combining all of them. Section~\ref{sec:real} gives the checks that need no ground truth, which are the ones the local scripts run on real data, and Section~\ref{sec:doctor} the diagnostics loop that turns their output into spec fixes. Sections~\ref{sec:v04}--\ref{sec:v06} add what the next three real databases required. Section~\ref{sec:v07} makes every synthetic dataset a draw from the posterior predictive distribution and adds rules that hold on most rows, group effects, link affinity, participation hurdles, cannot-link strata, type 2 slowly changing dimensions, unknown members and bridges (checks C57--C68). Sections~\ref{sec:v08} and \ref{sec:v09} add identification from the data, the exported model and a review of pooling and multiplicity (C69--C78). Section~\ref{sec:v10} (v1.0) fits fact tables larger than memory from statistics computed inside the database and proves that this leaves the relational model unchanged and the attribute models an exact posterior given a uniform sample.
\end{abstract}

\tableofcontents

%=====================================================================
\section{Setting and notation}\label{sec:setting}
A database has entity tables $e$ (rows $u\in U_e$, $n_e=|U_e|$) and fact tables $f$. Each entity row has a latent segment $z_u\in\{1..K_e\}$. Links are of five kinds (layer~1, Algorithm~3): \emph{edges} (many-to-one between entities, e.g.\ product$\to$category), \emph{facts} with an anchor $e_1$ and a partner $e_2$, \emph{extra} fact dimensions, \emph{single-dimension} facts, and \emph{trees} (self-references). Column values are mapped to a model space by a codec: discrete columns to codes $0..V-1$, continuous columns to reals (dates to days). $\Phi,\phi$ are the standard normal cdf and density.

The implementation follows the layer~1 pseudocode, with the changes listed in Section~\ref{sec:changes}. Every proposition below names the check that tests it.

%=====================================================================
\section{One table with mixed types}\label{sec:single}

\paragraph{Model.} Within a segment (or the whole table), $\mathbf z\sim\mathcal N_p(0,C)$ and $x_j=F_j^{-1}(\Phi(z_j))$ with $F_j^{-1}(u)=\inf\{x:F_j(x)\ge u\}$. A missing value is the lowest ``level'' of the latent scale: with null rate $q_j$,
\begin{equation}\label{eq:nullmap}
x_j=\begin{cases}\text{null}, & \Phi(z_j)<q_j,\\ G_j^{-1}\!\big((\Phi(z_j)-q_j)/(1-q_j)\big), & \text{otherwise,}\end{cases}
\end{equation}
where $G_j$ is the distribution of the non-null values.

\begin{proposition}[Marginal preservation]\label{prop:pit}
If $C$ has unit diagonal, then for every $j$ the synthetic $x_j$ has exactly the distribution $q_j\delta_{\text{null}}+(1-q_j)G_j$, whatever $C$ is. In particular the null rate is $q_j$ and discrete columns have their level probabilities exactly.
\end{proposition}
\begin{proof}
$z_j\sim\mathcal N(0,1)$, so $U=\Phi(z_j)$ is uniform on $(0,1)$ (probability integral transform). $P(\text{null})=P(U<q_j)=q_j$. Given $U\ge q_j$, $V=(U-q_j)/(1-q_j)$ is uniform, and $P(G_j^{-1}(V)\le x)=P(V\le G_j(x))=G_j(x)$ by the Galois inequality $G^{-1}(v)\le x\iff v\le G(x)$, which holds for any cdf, including step functions.
\end{proof}
\noindent\emph{Check C01}: KS or total-variation distance between real and synthetic, per column, is within the two-sample critical value at the 1\% level (Section~\ref{sec:thresholds}).

\paragraph{Bayesian marginals.} The synthesiser draws $F_j$ from its posterior, so that synthetic data reflect estimation uncertainty (posterior predictive synthesis):
\begin{itemize}
\item discrete: $\mathbf p\sim\Dir(\mathbf c+a\,\mathbf p_0)$ with exact level counts $\mathbf c$ and $\mathbf p_0$ the parent marginal (the whole table for a segment);
\item continuous: Dirichlet-process posterior $F\sim\DP\big(a+n,\ (aF_0+\sum_i\delta_{x_i})/(a+n)\big)$, drawn as $F=V\cdot\sum_iw_i\delta_{x_i}+(1-V)F_0$ with $w\sim\Dir(1,\dots,1)$ (Bayesian bootstrap, Rubin 1981) and $V\sim\mathrm{Beta}(n,a)$;
\item null rate: $q\sim\mathrm{Beta}(n_\text{null}+a q_0,\ n-n_\text{null}+a(1-q_0))$.
\end{itemize}

\begin{proposition}[Posterior mean of a segment marginal]\label{prop:dp}
Under the draws above, $\E[F(t)\mid\text{data}]=\big(aF_0(t)+nF_n(t)\big)/(a+n)$ for every $t$, where $F_n$ is the empirical cdf of the segment. Small segments are therefore shrunk towards the parent marginal, large ones follow their own data.
\end{proposition}
\begin{proof}
$\E[V]=n/(n+a)$, $\E[\sum_iw_i\mathbf 1\{x_i\le t\}]=F_n(t)$ because $\E w_i=1/n$, and $V$ is independent of $w$. The draw is the stick-breaking-free representation of the DP posterior (Ferguson 1973): the posterior base measure is the stated mixture, and the $\DP(a+n,\cdot)$ restricted to the atoms gives Dirichlet weights proportional to $(1,\dots,1,a)$, which is what $(Vw,1-V)$ is.
\end{proof}
{\sloppy\noindent\emph{Check C17} (400 draws, closed form within 0.03). The prior strength $a$ is \texttt{marginals.dirichlet\_prior} (default 0.25: a segment of 6 rows keeps 96\% of its own distribution).\par}

\paragraph{Correlation.} Stage 2 of layer 1's Algorithm~2: latent bins $(\ell_{ij},u_{ij}]$ from the full-data marginals, data augmentation $z_{ij}\mid\cdot\sim\TN(\mu_{ij},\sigma_j^2;\ell_{ij},u_{ij})$, and Metropolis on the C-vine partial correlations with the LKJ$(\eta)$ prior.

\begin{proposition}[The sampler targets the copula posterior]\label{prop:copula}
The chain alternating (i) the column-wise truncated-normal updates of $Z$ and (ii) Metropolis--Hastings steps on $y=\operatorname{atanh}(w)$ with target $\pi(y)=\sum_k(b_k-1)\log(1-w_k^2)-\tfrac n2\log|C|-\tfrac12\operatorname{tr}(C^{-1}Z^\top Z)+\sum_k\log(1-w_k^2)$ (up to constants, $b_k$ the LKJ Beta shapes) has the joint posterior $p(C,Z\mid\text{bins})$ as its invariant law.
\end{proposition}
\begin{proof}
(i) is an exact Gibbs step: given $C$ and the other columns, $z_{ij}$ is normal with mean $-\sum_{k\ne j}\Omega_{jk}z_{ik}/\Omega_{jj}$ and variance $1/\Omega_{jj}$ ($\Omega=C^{-1}$), restricted to its bin. (ii) The C-vine map $w\mapsto C$ is a bijection from $(-1,1)^{p(p-1)/2}$ onto correlation matrices, and under $\LKJ(\eta)$ the $w_k$ are independent with $(w_k+1)/2\sim\mathrm{Beta}(b_k,b_k)$ (Lewandowski, Kurowicka and Joe 2009); $\log(1-w_k^2)$ in $\pi$ combines the Beta density and the Jacobian $dw/dy=1-w^2$. Random-walk MH with a symmetric proposal leaves $\pi$ invariant; the adaptive proposal covariance is frozen after the burn-in, so the post-burn-in chain is an ordinary MH chain (Haario, Saksman and Tamminen 2001). A Gibbs/MH cycle of kernels each invariant for the joint law is invariant for it.
\end{proof}
\noindent\emph{Check C02}: on 4,000 rows from a known 7-dimensional copula with a binary, a 4-level nominal, an ordinal, a date and a 12\%-missing column, 20 of 21 true correlations fall in the 99\% posterior intervals; posterior means are within 0.08; split $\hat R\le1.1$.

\begin{remark}[Dependent missingness]\label{rem:null}
Because the null is the lowest latent bin of its column, $P(\text{null}_j\mid x_k)$ inherits the dependence of $z_j$ on $z_k$: missingness that is monotone in another column is reproduced (Check C03: $P(\text{cholesterol missing}\mid\text{age quartile})$ within 0.025). Non-monotone missingness needs a separate missingness indicator column, which the spec can add.
\end{remark}

\begin{remark}[Nominal columns]\label{rem:nominal}
A Gaussian copula needs an order on nominal levels. Two devices keep dependence:
(a) levels are ordered by the mean normal score of the most associated other column (largest correlation ratio), which makes the main dependence monotone;
(b) two low-cardinality categoricals with Cram\'er's $V\ge0.3$ and at most 64 joint levels are modelled as one categorical over their combinations, i.e.\ a Dirichlet--categorical clique as in a Bayesian network (layer~1 item~4). By Proposition~\ref{prop:pit} their joint law is then reproduced exactly.
\end{remark}

%=====================================================================
\section{Star schemas}\label{sec:star}
For a fact $f$ between anchor $e_1$ and partner $e_2$, the count of rows for the pair $(u,v)$ is $c(u,v)\sim\Pois(\eta_{z_uz_v})$ independently.

\begin{proposition}[Identifiability of segments]\label{prop:ident}
Let $m_l$ be the number of partner rows in segment $l$, and suppose the vectors $r_k=(\eta_{kl}m_l)_{l}$ are pairwise distinct. Then, given the partner partition, the anchor partition is identified up to relabelling from the expected row-statistic profiles $\E[C^{\mathrm{row}}_{u\cdot}]=r_{z_u}$, and the conditional of $z_u$ concentrates on the true segment as the row's total count grows.
\end{proposition}
\begin{proof}
$C^{\mathrm{row}}_{ul}=\sum_{v:z_v=l}c(u,v)\sim\Pois(\eta_{z_ul}m_l)$, independent over $l$. Distinct $r_k$ give distinct laws of $C^{\mathrm{row}}_{u\cdot}$, so the map $k\mapsto$ law is injective and the partition is determined up to a permutation of labels. For the posterior: by layer~1 Proposition~1, $\log p(z_u=k)-\log p(z_u=k')=\sum_l[C^{\mathrm{row}}_{ul}\log(\eta_{kl}/\eta_{k'l})-(\eta_{kl}-\eta_{k'l})m_l]+O(1)$. Its expectation under $z_u=k$ equals $\sum_l m_l\,\mathrm{KL}(\Pois(\eta_{kl})\|\Pois(\eta_{k'l}))>0$ times the exposure, and it grows linearly in $m_l$ while its standard deviation grows as $\sqrt{m_l}$.
\end{proof}
\noindent\emph{Check C05}: with caps of 8, the sampler and merge recover 4 customer segments (ARI 0.98) and 3 product segments (ARI 1.0).

\begin{proposition}[Block-rate posterior]\label{prop:eta}
Given the partitions, $\eta_{kl}\mid\cdot\sim\Ga(a+C^{\mathrm{blk}}_{kl},\ b+n_km_l)$, with mean $\to C^{\mathrm{blk}}_{kl}/(n_km_l)$ and coefficient of variation $(a+C^{\mathrm{blk}}_{kl})^{-1/2}$.
\end{proposition}
\begin{proof}
Gamma--Poisson conjugacy over the $n_km_l$ independent pair counts of the block, whose sum is $C^{\mathrm{blk}}_{kl}$; zero pairs enter through the exposure only.
\end{proof}
\noindent\emph{Check C06}: all 12 true block rates lie in their 99\% posterior intervals.

\begin{proposition}[Splitting and degree correction]\label{prop:split}
(a) Drawing $N\sim\Pois(\eta\,W_1W_2)$ rows for a block and assigning each to $(u,v)$ with probabilities $\omega_u/W_1\cdot\omega'_v/W_2$ ($W_1=\sum\omega_u$, $W_2=\sum\omega'_v$) gives independent $c(u,v)\sim\Pois(\eta\,\omega_u\omega'_v)$.
(b) If $\omega_u\sim\Ga(\alpha,\alpha)$, the number of rows of an anchor row with expected count $\mu$ is negative binomial with $\E=\mu$ and $\Var=\mu+\mu^2/\alpha$.
(c) Rescaling the $\omega$ of a segment to sum to the segment size leaves every block's expected total unchanged, so the scale rule (Proposition~\ref{prop:scale}) stays exact.
\end{proposition}
\begin{proof}
(a) Poisson colouring: a Poisson total thinned multinomially gives independent Poisson cells with means $\eta W_1W_2\cdot\omega_u\omega'_v/(W_1W_2)$. (b) $\E[N\mid\omega]=\mu\omega$, so $\E N=\mu$ and $\Var N=\E\Var(N\mid\omega)+\Var\E(N\mid\omega)=\mu+\mu^2/\alpha$. (c) The expected total is $\eta W_1W_2$, and after rescaling $W_1=n_k$, $W_2=m_l$ as without frailty.
\end{proof}
The frailty (a degree-corrected block model, Karrer and Newman 2011) is switched on per segment when the dispersion index of per-row totals exceeds 1.5 and the segment has at least 30 rows. \emph{Checks C07} (rows-per-customer distribution within the KS critical value; repeated-pair share equal) \emph{and C16} (variance formula within 2\%).

\begin{proposition}[Scale rule]\label{prop:scale}
With $\eta^*_{kl}=\eta_{kl}m_l/m^*_l$, the expected number of rows per anchor row is unchanged when any table is resized; resizing the anchor by $s$ multiplies the fact's expected size by $s$.
\end{proposition}
\begin{proof}
Layer~1, Proposition~4: $\sum_l\eta^*_{kl}m^*_l=\sum_l\eta_{kl}m_l$. With $n^*_k=sn_k$, $\E N^*=\sum_{kl}\eta^*_{kl}n^*_km^*_l=s\sum_{kl}\eta_{kl}n_km_l$.
\end{proof}
\noindent\emph{Check C08}: products $\times2$ and $\times0.5$ keep rows per customer within 4\%; customers $\times2$ doubles the fact (2.03).

%=====================================================================
\section{Snowflake chains and joins}\label{sec:snow}
An edge child$\to$parent is the same Poisson block model with $c(u,v)\in\{0,1\}$. At synthesis each child row picks its parent segment from $\varphi_{kt}\propto\eta_{kt}m_t$, then a parent inside the segment (with frailty weights).

\begin{proposition}[Chained segments compose]\label{prop:chain}
For a chain fact$\to$product$\to$category$\to$department, the synthetic law of the department segment of a fact row with product segment $l$ is $\sum_t\varphi^{(p\to c)}_{lt}\varphi^{(c\to d)}_{ts}$. If the source satisfies the Markov property ``a row's ancestor two levels up depends on the row only through its parent's segment'', then every joint law of a fact attribute and an ancestor attribute is reproduced in expectation.
\end{proposition}
\begin{proof}
Each synthetic step draws the next segment from $\varphi$ given the current segment only, so the chain of segments is a Markov chain whose transition matrices are the fitted $\varphi$; the two-step law is their product (Chapman--Kolmogorov). Attributes are drawn given their own row's segment. Under the stated Markov property the source has the same factorisation, so the joint law of (fact attribute, ancestor attribute) is $\sum$ over segment paths of the product of the same factors.
\end{proof}

\paragraph{Fact measures that depend on a parent's value.} Segments carry a parent's attributes only coarsely; a sale amount that follows the product's own price needs more. For every fact, continuous attributes of its parents (at most two per parent) are joined onto the fact rows and enter the fact's copula as conditioning columns $a$. At synthesis the parent is chosen first; its synthetic attribute values give $z_a$, and the fact's own columns are drawn from $z_r\mid z_a\sim\mathcal N(C_{ra}C_{aa}^{-1}z_a,\ C_{rr}-C_{ra}C_{aa}^{-1}C_{ar})$, which is the exact Gaussian conditional, so Proposition~\ref{prop:pit} holds for the pair (fact column, parent column) jointly.

\noindent\emph{Check C09}: segments of department, category and customer are recovered (ARI 1.0, 1.0, 0.997). Three-hop statistics (share of sales by department division, mean amount per division, Spearman correlation of amount with the product's price) fall inside the posterior predictive spread of 20 synthetic databases ($|z|\le3$; the amount--price correlation is 0.909 real vs 0.910 synthetic).

%=====================================================================
\section{Self-referencing hierarchies}\label{sec:tree}
A self-reference (employee$\to$manager) is a multi-type Galton--Watson process: a node of type $k$ has $N\sim\Pois(\lambda_k)$ children (with optional Gamma frailty), each of type $l$ with probability $\psi_{kl}$; roots have types $\Cat(\pi^{\mathrm{root}})$. Let $M_{kl}=\lambda_k\psi_{kl}$ be the mean matrix.

\begin{proposition}[Generation sizes, finiteness and expected size]\label{prop:gw}
Let $r$ be the row vector of expected root counts per type and $Z_n$ the vector of type counts in generation $n$. Then $\E Z_n=rM^n$. If the spectral radius $\rho(M)<1$, every tree is finite almost surely and $\E|T|=r(I-M)^{-1}\mathbf 1$.
\end{proposition}
\begin{proof}
$\E[Z_{n+1}\mid Z_n]=Z_nM$ by linearity over parents, so $\E Z_n=rM^n$ by induction. If $\rho(M)<1$ then $\sum_nM^n=(I-M)^{-1}$ converges, so $\E\sum_n Z_n\mathbf 1=r(I-M)^{-1}\mathbf 1<\infty$, and a non-negative variable with finite mean is finite a.s.
\end{proof}

\begin{proposition}[Fit and validity]\label{prop:treefit}
(a) Given the types, the posteriors are $\psi_k\sim\Dir(\tfrac12+B_{k\cdot})$ with $B_{kl}$ the number of type-$k$ parents' type-$l$ children, and $\lambda_k\sim\Ga(a+\sum_{u:z_u=k}N_u,\ b+n_k)$. (b) Nodes at even depth are conditionally independent given nodes at odd depth, so the two-colour blocked Gibbs sweep is valid. (c) Synthesis generates each node once, as the child of an already generated node, so the synthetic hierarchy is a forest: every non-root has exactly one existing parent and there are no cycles. Growth stops at a maximum depth (by default the deepest level of the source, so no node is deeper than any source node; \texttt{trees.max\_depth} overrides it).
\end{proposition}
\begin{proof}
(a) Multinomial--Dirichlet and Poisson--Gamma conjugacy. (b) The type of $u$ appears only in factors with its parent and its children, whose depths have the other parity. (c) By construction of the breadth-first generation.
\end{proof}
\noindent\emph{Checks C10, C11}: types recovered (ARI 1.0); $\rho(\hat M)$ within 0.05 of the source tree's realised value and below 1; the formula for $\E|T|$ agrees with 300 simulations within three standard errors; synthetic hierarchies have no cycles and no dangling managers; the depth histogram passes a posterior predictive check against 40 re-simulations ($p\ge0.05$); children-per-node KS 0.016; grade given depth within 0.10.

\begin{remark}
Near-critical hierarchies ($\rho(M)$ close to 1, as in the test tree: managers managing managers with mean 0.9) have depth profiles that vary a lot between realisations. A fixed tolerance on the depth histogram is then wrong; the posterior predictive check (Gelman, Meng and Stern 1996) is the right test, and is what the evaluator uses.
\end{remark}

%=====================================================================
\section{Temporal facts}\label{sec:time}
Rows carry a time bucket $b$ (day/week/month/quarter/year). Following layer~1's design, $c(u,v,b)\sim\Pois(\eta_{kl}\lambda_{kb}\Delta_b)$ with $\sum_b\lambda_{kb}\Delta_b=\sum_b\Delta_b$.

\begin{proposition}[Seasonality by multinomial thinning]\label{prop:time}
Given a block total $N_{kl}$, bucket counts are $\Mult(N_{kl};\ \theta_{k\cdot})$ with $\theta_{kb}\propto\lambda_{kb}\Delta_b$; the posterior of the profile is $\theta_k\sim\Dir(\theta_0+n_{k\cdot})$ with $n_{kb}$ the source rows of anchor segment $k$ in bucket $b$. Synthesis that draws a bucket per row from $\theta_k$, then a uniform instant inside the bucket, therefore has the time law of the model.
\end{proposition}
\begin{proof}
Independent Poisson cells conditioned on their sum are multinomial with probabilities proportional to their means; the factor $\eta_{kl}$ cancels. Dirichlet--multinomial conjugacy gives the posterior.
\end{proof}
\noindent\emph{Check C12}: two customer segments with opposite seasonality are recovered (ARI 0.995); fitted monthly profiles are within TVD 0.05 of the true ones; the synthetic monthly histogram is within 0.03 overall and 0.04 within each plan.

%=====================================================================
\section{Text}\label{sec:text}
Free-text columns use a class-conditional hierarchical Dirichlet back-off $n$-gram model (MacKay and Peto 1995). With $h_n$ the last $n$ tokens, $c$ the value of the conditioning column (or the row's segment) and $p_G$ the add-one global unigram,
\begin{equation}\label{eq:lm}
p_c(w\mid h_n)=\frac{n_c(h_n,w)+\beta\,p_c(w\mid h_{n-1})}{n_c(h_n)+\beta},\qquad p_c(w\mid h_0)=\frac{n_c(w)+\beta\,p_G(w)}{n_c+\beta}.
\end{equation}

\begin{proposition}[Proper predictive]\label{prop:lm}
For every class and context, $p_c(\cdot\mid h_n)$ is a probability distribution on the vocabulary, and it is the posterior predictive of a Dirichlet prior with mean $p_c(\cdot\mid h_{n-1})$ and concentration $\beta$ on the next-token distribution after $h_n$.
\end{proposition}
\begin{proof}
Induction on $n$: $p_G$ sums to one; if $p_c(\cdot\mid h_{n-1})$ does, the numerator of \eqref{eq:lm} sums to $n_c(h_n)+\beta$. The second statement is Dirichlet--multinomial conjugacy. Contexts never seen use the shorter context, which is the same formula with zero counts.
\end{proof}

\begin{proposition}[Memorisation guards]\label{prop:guard}
Let $k$ be \texttt{min\_token\_rows} and $L$ \texttt{max\_copy\_ngram}. (a) Every token in synthetic text occurs in at least $k$ distinct source rows. (b) No synthetic text contains a run of $L$ tokens that occurs in fewer than $k$ source rows, nor equals a source text shared by fewer than $k$ rows. (c) Accepted texts are distributed as the model law conditioned on the guard event.
\end{proposition}
\begin{proof}
(a) Rarer tokens are removed from the vocabulary before counting, so they have probability zero. (b) Candidates that violate the rule are rejected. (c) Rejection sampling: a candidate is accepted with probability $\mathbf 1\{\text{guard}\}$, so accepted draws have density proportional to $p\cdot\mathbf 1\{\text{guard}\}$. After \texttt{max\_retries} failures the cell is left empty rather than emitting a violating text.
\end{proof}
Emails and long digit runs are replaced by typed placeholders before training and refilled with fresh values. These guards bound verbatim memorisation; they are not differential privacy. For DP text the layer~1 route (Private Evolution over a foundation model, Xie et al.\ 2024) stays the option, and the same guards apply to its output. \emph{Check C13}: the predictive sums to one within $10^{-9}$ over 125 random contexts; none of the uniquely named tickets' names appears; no rare token or rare 8-gram is copied; a classifier trained on synthetic tickets predicts the real ticket category as well as one trained on real tickets (ratio 1.0).

%=====================================================================
\section{Business rules and the whole schema}\label{sec:rules}
Rules come from the spec: \texttt{derive} rules (\texttt{total = quantity * unit\_price - discount}) are computed after sampling; \texttt{check} rules are enforced by redrawing the attributes of violating rows from the same model (same segment and parents).

\begin{proposition}[Rules by rejection]\label{prop:rules}
If a row's attributes are drawn from $p$ and redrawn until the check $A$ holds, the result has law $p(\cdot\mid A)$, and the expected number of draws is $1/p(A)$. Derived columns satisfy their equation exactly.
\end{proposition}
\begin{proof}
Standard rejection sampling with acceptance indicator $\mathbf 1_A$; the number of draws is geometric with success probability $p(A)$.
\end{proof}
\noindent\emph{Checks C14, C15}: on the 9-table enterprise schema (regions, customers, stores, a three-level product hierarchy, a 4-dimension sales fact with dates, tickets with text, and an employee hierarchy) the full scorecard passes, all 9 foreign keys resolve, derived totals are exact and no check rule is violated; on a test distribution the redrawn sample matches the exact truncated law (KS 0.003). \emph{Check C18}: the same seed gives a bit-identical database.

\paragraph{Reference tables.} Small lookup tables (regions, departments) are usually not personal data and their individual rows matter (one region is far bigger than another). A table with \texttt{mode: copy} is kept as is; each of its rows is its own fixed segment, so Proposition~\ref{prop:eta} then models every child's choice of that exact parent, which reproduces each parent's popularity. Introspection proposes this for tables of at most 100 rows whose parents are also copied.

%=====================================================================
\section{Grain, structural zeros and computed columns (v0.3)}\label{sec:v03}
The first real-database run (a business model of a supervisory reporting framework: reporters, submissions, templates, cells and reported values) failed in ways the fake scenarios had not exercised: duplicate reports per reporter and date, duplicate values on fixed cells, templates filed partially or twice, values in cells of the wrong framework version, and a stored count that no longer matched the rows it counts. None of these is specific to that model. Each comes from one of four general situations: a table has a \emph{grain} (a unique key made of foreign keys); a combination of parents is \emph{structurally absent} (never allowed) rather than rare; a column is \emph{computed} from other tables; or a key is \emph{partly null}. This section gives the model for each and proves what it guarantees. Scenario S8 (a generic reporting model with known truth: reporters, periods, two framework versions, forms with sheets, items with open rows, omitted and partial filings, an aggregate and a derived status) tests them together.

\subsection{Rules and computed columns}
Rules are parsed into a restricted syntax tree (no attribute access, no calls outside a fixed list), evaluated column-wise with null semantics. The string functions (\texttt{regex\_replace}, \texttt{regex\_match}, \texttt{substr}, \texttt{concat}, \texttt{upper}, \texttt{lower}, \texttt{strip}, \texttt{length}) require constant patterns.

\begin{proposition}[Three-valued rules agree with the database]\label{prop:k3}
Logical connectives are evaluated in Kleene's strong three-valued logic, comparisons with a null are unknown, \texttt{where} takes its else branch on unknown, functions of a null are null, and a \texttt{check} fails only when it evaluates to false. These are the semantics of SQL \texttt{AND}, \texttt{OR}, \texttt{NOT}, \texttt{CASE} and \texttt{CHECK}; hence a row accepted by a database \texttt{CHECK} constraint is accepted by the same rule here, and conversely.
\end{proposition}
\begin{proof}
SQL defines \texttt{AND}/\texttt{OR}/\texttt{NOT} by the strong Kleene truth tables (ISO/IEC 9075-2, 8.20) and a \texttt{CHECK} constraint as violated only when its condition is false. The evaluator implements exactly these tables elementwise (unknown is a third value, not false), so both evaluations agree row by row.
\end{proof}
\noindent\emph{Check C20}: the three connectives on all nine pairs of truth values, nulls in \texttt{where}, \texttt{between}, \texttt{concat} and \texttt{regex\_replace}, and refusal of non-constant patterns and of calls outside the list.

A column is \emph{computed} when it is a \texttt{lookup} (\texttt{fk.col}, possibly several hops), a \texttt{derive} rule, or an \texttt{aggregate} rule (count, sum, min, max, mean, exists or count of distinct values of the child rows that reference the row, through a plain or a composite key). Computed columns form a directed dependency graph $G$: $c\to c'$ when the rule of $c'$ reads $c$.
\begin{proposition}[Computed columns are exact]\label{prop:computed}
(a) If $G$ is acyclic, the fixed-point pass of \texttt{prepare\_source} fills every computed column missing from the source, each with the value of its rule on the final inputs, and reports, for each stored computed column, the share of rows where the stored value differs from its rule. (b) On the synthetic database every computed column is recomputed after the tables it reads are generated (aggregates and the derives that read them last, then lookups and derives of the dependent tables again), so every derive, lookup and aggregate identity holds exactly. If $G$ has a cycle, the columns on it are reported and no value is invented.
\end{proposition}
\begin{proof}
(a) Each pass visits tables in foreign-key order and computes a column once all columns it reads exist; a column at depth $d$ of $G$ is therefore filled by pass $d+1$, and a pass without progress certifies that the remaining columns lie on or above a cycle. Values are computed from the final inputs because inputs are filled before readers. (b) Generation never samples a computed column; the late recomputation runs in topological order of $G$ restricted to synthetic tables, so each value equals its rule applied to the final values it reads.
\end{proof}
\noindent\emph{Checks C26, C32}: on S8 the missing \texttt{form\_code} is filled from its rule, 10\% planted mismatches in a stored count are reported as 0.0976 (the planted share), and on the synthetic data the count, the derived status and the form code match their rules on every row.

\subsection{Structural zeros}
A real schema has combinations that never occur because they are not allowed (a cell of another framework version, a header item missing from a filed form, a null in a mandatory column), mixed with combinations that are merely rare. A Beta or Dirichlet prior gives every combination positive probability, so with enough synthetic rows the forbidden ones appear. A spike-and-slab prior separates the two.

\begin{proposition}[Spike-and-slab Bernoulli]\label{prop:zoib}
Let $p\sim\tfrac{\pi_s}{2}\delta_0+\tfrac{\pi_s}{2}\delta_1+(1-\pi_s)\,\mathrm{Beta}(a,b)$ and $S\mid p\sim\mathrm{Bin}(T,p)$. The posterior is
\[
p\mid S\;\propto\;\tfrac{\pi_s}{2}[S=0]\,\delta_0+\tfrac{\pi_s}{2}[S=T]\,\delta_1+(1-\pi_s)\tfrac{B(a+S,\,b+T-S)}{B(a,b)}\,\mathrm{Beta}(a+S,\,b+T-S),
\]
with mean $w_1+w_2\,\frac{a+S}{a+b+T}$ for the normalised weights $w_0,w_1,w_2$. The same holds for a null rate (\texttt{\_null\_rate}) and for each occupancy probability (\texttt{zoib\_draw}).
\end{proposition}
\begin{proof}
Bayes' rule with a mixture prior: each component's posterior weight is its prior weight times its marginal likelihood; the binomial coefficient is common to all three and cancels; the point masses have likelihood $[S=0]$ and $[S=T]$, the slab $B(a+S,b+T-S)/B(a,b)$.
\end{proof}

When the rows $k$ of a group $g$ (the segments of one stratum, e.g.\ the submissions of one framework version) share the structure, the evidence must be pooled: each row alone may have seen a combination only a few times.
\begin{proposition}[Grouped structural support]\label{prop:grouped}
Let $G_{gl}\in\{\text{all zero},\text{all one},\text{free}\}$ with prior $(\tfrac{\pi_s}{2},\tfrac{\pi_s}{2},1-\pi_s)$, $p_{kl}=0$ (resp.\ 1) for all $k\in g$ if $G_{gl}$ is all zero (all one), and $p_{kl}$ independent with the prior of Proposition~\ref{prop:zoib} if free. Then $P(G_{gl}\mid S)\propto\big(\tfrac{\pi_s}{2}\prod_{k}[S_{kl}=0],\ \tfrac{\pi_s}{2}\prod_k[S_{kl}=T_{kl}],\ (1-\pi_s)\prod_k m_{kl}\big)$ with $m_{kl}$ the per-row marginal of Proposition~\ref{prop:zoib}, and given ``free'' the $p_{kl}$ follow their per-row posteriors. \texttt{grouped\_zoib\_draw} samples this posterior exactly. In particular, a combination never observed in a group is absent from every row of the group with posterior probability tending to 1 as the number of rows with exposure ($T_{kl}\ge1$) grows, even when every $T_{kl}$ stays bounded.
\end{proposition}
\begin{proof}
The likelihood factorises over $k$ given $G_{gl}$; under the two pooled states it is the product of indicators, under ``free'' the product of per-row marginals because the $p_{kl}$ are independent. Given ``free'', the posterior of $(p_{kl})_k$ factorises, which is what the per-row draw returns. For the limit: if all $S_{kl}=0$ and $T_{kl}\ge1$, then $m_{kl}=\tfrac{\pi_s}{2}+(1-\pi_s)B(a,b+T_{kl})/B(a,b)\le\tfrac{\pi_s}{2}+(1-\pi_s)\tfrac{b}{a+b}<1$, because $B(a,b+T)/B(a,b)$ decreases in $T$ and equals $b/(a+b)$ at $T=1$. So $\prod_km_{kl}\to0$ geometrically while the ``all zero'' weight $\pi_s/2$ is constant. A group without exposure gets its prior back ($m_{kl}=1$, both indicators hold).
\end{proof}
\noindent This is used for the occupancy probabilities $p_{kl}$ and for the anchor activity $\pi$ (which parents have any rows) with the strata as groups.

\begin{proposition}[Sparse Dirichlet, exact inclusion]\label{prop:sparsedir}
Let each level $j$ be structurally absent with prior probability $\pi_s$ independently, $p_J\sim\Dir(\alpha_J)$ on the included set $J$, and counts $c\sim\Mult(n,p)$. Observed levels are included; for the set $Z$ of unobserved levels, with $\rho=(1-\pi_s)/\pi_s$, $A_o=\sum_{c_j>0}\alpha_j$,
\[
P(J\cap Z=I\mid c)\;\propto\;\rho^{|I|}\,\frac{\Gamma(A_o+\alpha_I)}{\Gamma(A_o+\alpha_I+n)},\qquad \alpha_I=\textstyle\sum_{j\in I}\alpha_j,
\]
and with an auxiliary $u\in(0,1)$ of density $\propto u^{A_o-1}(1-u)^{n-1}\prod_{j\in Z}(1+\rho u^{\alpha_j})$, the indicators $[j\in J]$ are independent given $u$ with probabilities $\rho u^{\alpha_j}/(1+\rho u^{\alpha_j})$. \texttt{sparse\_dirichlet\_draw} samples $u$ by inverse cdf on a grid of $\operatorname{logit}u$, then $J$, then $p_J\sim\Dir(c_J+\alpha_J)$.
\end{proposition}
\begin{proof}
The Dirichlet-multinomial likelihood is $\frac{n!}{\prod c_j!}\frac{\Gamma(A_J)}{\Gamma(A_J+n)}\prod_{j\in J}\frac{\Gamma(\alpha_j+c_j)}{\Gamma(\alpha_j)}$; unobserved levels contribute factor 1 to the product, giving the display. Since $\Gamma(x)/\Gamma(x+n)=B(x,n)/\Gamma(n)=\Gamma(n)^{-1}\int_0^1u^{x-1}(1-u)^{n-1}du$, the joint of $(u,I)$ is $\propto u^{A_o-1}(1-u)^{n-1}\prod_{j\in I}\rho u^{\alpha_j}$, whose $I$-marginal is the display and whose conditional factorises over $j$. Given $J$, conjugacy gives the Dirichlet.
\end{proof}
\noindent The earlier one-level-at-a-time Bayes factor is exact only when one level is unobserved; with five unobserved levels it under-stated absence by up to 0.07 (rejection reference), which is why it was replaced. \emph{Check C21}: zoib, grouped zoib, sparse Dirichlet (one and five unobserved levels) and null rates agree with exact posteriors obtained by rejection sampling from the generative model (all within Monte Carlo error, e.g.\ 0.857 vs 0.857, 3.242 vs 3.228 absent levels).

\subsection{Partial pooling}
\begin{proposition}[Empirical-Bayes Beta slab]\label{prop:ebbeta}
For row $k$ of a block matrix of presence counts, on its supported entries ($S_{kl}>0$) let $m=\sum S/\sum T$ and $\hat\tau^2=\widehat{\Var}(S_l/T_l)-\overline{m(1-m)/T_l}$. If $p_{kl}\sim\mathrm{Beta}(cm,c(1-m))$ with $S_{kl}\sim\mathrm{Bin}(T_{kl},p_{kl})$, then $\hat c=m(1-m)/\hat\tau^2-1$ is the moment estimator of $c$ (Kleinman 1973): as the number of entries grows it converges to $c$ up to a relative error $O(1/\bar T)$. The slab $\mathrm{Beta}(\hat cm,\hat c(1-m))$, bounded to $\hat c\in[1,10^4]$, shrinks entries whose rates differ only by binomial noise towards $m$ and leaves real differences.
\end{proposition}
\begin{proof}
$\E[S/T]=m$ and $\Var(S/T)=\E\Var(S/T\mid p)+\Var(p)=\E[p(1-p)]/T+\Var(p)$ with $\Var(p)=m(1-m)/(c+1)$ and $\E p(1-p)=m(1-m)-\Var p$. Hence $\Var(S/T)-m(1-m)/T=\Var(p)\,(1-1/T)$, so $\hat\tau^2\to\Var(p)(1+O(1/\bar T))$ by the law of large numbers, and $c=m(1-m)/\Var(p)-1$ gives the limit of $\hat c$ by continuity. The bounds keep the prior proper and stop $\hat\tau^2\le0$ (no excess variance) from producing an infinite $c$.
\end{proof}

\begin{proposition}[Multiplicities: spike-and-slab Poisson]\label{prop:mult}
Excess rows $X_g$ over $E_g$ present pairs of a group $g$ (partner, anchor or block) have prior $\lambda_g\sim\pi_s\delta_0+(1-\pi_s)\Ga(a_m,b_m)$ and $X_g\mid\lambda_g\sim\Pois(\lambda_gE_g)$. The collapsed marginal is $\pi_s[X_g=0]+(1-\pi_s)\frac{b_m^{a_m}\Gamma(a_m+X_g)}{\Gamma(a_m)(b_m+E_g)^{a_m+X_g}}$ per group (up to a factor common to all groupings), the grouping with the highest product of these over its groups is chosen, and $b_m$ is set so that the slab mean $a_m/b_m$ equals $\sum_{X_g>0}X_g/\sum_{X_g>0}E_g$: the groups with $X_g>0$ are certainly in the slab, so pooling the spikes into the mean would dilute it.
\end{proposition}
\begin{proof}
Each present pair $i$ of group $g$ has $x_i\sim\Pois(\lambda_g)$, so the likelihood is $\lambda_g^{X_g}e^{-\lambda_gE_g}/\prod_ix_i!$; Poisson-Gamma conjugacy gives the slab marginal shown and the mixture follows as in Proposition~\ref{prop:zoib}. The factor $\prod_i1/x_i!$ is the same for every grouping of the same pairs, so comparing groupings is a Bayes-factor comparison. Groups with $X_g>0$ have zero spike likelihood, so they belong to the slab; selecting on $X_g>0$ raises the estimate by at most the factor $1/(1-(b_m/(b_m+E_g))^{a_m})$, which is close to 1 for the exposures used.
\end{proof}
\noindent\emph{Checks C22, C28}: the concentration is recovered (median 20.5 for true 20), shrinkage lowers the error of $p$ by 27\%, homogeneous rates pool almost fully; S8 open rows repeat 4.10 times on average (truth 4) and the partner grouping is selected.

\subsection{Grain}
\begin{proposition}[Repeat test]\label{prop:repeat}
Under the degree-corrected Poisson model with $\lambda_{ij}=r_is_j/N$ (plug-in row and column totals), the number of pairs with two or more rows is a sum of independent Bernoulli variables with probabilities $q_{ij}=1-e^{-\lambda_{ij}}(1+\lambda_{ij})$, so $z=(\text{obs}-\sum q)/\sqrt{\sum q(1-q)}$ is asymptotically standard normal (Lyapunov). Under an occupancy model with a unique key, obs $=0$. The engine selects occupancy when $z<-4$ and $\sum q\ge10$, or when a unique key on the pair is declared.
\end{proposition}
\begin{proof}
Counts of distinct pairs are independent Poisson variables; $P(\text{count}\ge2)=q_{ij}$; the variances are bounded and their sum diverges with the number of pairs, which is Lyapunov's condition.
\end{proof}
\begin{remark}
A block structure on top of the degree correction concentrates the rates, and $q(\lambda)\approx\lambda^2/2$ is convex, so it tends to raise the expected repeats: the rank-one reference is the conservative choice. On S8 the filing and value tables give $z=-29$ and $z=-61$; the S2 sales fact keeps the Poisson model (check C24).
\end{remark}

\begin{proposition}[Grain entities]\label{prop:grainent}
Let an entity table $R$ have a unique key $(a,b)$ of two foreign keys to $A$ and $B$. The model: pair $(a,b)$ is present with probability $p_{z_az_b}$ (occupancy, priors of Propositions~\ref{prop:grouped} and \ref{prop:ebbeta}), and the row's own segment is drawn from $P(z\mid z_a,z_b)\propto\pi_z\,\phi^A_{z z_a}\phi^B_{z z_b}$. Then (a) $R^*$ has no duplicate key; (b) the expected number of rows is $\sum_{a,b}p_{z_az_b}$, which scales with $s_A s_B$ when $A,B$ are scaled by $s_A,s_B$ (a copied $B$ has $s_B=1$); (c) the segment draw is the exact conditional of a latent class model in which $z_a$ and $z_b$ are independent given $z$.
\end{proposition}
\begin{proof}
(a) Each pair is generated at most once. (b) Linearity of expectation over the $n_An_B$ pairs, whose segment frequencies are preserved by scaling (Proposition~\ref{prop:scale}). (c) Bayes' rule with prior $\pi$ and the two conditionally independent likelihood factors.
\end{proof}
\noindent\emph{Checks C26, C29}: no duplicate report at scale 1 or 2; reports per reporter 7.675 (real 7.650) at scale 2 with the periods copied.

\begin{proposition}[Unique keys: de-duplication is exact conditioning]\label{prop:unique}
Let a pair be present with probability $p$ and, when present, carry $M=1+N$ rows with exchangeable attributes. Keeping the first row of each pair gives (a) presence exactly $\mathrm{Bern}(p)$, (b) kept attributes with exactly the law of one row. For a filtered key (\texttt{unique: \{columns, where\}}) the same holds on the rows that satisfy the filter, which the occupancy model generates with multiplicity one; rows outside the filter (open rows) keep their multiplicity.
\end{proposition}
\begin{proof}
(a) $\{M\ge1\}=\{\text{present}\}$. (b) By exchangeability the first row has the marginal law of any row and does not depend on $N$.
\end{proof}
\noindent\emph{Checks C23, C28}: presence 0.3522 for $p=0.35$ (within 3 s.d.) and KS 0.005 for the kept rows; on S8 no closed item is ever duplicated while open rows keep their multiplicity.

\begin{proposition}[Strata]\label{prop:strata}
Let every segment of an entity belong to exactly one level of a stratum key (e.g.\ the framework version of a submission). (a) Setting the key from the segment makes every synthetic (segment, level) pair one that occurs in the real data. (b) With the strata as groups in Proposition~\ref{prop:grouped}, a partner never seen with a level is absent from that level with posterior probability tending to 1 as the level's exposed segments grow; for grain entities the segment draw is also restricted to segments compatible with both parents' levels, which is exact.
\end{proposition}
\begin{proof}
(a) The key is a function of the segment. (b) Proposition~\ref{prop:grouped}; the restriction of the segment draw puts zero probability outside the compatible set.
\end{proof}
\noindent\emph{Check C26}: the item$\to$version link of S8 is never declared, yet every synthetic value uses an item of its report's version (share 1.0).

\begin{proposition}[Nullable extra keys]\label{prop:nullfk}
Let a fact row have anchor and partner keys and further keys that may be null. If rows with a null extra key are dropped before fitting, the fitted occupancy targets $P(\text{present},\ \text{key}\ne\text{null}\mid\text{block})$ instead of $P(\text{present}\mid\text{block})$, a bias by the factor $P(\text{key}\ne\text{null}\mid\text{block})$ that varies across blocks. Treating null as one more level of the extra key's distribution, $P(\text{segment or null}\mid\text{block})$ with the prior of Proposition~\ref{prop:sparsedir}, removes it.
\end{proposition}
\begin{proof}
$P(\text{present},\text{key}\ne\text{null})=P(\text{present})P(\text{key}\ne\text{null}\mid\text{present, block})$. With null as a level the likelihood of each row factorises into the pair term and the extra-key term on all rows, so both are fitted on the full data.
\end{proof}
\noindent S8 found this: 84\% of its values have no code, the earlier fit kept only coded values and generated 5{,}027 rows instead of 31{,}817. With the fix, 31{,}622 (check C26).

\begin{proposition}[Nested refinement]\label{prop:nested}
When a fact is nested in a parent fact (values inside a filing, through a composite key), the parent-fact rows inherit their block from the parent model and are refined into sub-segments inside each block with the hierarchical prior $p_{k'l}\sim\mathrm{Beta}(c\mu_{kl},c(1-\mu_{kl}))$, $c=2$, centred on the block's rate $\mu_{kl}$. The collapsed score of a partition is a sum of per-segment Beta-binomial terms, so a merge of two sub-segments changes only their terms and merges in different blocks are independent.
\end{proposition}
\begin{proof}
Given the block, sub-segments have independent rates under the prior; integrating them out gives a product over sub-segments and partner segments of $B(c\mu+S,c(1-\mu)+T-S)/B(c\mu,c(1-\mu))$.
\end{proof}
\noindent Without the centring, unrelated cells penalised every split and each block collapsed to one segment; with it, partial and complete filings are separated. \emph{Check C27}: declared-not-filed and partial filings match within sampling error.

\subsection{Dependence and posterior predictive draws}
\begin{proposition}[Rank likelihood with parameter expansion]\label{prop:rank}
Let $\mathcal D(Y)=\{Z:z_{ij}<z_{i'j}\text{ whenever }y_{ij}<y_{i'j}\}$. (a) $P(Z\in\mathcal D(Y)\mid C)$ depends on $Y$ only through the within-column ranks, so the posterior of $C$ is invariant to increasing transformations of any column. (b) $\mathcal D$ is invariant to positive rescaling of columns, so with $\Sigma=DCD$ ($D$ diagonal) and $\Sigma\sim\mathrm{IW}(p+1,I)$ (uniform marginal priors on every correlation, Barnard, McCulloch and Meng 2000) the Gibbs sampler $Z\mid\Sigma$ (truncated normals within the rank bounds), $\Sigma\mid Z\sim\mathrm{IW}(p+1+n,I+Z^\top Z)$ has the posterior of $C$ as the law of $D^{-1}\Sigma D^{-1}$ (Hoff 2007). (c) Extending the chain in blocks until split $\hat R\le1.1$ and discarding the first half of the post-burn-in draws keeps every draw a state of the same chain.
\end{proposition}
\begin{proof}
(a) $\mathcal D(Y)$ is defined by order relations only. (b) $P(Z\in\mathcal D\mid\Sigma)=P(D^{-1}Z\in\mathcal D\mid C)$ because $D^{-1}Z\sim\mathcal N(0,C)$ and $D^{-1}\mathcal D=\mathcal D$; the IW update is conjugate for $Z$ given $\Sigma$. (c) Run-length control does not change the transition kernel; stopping by a diagnostic can bias finite runs (Cowles, Roberts and Rosenthal 1999), which the half-discard and the R2 drift check bound in practice.
\end{proof}
\noindent The implementation starts from the normal scores of the bins' mid-ranks, so the whole chain is a function of the bin order: \emph{check C25} finds the draws bit-identical after a strictly increasing transform of every column, error 0.034 against the true correlations, $\hat R=1.00$; on the real submissions table $\hat R$ fell from 1.58 to below 1.1.

\begin{proposition}[Rare columns kept independent]\label{prop:rare}
Let column $X$ take a value other than its mode with probability $q$. For any $Y$, $\mathrm{TV}(P_{XY},P_X\otimes P_Y)\le q(2-q)\le2q$. Hence sampling $X$ independently of the other columns (given the segment) when fewer than $m$ rows are non-modal ($\hat q<m/n$, default $m=5$) changes every joint law involving $X$ by at most $2m/n$ in total variation, while the copula correlations of such a column are not identified by the data.
\end{proposition}
\begin{proof}
$\mathrm{TV}=\frac12\sum_xP(x)\|P_{Y\mid x}-P_Y\|_1$. For the mode $x_0$, $P_{Y\mid x_0}-P_Y=\sum_{x\ne x_0}P(x)(P_{Y\mid x_0}-P_{Y\mid x})$ has norm at most $2q$; for the other values the norm is at most 2. So $\mathrm{TV}\le\frac12[(1-q)2q+2q]=q(2-q)$.
\end{proof}
\noindent\emph{Check C34}: the resubmission flag of S8 (3 rows) is sampled independently; $\hat R=0.99$.

\begin{proposition}[Posterior predictive draws]\label{prop:ppc}
Given the fitted partitions, \texttt{redraw()} draws all structural parameters (segment weights, block, occupancy and activity probabilities, multiplicity rates) from their conditional posteriors; sampling after each redraw gives draws from the posterior predictive distribution of table sizes and link counts conditional on the partitions. Posterior predictive $p$-values are conservative (concentrated near $\tfrac12$ under the model), so extreme values signal misfit (Gelman, Meng and Stern 1996).
\end{proposition}
\begin{proof}
Composition: $\theta\sim p(\theta\mid\text{data},z)$ then $D^*\sim p(D\mid\theta,z)$ gives $D^*\sim p(D\mid\text{data},z)$.
\end{proof}
\noindent\emph{Check C33}: over 20 redraws of S8 the number of values has s.d.\ 498 and $p$-values 0.35 and 0.55 for values and reports.

\subsection{Privacy on coarse tables}
\begin{proposition}[Distances with an atom at zero]\label{prop:dcr}
When the real rows duplicate each other (a coarse table: a few low-cardinality columns), the leave-one-out distance distribution is $a_r\delta_0+(1-a_r)G_r$ with $a_r>0.05$, its 5th percentile is 0 and the ratio of 5th percentiles is undefined. Writing the synthetic distances as $a_s\delta_0+(1-a_s)G_s$, the pair of tests (i) exact-match rate $a_s\le\max(0.01,2a_r)$ and (ii) 5th-percentile ratio of $G_s$ to $G_r$ at least 0.5 compares both parts of the mixture and reduces to the original test when $a_r$ is small.
\end{proposition}
\begin{proof}
A mixture law is determined by its weight and its components; each test compares one of them. For $a_r<0.05$ the 5th percentile is attained in $G_r$ and the original ratio applies.
\end{proof}
\noindent The earlier code compared synthetic distances including the atom with the real positive part only, and so failed every coarse table (S8 reports: ratio 0 although the exact-match rate was within its limit).

\subsection{Defaults}
$\pi_s=0.9$ for occupancy, activity, null-rate and category-level priors, i.e.\ a combination that is absent from the data is a priori more likely to be forbidden than rare (\texttt{structure.structural\_prior}). Copula mode \texttt{rank} with chain extension up to 600 iterations, $\hat R$ target 1.1. Rare columns: $m=5$. Grain inference in \texttt{introspect}: a pair of foreign keys without duplicates where the chance-expected number of repeats $n(n-1)/(2\cdot\text{space})$ is at least 10.

%=====================================================================
\section{Small operational databases (v0.4)}\label{sec:v04}
The second real database was a small operational schema (11 tables, 15{,}607 rows: artists, albums, tracks, genres, media types, playlists, customers, employees, invoices and invoice lines; the catalog declares every key). It exercises situations that neither the fake scenarios nor the reporting model did: counts per parent that sit on a few fixed values (every customer has exactly 7 invoices; invoices have 1, 2, 4, 6, 9 or 14 lines), hierarchies of categorical columns (city, state, country), a hierarchy of 8 people of whom only 3 serve customers, price points (0.99, 1.99), long-tail labels (composers), a genre list just above the stratification limit, empty reference rows (playlists without tracks) and parents that always have children (every album has a track). Scenario S9 (a small shop with known truth) reproduces each of them; checks C35--C41 test the results below on it.

\subsection{Rows per parent}
\begin{proposition}[Count-model evidence]\label{prop:countbf}
Let $t_1,\dots,t_n$ be the counts of one segment. Under $\mathrm{DP}(a,G_0)$ with $G_0=\mathrm{Poisson}(\bar t)$ the marginal is the P\'olya-urn product
\[ p_{\mathrm{DP}}(t_{1:n})=\prod_{i=1}^n\frac{a\,g_0(t_i)+\#\{j<i:t_j=t_i\}}{a+i-1}, \]
which does not depend on the order of the counts. The engine compares $\log p_{\mathrm{DP}}-\tfrac12\log n$ with the Schwarz scores of Poisson ($\tfrac12\log n$) and negative binomial ($\log n$) at their maximum likelihood, summed over segments, and uses the DP (posterior predictive: Bayesian bootstrap of the observed counts with mass $a/(n+a)$ on $G_0$) when the DP wins by a factor of at least 100. (a) If the truth $F$ has finite support and is not in the parametric family, $\tfrac1n\log\mathrm{BF}\to\mathrm{KL}(F\|G^\star)>0$ for the best parametric $G^\star$, so the DP is chosen with probability $\to1$. (b) Under a parametric truth every first occurrence of a value costs the DP $\log\{(a+i-1)/(a\,g_0(t_i))\}$, i.e.\ about $\log n$ per distinct value, while later occurrences approach the empirical frequency; the parametric score pays a fixed $\tfrac d2\log n$. The DP is then chosen only when the counts repeat a small set of values far more than any Poisson or NB mixture allows.
\end{proposition}
\begin{proof}
The urn product is the Blackwell--MacQueen joint law of an exchangeable sequence, hence symmetric. (a) The urn predictive converges to $F$ (Ferguson 1973), so $\tfrac1n\log p_{\mathrm{DP}}\to\mathbb E_F\log F$; the parametric maximum likelihood gives $\tfrac1n\log\to\mathbb E_F\log G^\star$; penalties are $O(\log n)$. (b) Term by term from the urn product.
\end{proof}
\noindent\emph{Check C35}: in 70 Poisson, NB and two-segment data sets the DP is never chosen; in 20 sets with counts on $\{1,2,4,6,9\}$ or constant it always is; the marginal is invariant under permutation. On the real database: lines per invoice $\log\mathrm{BF}=254$, invoices per customer 101 (DP chosen); albums per artist $-17$, tracks per album $-92$ (not chosen).

\begin{proposition}[Counts conditioned on a segment total]\label{prop:condcounts}
Let $M$ parents of one segment receive $R$ children in total (from the block model, which keeps how child segments choose parent segments). Drawing the parents' counts from $\pi(N)\propto\prod_i f(N_i)\,\mathbf 1\{\sum_iN_i=R\}$, with $f$ the DP predictive of Proposition~\ref{prop:countbf}, by repeatedly picking random disjoint pairs $(i,j)$ and redrawing $N_i$ from $f(x)f(s-x)$, $x=0..s$, with $s=N_i+N_j$, leaves $\pi$ invariant and is irreducible; the chain converges to $\pi$ from any start. When $f$ is Poisson, $\pi$ is the multinomial $(R;1/M,\dots,1/M)$. For a fact with two parents, drawing rows per anchor from the DP predictive and then each row's partner segment with probability $\propto\eta^*_{kl}W_{2l}$ changes only the law of the anchor total: under the Poisson block model, given the total, the split over partner segments is exactly this multinomial (Proposition~\ref{prop:split}), and the scale rule is kept.
\end{proposition}
\begin{proof}
$\pi(N_i,N_j\mid N_{-ij})\propto f(N_i)f(N_j)\mathbf 1\{N_i+N_j=s\}$ is the exact full conditional of the pair, so each update is a Gibbs step. $f>0$ on $0..R$ because of the mass $a/(n+a)$ on the Poisson base, so any allocation reaches any other by moving one unit at a time between pairs. Independent Poisson variables with equal means, conditioned on their sum, are multinomial with equal cells.
\end{proof}
\noindent\emph{Check C36}: with a Poisson $f$, $N_1$ over 4000 draws has TVD 0.019 from Binomial$(10,\tfrac14)$; every draw sums to $R$; with observed counts all 7, every count is 7.

\begin{proposition}[Mandatory participation]\label{prop:mandatory}
If every one of $n$ parents has at least one child while the fitted Poisson rates leave $E_0=\sum_ie^{-\mu_i}\ge3$ parents empty in expectation, the Poisson model gives the observed ``no empty parent'' probability $\prod_i(1-e^{-\mu_i})\le e^{-E_0}\le e^{-3}\approx0.05$. The engine then uses the shifted model $N=1+M$: for links between entities, each parent of a segment first receives one child (while the segment's rows last), the rest follow the usual weights; for facts, rows per anchor are $1+\mathrm{Poisson}(\max(\mu-1,0))$, which keeps the mean.
\end{proposition}
\begin{proof}
$1-x\le e^{-x}$. Given a total $R\ge M$, assigning one child to each parent and the remaining $R-M$ by the multinomial of the weights is the conditional law of $1+M_i$ with $M$ multinomial.
\end{proof}

\begin{proposition}[Copied anchors keep their activity]\label{prop:copiedactive}
A copied table reappears row for row, each row its own segment. Whether such a row takes part in an occupancy fact at all (a playlist with no tracks) is a property of that known unit, so its activity is set to the observed indicator rather than drawn from a Beta posterior based on one observation; its pairs are still drawn from the block probabilities.
\end{proposition}
\begin{proof}
Conditioning on the unit's identity: the posterior predictive of an observed unit's own fixed attribute is a point mass at its observed value. A Beta$(a+A,b+1-A)$ draw from a single observation would make an empty playlist active with probability about $a/(a+b+1)$, and its tracks would then come from the prior (the shop data: 2140 rows in a playlist that has none).
\end{proof}

\subsection{Hierarchies of categorical columns}
\begin{proposition}[Dependency groups]\label{prop:fd}
For categorical columns $a,b$ (null a level of its own) let $M_0$: $b$ i.i.d.\ Dirichlet(1)-categorical, and $M_1$: the first row of each level of $a$ draws $b$ from the same law and every later row repeats that value with probability $1-\varepsilon$ ($\varepsilon=10^{-3}$). (i) If $a$ has no repeated level, $\log\mathrm{BF}=0$, so a near-unique column never determines another. (ii) If $b$ is independent of $a$, each repeated row agrees with probability $\sum_lp_l^2<1$ and $\log\mathrm{BF}\to-\infty$ linearly in the number of repeats. (iii) Columns linked by decisive dependencies are joined into one categorical over their observed combinations only if the combinations are no more than the levels of the finest member (plus $\lceil10^{-3}n\rceil$ exceptions). Then the joint releases the information of that one column, every synthetic row satisfies the dependency, and no unseen combination (a city in the wrong country) can be generated.
\end{proposition}
\begin{proof}
(i) Without repeats $M_1$ and $M_0$ assign the same Dirichlet-multinomial marginal to the same values. (ii) Each repeat contributes $\log(1-\varepsilon)-\log p_{b}$ when it agrees and $\log\{\varepsilon/(L-1)\}-\log p_b$ otherwise; under independence the expected contribution is $\sum_lp_l\{p_l\log(1-\varepsilon)+(1-p_l)\log(\varepsilon/(L-1))\}-\sum_lp_l\log p_l$, which is negative for $\varepsilon\le10^{-3}$ unless $b$ is nearly constant. (iii) The joint codec's levels are the observed combinations, in bijection with the finest member's levels when the dependencies are exact; its Dirichlet posterior puts mass only on them.
\end{proof}
\noindent\emph{Check C37}: on the shop data the group (city, state, country) is found, every synthetic customer's state and country follow its city, the independent tier stays out, $\log\mathrm{BF}=-4851$ for independent columns and $0$ for a unique determinant. Distance-to-closest-record keeps one column per group (the others add no information), so a table whose only attributes are one geographic hierarchy is not flagged as copied.

\subsection{Tiny parents}
\begin{proposition}[Strata from an own column]\label{prop:attrstrata}
An entity with fewer than $2\times$\texttt{min\_rows\_per\_segment} rows has one latent segment, so it cannot learn which of its rows receive children. For such an entity, each categorical column $c$ with $2\le L_c<n$ levels is scored by $\sum_r\max\{0,\log\mathrm{BF}_r(c)\}-\log(\#\text{candidates}\times\#\text{relations})$, where $\mathrm{BF}_r$ compares a Poisson--Gamma model of children per row in relation $r$ with one rate per level of $c$ against one pooled rate (Gamma$(\tfrac12,\tfrac12/\bar y)$ prior). If the best score exceeds $\log100$, the levels of $c$ become strata with one segment each, $c$ is set from the segment, and the table is generated with its observed stratum counts (scaled). Children then go only to strata that had children, up to the structural prior (Proposition~\ref{prop:zoib}), and generating with fixed stratum counts is exact conditioning on the sufficient statistic of the composition.
\end{proposition}
\begin{proof}
The Poisson--Gamma marginal is $\frac{b^a}{\Gamma(a)}\frac{\Gamma(a+\sum y)}{(b+n)^{a+\sum y}}\prod\frac1{y!}$; the BF is a ratio of such marginals. A relation the levels do not explain is given score 0, not a negative one, because the stratified model contains the pooled one for that relation (equal rates across strata). Given the counts, all arrangements of an exchangeable categorical sample are equally likely, so a uniform permutation of the scaled counts is the conditional law.
\end{proof}

\begin{proposition}[Hierarchy of a given size]\label{prop:treesize}
Every non-hierarchical entity is generated with exactly $\mathrm{round}(s\,n)$ rows; a multi-type Galton--Watson forest has a random size whose relative spread is large when the forest has few roots (one root: the size is a geometric-like variable). Repeating the branching process until $|T|=\mathrm{round}(s\,n)$ (within 1\% for large forests, and with the stratum counts of Proposition~\ref{prop:attrstrata} for tiny ones) draws from the process conditioned on its size, by rejection.
\end{proposition}
\begin{proof}
Rejection with acceptance indicator $\mathbf 1\{|T|=t\}$ samples the conditional law (Proposition~\ref{prop:rules}); the expected number of attempts is $1/P(|T|=t)$. When no attempt in the budget is accepted, the closest draw is used and the deviation is reported.
\end{proof}
\noindent\emph{Check C38}: the shop's staff (8 rows) is stratified by office (agents' branch), all synthetic customers are served by sales agents, the staff has 8 rows with the source's stratum counts, one root and no cycle. \emph{Check C10} still compares the size formula with the unconditioned process.

\subsection{Stratification by larger code lists}
\begin{proposition}[Code lists that explain links]\label{prop:partner}
A copied parent with more levels than \texttt{stratify\_max\_levels} is used as a stratifier when its levels explain how many children the entity's rows have (Poisson--Gamma, as above) or which partners they meet in two-parent facts: a symmetric Dirichlet-multinomial over partner rows per level against one pooled, each relation contributing $\max\{0,\log\mathrm{BF}\}$, minus $\log(\#\text{relations})$, above $\log100$. Segments are nested in the strata, so a block of a fact that pairs a partner row with a segment of one genre never generates tracks of another genre.
\end{proposition}
\begin{proof}
The Dirichlet-multinomial marginal is exact for exchangeable partner draws within a level; the comparison is a ratio of marginals. Nesting: Proposition~\ref{prop:strata}.
\end{proof}
\noindent\emph{Check C39}: the shop's playlists hold only tracks of their own genres; with genre (30 levels) as stratum the share of synthetic playlist rows with a genre the playlist never has is at most 2\% (82\% before).

\subsection{Values}
\begin{proposition}[Tail levels and few-valued numbers]\label{prop:tail}
(a) A categorical with more than \texttt{max\_levels} levels models the rarest as one level. Decoding that level by a Bayesian bootstrap over the tail levels seen in at least $k$ rows (default 5) gives each of them the marginal of the full Dirichlet posterior restricted to those levels; levels seen in fewer than $k$ rows are never released, and the placeholder is never written. (b) An integer or decimal column with at most 30 distinct values is modelled as an ordered categorical on those values, so no value between them is generated.
\end{proposition}
\begin{proof}
(a) Aggregation property: if $(p_1,\dots,p_K,p_O)\sim\mathrm{Dir}(\alpha_1,\dots,\alpha_K,\sum_{j\in O}\alpha_j)$ and, independently, $w\sim\mathrm{Dir}(\alpha_j,j\in O)$, then $(p_1,\dots,p_K,p_Ow)\sim\mathrm{Dir}(\alpha_1,\dots,\alpha_K,\alpha_j,j\in O)$. (b) A continuous Bayesian-bootstrap marginal interpolates its quantile knots, so it puts positive mass between two atoms (the shop's prices came out as 1.54 or 4.26); a categorical marginal is supported on its levels.
\end{proof}
\noindent Aggregates are rounded to the precision of the stored column (a sum of prices is stored to the cent); without this, totals $5.9399\ldots$ and $5.94$ count as different values in every comparison. \emph{Checks C39, C41}: lines per invoice stay on $\{1,2,4,6,9\}$, every customer has 5 invoices, totals equal the sum of their lines to the cent, prices are 0.99 or 1.99, the empty playlist stays empty, the shop scorecard passes; a categorical's tail decodes only to levels seen in at least 5 rows.

\paragraph{Labels.} The text model removes every token seen in fewer than $k$ rows (Section~\ref{sec:text}). For names and titles most words are unique to their row; when fewer than half of the rows keep a word after that removal, the model could only recombine punctuation and filler words (``, \&''), so it generates surrogate labels (``Artist 04217'') instead; otherwise generated texts must contain a word and may not open or close on punctuation or leave a bracket open (a rejection step, as in Proposition~\ref{prop:rules}). Identifiers without a format shared by 5 rows (company names, street addresses) are surrogate labels too, and person-name columns give first or last names according to the column. \emph{Check C41}.

\begin{proposition}[Evidence for a grain]\label{prop:grainev}
Let a table have $n$ rows whose columns $(A,B)$ are drawn independently with level probabilities $p_a$, $q_b$. The number of row pairs that agree on $(A,B)$ has mean $\lambda=\binom n2\sum_a p_a^2\sum_b q_b^2$, and $\hat\lambda=\binom n2\,\hat S_A\hat S_B$ with $\hat S_A=\sum_a c_a(c_a-1)/(n(n-1))$ estimates it; each factor is unbiased for its sum of squares. When the probabilities are small, the count is approximately Poisson($\lambda$) (Chen--Stein), so a table with no repeated pair gives a likelihood ratio of about $e^{\lambda}$ for ``each pair occurs at most once'' against independence. The doctor proposes a unique key only when $\hat\lambda\ge10$ ($e^{-10}<5\cdot10^{-5}$). Below that, regular counts are reported as a property of the source and nothing is changed.
\end{proposition}
\begin{proof}
$\mathbb E[c_a(c_a-1)]=n(n-1)p_a^2$ for a multinomial count, which gives unbiasedness. Two rows agree with probability $\sum_{a,b}(p_aq_b)^2=\sum p_a^2\sum q_b^2$, so linearity over the $\binom n2$ pairs gives $\lambda$. The Poisson limit for the number of matching pairs (the birthday problem) is Theorem~1 of \cite{agg}: the total variation distance is at most $\min(1,1/\lambda)\,(b_1+b_2)$, and both terms vanish when $\max_{a,b}p_aq_b\to0$. \emph{Check C42} (pair count within 0.3\% of Monte Carlo, plug-in within 0.2\%, $P(0)$ within 0.01 of $e^{-\lambda}$).
\end{proof}
On the music store, lines of an invoice never repeat a track, but with 2240 lines spread over 3503 tracks only $\hat\lambda=0.9$ repeats were expected, so this is not evidence of a rule; on the playlists, $\hat\lambda=2241$ and the grain is real.

\subsection{Introspection}
Type inference now uses the table's context: a bare \texttt{name} is a person's name only in a table with personal columns; near-unique labels of a small table that is not a party (genres, media types) are reference labels; near-unique labels of a party table (reporters, customers, staff: a lexicon of party nouns) are identifiers; other near-unique or long-tail multi-word labels (titles, composers: at least 20\% of rows hold values seen fewer than 5 times) are short text. Columns equal to a parent's column on 99.9\% of rows (null equal to null) become \texttt{lookup} columns; parent columns equal to a count, a sum or a sum of products of child columns become \texttt{aggregate} rules; the anchor of a fact is the parent whose name the fact extends and whose keys it covers. A column named like a key whose parent the names do not reveal is listed as a question (surrogate keys of most tables share the range $1..n$, so values alone cannot decide). \emph{Check C40}.

\subsection{Defaults (v0.4)}
\texttt{facts.count\_model\_bf} 100; \texttt{marginals.fd\_bf} 100; \texttt{segments.attribute\_strata} \texttt{auto} (a table name to column map forces one); mandatory participation at $E_0\ge3$; tree size tolerance 1\%.

%=====================================================================
\section{Data warehouses with order headers (v0.5)}\label{sec:v05}
The third real database was a sales data warehouse (AdventureWorks DW: 27 tables, 1.05 million rows; composite keys, a composite foreign key into a fact, order numbers stored on the lines, three role-playing date keys per order, amounts that are identities of other columns, a running stock balance). Its first two runs exposed 38 findings (report: \texttt{realdata/adventureworks/README.md}). This section states the model changes that remove their causes and the properties they rely on. Scenario S10 (a sales warehouse with known truth) reproduces each structure; checks C43--C45 test them.

\subsection{Identifying the data model}
A drafted model is a set of dependencies found in samples: unique column combinations (keys), inclusion dependencies (links), functional dependencies (lookups, order headers), algebraic constraints (derive rules) and order dependencies (running balances). Each candidate is accepted only by the following test.

\begin{proposition}[Significance of a discovered dependency]\label{prop:signif}
Let a candidate dependency hold on $s$ of $n$ eligible sampled rows, and let $p_0$ be the probability that it holds on one row by chance: the collision probability $\sum_v p_v^2$ of the column for an equality with another column, $\sum_v p_v^{m}$ over the value shares of a determined column for a functional dependency checked on groups of $m$ rows, and the agreement rate with a randomly permuted input for algebraic and cumulative rules. Introspect accepts the candidate iff $n\ge30$, $s/n\ge0.999$ and $P(\mathrm{Bin}(n,p_0)\ge s)\le\alpha=10^{-9}$. With $M$ candidates tested, the probability of accepting any dependency that holds only by chance is at most $M\alpha$ (Bonferroni); on AdventureWorks $M<10^5$, so at most $10^{-4}$. A rule that is true on the whole table is accepted with probability $\to1$ as soon as $n\,\mathrm{KL}(1\|p_0)=-n\log p_0>\log(1/\alpha)+O(\log n)$, i.e.\ about $21/(-\log p_0)$ rows.
\end{proposition}
\begin{proof}
Under the null each row agrees independently with probability at most $p_0$, so the binomial tail bounds the false-acceptance probability of one test; the union bound gives $M\alpha$. If the rule holds on every row, $s=n$ and the tail is $p_0^n$, which is $\le\alpha$ iff $n\ge\log(1/\alpha)/(-\log p_0)$.
\end{proof}
\noindent The fixed threshold used before (chance share below 0.5) rejected true rules on skewed columns (89\% of reseller lines carry no discount, so the discount-rate lookup has $p_0=0.89$; with $n=20{,}000$ the tail is $10^{-1000}$) and accepted chance agreements on small samples. Samples are simple random samples drawn with a seeded generator from the table's primary key (Section~\ref{sec:v05}, C44: two runs give the same spec); groups (orders, products) are read whole as cluster samples for within-group rules. Personal and free-text columns are candidates for lookups too: a billing address equal to the customer's address on every row is a functional dependency like any other, and generating it as a lookup copies the synthetic customer's generated address, which keeps both the identity and the privacy (an earlier v0.5 draft excluded them and lost this rule on Chinook).

\begin{proposition}[Order headers: lossless split]\label{prop:header}
Let a fact $F$ with a document number $d$ satisfy the functional dependency $d\to H$ for a set of header columns $H$. Then $F=\pi_{d,H}(F)\bowtie\pi_{d,L}(F)$ for the remaining columns $L$ (Heath 1971). Generating $\pi_{d,H}$ as an entity (one row per document, with its own links) and the lines as a fact anchored on it, then joining, gives synthetic data in which every document has exactly one header and line numbers $1..n$ by construction; the scale rule applies to documents per customer and lines per document.
\end{proposition}
\begin{proof}
Heath's theorem: a decomposition along a functional dependency whose left side is a key of one component is lossless. The synthetic header table has one row per document by construction, so the join repeats it on each line.
\end{proof}

\begin{proposition}[Offset keys]\label{prop:offset}
Let $b$ and $k$ be keys into one parent whose order column $o$ is injective (a calendar). The map $(b,k)\mapsto(b,\,o(k)-o(b))$ is a bijection from pairs of parent rows onto pairs (row, attainable gap). Modelling the gap as an attribute therefore loses nothing, keeps the gap distribution (ship $=$ order $+7$ days on 100\% of rows) exactly, and a draw whose $o(b)+\text{gap}$ is not an order value of the parent is outside the support of the source; redrawing it (Proposition~\ref{prop:refill}) gives the model conditioned on the support.
\end{proposition}
\begin{proof}
Given $b$, $k\mapsto o(k)-o(b)$ is injective because $o$ is; its inverse is the lookup of $o(b)+\text{gap}$.
\end{proof}

\subsection{Generating dependent records}
\begin{proposition}[Redraw instead of delete]\label{prop:refill}
Let the rows of a fact be independent given their anchors, row $i\sim Q(\cdot\mid a_i)$, and let $V$ be the set of valid rows (rules, unique keys, offsets landing on a parent row). The law that keeps the number of rows per anchor and puts all mass on valid rows, closest to $Q$ in Kullback--Leibler divergence, is $\prod_iQ(\cdot\mid a_i,V)$, and redrawing each invalid row from $Q(\cdot\mid a_i)$ until it is valid samples it exactly. Deleting invalid rows instead conditions the counts as well: a customer's expected number of kept rows is multiplied by $Q(V\mid a_i)$.
\end{proposition}
\begin{proof}
For a product reference measure and the constraint that each factor lives on $V$, the I-projection factorises and each factor is $Q(\cdot\mid a_i)$ restricted to $V$ and renormalised (Csisz\'ar 1975). Rejection sampling from $Q(\cdot\mid a_i)$ returns exactly that restriction. Under deletion the number of kept rows of anchor $i$ is binomial with success probability $Q(V\mid a_i)$.
\end{proof}
\noindent On AdventureWorks deletion removed 57{,}750 of 60{,}398 internet sales (a unique key over an order number drawn as a 200-level category) and 46{,}224 of 60{,}855 reseller sales (date checks over independently drawn keys). Both causes are now modelled (Propositions~\ref{prop:header}, \ref{prop:offset}); leftovers after the last round are deleted, counted and scored (\texttt{rows\_removed\_share} $\le0.01$, doctor D23).

\begin{proposition}[Shares of copied rows]\label{prop:shares}
Rows $d$ of a copied table in segment $k$ receive $c_d$ links. Model their rates as $\lambda_d\sim\mathrm{Gamma}(a_k,\beta_k)$, $c_d\mid\lambda_d\sim\mathrm{Pois}(\lambda_d)$, with $a_k$ estimated by moments, $a_k=\bar c^2/(s^2-\bar c)$ when $s^2>\bar c$ and $a_k=\infty$ otherwise. Drawing each row's share from the posterior $\mathrm{Gamma}(c_d+a_k)$ and the synthetic counts from the resulting Poisson gives counts whose marginal law over rows equals the source's (negative binomial with the fitted $a_k$). The previous plug-in, shares $\propto\mathrm{Gamma}(c_d+1/m_k)$, adds the source's Poisson noise a second time: on a calendar with 4.6 orders per day the variance/mean of orders per day became about 2 instead of 1.
\end{proposition}
\begin{proof}
If $\lambda\sim\pi$ and $c\mid\lambda\sim p(c\mid\lambda)$, then a draw $\lambda^*\sim\pi(\cdot\mid c)$ is again distributed as $\pi$ (the data-averaged posterior equals the prior), so $c^*\mid\lambda^*$ has the marginal law of $c$. With $a=1/m$ the posterior mean of $\lambda$ is about $c$, and $c^*\sim\mathrm{Pois}(\approx c)$ has variance $\mathrm{Var}(c)+\mathbb E c$.
\end{proof}

\begin{proposition}[Per-row segments need data per row]\label{prop:perrow}
If a copied table with $m$ rows keeps one segment per row, every segment $k$ of an entity linked to it carries $m$ free rates $\eta_{k\cdot}$. Splitting one segment into two adds $m$ rates; under the Gamma prior each costs an Occam factor of order $\tfrac12\log(1+n_{kd}/\beta)$ in the collapsed marginal, so a split that the data support by a log-likelihood gain $G$ is rejected when $G<\tfrac m2\overline{\log(1+n_{kd}/\beta)}$. With few links per row this blocks every split of the other side (S10: the products collapsed from four categories to two segments, and the price--category and buyer--category dependences were lost). Rows are therefore kept as segments only when each relation carries at least \texttt{min\_rows\_per\_segment} links per row; otherwise the copied table is clustered like any entity and its rows keep their own shares (Proposition~\ref{prop:shares}): a degree-corrected block model (Karrer and Newman 2011), whose cost per extra segment is $K_{\mathrm{copied}}$ rates instead of $m$.
\end{proposition}
\begin{proof}
Laplace approximation of the Gamma--Poisson marginal of one block: $\log p(C\mid E)=\log p(C\mid\hat\eta)-\tfrac12\log(1+E/\beta)+O(1)$.
\end{proof}

\begin{proposition}[Tied header segments]\label{prop:tied}
Let the header entity $H$ take its parent's segment, $z_j=z_{p(j)}$. The Gibbs full conditional of a parent row $i$ is
\[ P(z_i=k\mid\cdot)\propto\pi_k\,f(x_i\mid k)\prod_{r\ni i}L_r(i,k)\;\mathrm{Pois}(m_i\mid\mu_k)\prod_{j:p(j)=i}f(x_j\mid k)\prod_{r\ni j}L_r(j,k), \]
which the sampler evaluates by summing the children's row log-likelihoods into their parent's row. The collapsed merge score of a parent partition adds the children's attribute and relation marginals under the induced partition and no partition prior of their own. The model of lines given the customer is then the one the unsplit fact had (rates per customer segment and product segment), while documents keep their own header columns and line counts.
\end{proposition}
\begin{proof}
The children's segments are a deterministic function of $z_i$, so changing $z_i$ changes every factor that involves $z_j$, $p(j)=i$; all other factors do not depend on $z_i$. The partition of the children is not a free variable, so it has no prior.
\end{proof}
\noindent Without the tie the header segments collapsed to one or two (about two lines per order give little evidence per order) and product choice stopped depending on the customer (S10: share of bikes in segment C orders 0.24 instead of 0.03; with the tie 0.07; check C43).

\begin{proposition}[Running balances]\label{prop:cumulative}
A cumulative rule $y_t=y_{t-1}+s_t$ within groups is generated from the steps. If every source group opens at the same value ($y_1-s_1=c$ on 99.9\% of groups) the synthetic groups open at $c$. If the source balance is never negative and the step subtracts a column ($s=a-b$), the subtracted column is reduced to $\min(b,\,y_{t-1}+a-\text{floor})$, the rule of an account that cannot pay out more than it holds; then $y_t\ge$ floor on every row and the identity still holds.
\end{proposition}
\begin{proof}
Induction over $t$: if $y_{t-1}\ge$ floor and $y_{t-1}+a-b<$ floor, the reduced $b'=y_{t-1}+a-\text{floor}\in[0,b)$ gives $y_t=$ floor.
\end{proof}

\subsection{Run time of the segment merge}
\begin{proposition}[Merge from block statistics]\label{prop:mergestats}
Every term of the collapsed merge score (attributes, Poisson and occupancy blocks, counts, trees, tied children) depends on the partition $z$ of the merged entity only through statistics that are sums over rows within segments. Merging segments $a,b$ adds row $b$ of each statistic to row $a$ along the axes indexed by that entity. Scoring a candidate therefore costs $O(K\cdot K_{\mathrm{other}})$ per relation instead of a pass over its rows, and gives the same value (checked to $2\times10^{-16}$ relative on S4, S7--S10).
\end{proposition}
\begin{proof}
Block counts, sums of measures and segment sizes are additive over the rows of a segment; the remaining factors (e.g.\ $E=n_1n_2^\top$, $T=A\,m_2^\top$) are products of additive statistics and are formed after merging.
\end{proof}
\begin{proposition}[Merge gain from two slices]\label{prop:mergegain}
Call a term of the merge score separable in entity $e$ if it equals $\sum_k g(S_{k\cdot})+h(\sum_k S_{k\cdot})$, where $S_{k\cdot}$ is the slice of its statistics for segment $k$ of $e$. Attribute terms, segment sizes, link, single-count, tied-count and occupancy terms (without a per-stratum prior) are separable; tree terms and links of an entity to itself are not. For a separable term the gain of merging $a$ and $b$ is $g(S_{a\cdot}+S_{b\cdot})-g(S_{a\cdot})-g(S_{b\cdot})$, i.e.\ the term evaluated on the two slices merged minus on the two slices apart, at cost $O(K_{\mathrm{other}})$ instead of $O(K\,K_{\mathrm{other}})$.
\end{proposition}
\begin{proof}
Merging changes only rows $a,b$ of each sliced statistic, and leaves their column sums, hence $h$, unchanged.
\end{proof}
\noindent Checked against full rescoring on 537 candidates of five scenarios and every entity of AdventureWorks (largest relative difference $1.5\times10^{-10}$, same best pair everywhere). On AdventureWorks the merge took more than 8 minutes when rescoring over rows (stopped), 276 s from block statistics, and 84 s with two-slice gains; the customer table carries 709 segments because it is stratified by its 655 geographies.

\subsection{Structural zeros, mandatory participation and hierarchies}
\begin{proposition}[Structural zeros in an occupancy fact]\label{prop:structzero}
Let a nested fact's key require that a partner column equals a column of its anchor (a cell's template equals the filing's template). Pairs with different values are impossible. In the occupancy model put group codes $g_i$ on anchors and $h_j$ on partners and replace every count of pairs at risk by eligible pairs: $T_{kl}=\sum_g N^{(1)}_{kg}N^{(2)}_{lg}$, and the non-present pairs of row $i$ by the partners of its own group. This is the quasi-independence model of a table with structural zeros (Goodman 1968). Sampling only eligible pairs then needs no deletion, and $\mathbb E$[rows per anchor] equals the source's within each block, whereas the full-grid model underestimates $p_{kl}$ by the share of ineligible pairs and loses those rows when they are deleted.
\end{proposition}
\begin{proof}
Under the full-grid model $\hat p_{kl}=E_{kl}/(A_km_l)$, while only $T_{kl}\le A_km_l$ pairs can occur; drawing on the eligible pairs with $\hat p=E/T$ gives $\mathbb E=E$ in every block. The statistics $N^{(1)}$, $N^{(2)}$ are additive over segments, so Proposition~\ref{prop:mergestats} still applies.
\end{proof}
\noindent On the bi schema 784--1{,}088 reported values (1.2--1.6\%) were deleted in every earlier version; now none are.

\begin{proposition}[Mandatory participation in an occupancy fact]\label{prop:occminone}
If every source anchor has at least one row (an order header always has lines) and the fitted model expects at least 3 anchors without pairs, $\sum_i\prod_l(1-p_{z_il})^{m_l}\ge3$, the pairs of each active anchor are drawn conditioned on at least one being present, by redrawing empty anchors (rejection) and, for anchors still empty after 20 rounds, by the sequential method below; both are exact for the zero-truncated law. The rates are then fitted for that law: in segment $k$ the free rates ($0<p_{kl}<1$) are scaled by the $c_k$ that solves
\[\sum_{i:z_i=k}\frac{\mu_i(c_k)}{1-P_{0,i}(c_k)}=\sum_{i:z_i=k}n_i,\qquad \mu_i(c)=\sum_l c\,p_{kl}m_l,\quad P_{0,i}(c)=\prod_l(1-c\,p_{kl})^{m_l},\]
with $n_i$ the observed pairs of anchor $i$ ($m_l$ replaced by the anchor's eligible partners under Proposition~\ref{prop:structzero}).
\end{proposition}
\begin{proof}
For any count $N\ge0$, $\mathbb E[N\mid N\ge1]=\mathbb E[N\,\mathbf 1\{N\ge1\}]/P(N\ge1)=\mathbb E[N]/(1-P_0)$, exactly for the Poisson-binomial count of an anchor's pairs. The left side increases in $c$ (the zero-truncated mean of a sum of Bernoullis increases with each success probability), so the root is unique and found by bisection. This is the method-of-moments estimator of the truncated model; in the Poisson limit it is the maximum-likelihood estimator of the zero-truncated Poisson. Rates fitted without it give $\mathbb E[N\mid N\ge1]=\bar n/(1-P_0)$, too high by the factor $1/(1-P_0)$.
\end{proof}
\noindent\emph{Sequential draw.} With the partner blocks in a fixed order and $P_{0l}=(1-p_l)^{m_l}$, the first non-empty block is $l^*$ with probability $(1-P_{0l^*})\prod_{l<l^*}P_{0l}/(1-P_0)$, its count is binomial truncated at zero, later blocks are unconditioned binomials, and partners are uniform within blocks. This factorises the conditional law exactly, and its cost does not depend on $P_0$; rejection alone needs $1/(1-P_0)$ rounds, unbounded when an anchor segment always has exactly one pair (then $c_k\to0$ and the law tends to one pair with block probabilities $\propto p_lm_l$).
\noindent Without it about 15\% of AdventureWorks internet orders received no line and vanished (23{,}565 orders instead of 27{,}659, 2.54 lines per order instead of 2.18), and with them 12\% of the customers' only purchases. Conditioning alone restored the orders but gave 2.46 lines per order (11\% too many lines, 67{,}994 instead of 60{,}398): the fitted rates already averaged the source's 2.18 lines over orders that all had lines. With the rates fitted for the truncated law the run gives 27{,}659 orders, 2.18 lines per order and 60{,}250 lines (source 60{,}398); the scale factors range from 0.000 (segments of one-line orders, drawn by the sequential method) to 2.4 (a value above 1 means the posterior draw of the rates sat below the segment's observed mean). On simulated data with known rates (Check C43) the uncorrected rates overstate the pairs per anchor by 20--25\%, the corrected ones are within 4\% (sampling noise) over 16 seeds.

\begin{proposition}[Hierarchies generated level by level]\label{prop:treelevels}
Let a self-reference be a multi-type Galton--Watson process whose type is the node's segment and depth, with rates $\lambda_{kd}$ and child-segment probabilities $\psi_{kd}$. Given the numbers $n_{k'd}$ of nodes of each type at depth $d$, the parents of the depth-$d$ nodes of segment $k'$ are i.i.d.\ among the depth-$(d-1)$ nodes with probabilities $\propto w_i\lambda_{z_i,d-1}\psi_{z_i,d-1}(k')$ ($w_i$ the node's frailty). The generator therefore fixes $\mathrm{round}(s\,n_{kd})$ and draws the parents; size, depth profile and composition match the source at every scale.
\end{proposition}
\begin{proof}
Independent Poisson counts conditioned on their total are multinomial with probabilities proportional to their means; children of type $k'$ of node $i$ are Poisson with mean $w_i\lambda_{z_i,d-1}\psi_{z_i,d-1}(k')$, independently over $i$.
\end{proof}
\noindent The rates by level have the prior $\mathrm{Gamma}(\lambda_k,1)$ (one pseudo-node at the segment's rate). Before, the process was conditioned on its size only: on AdventureWorks a segment mixing managers and staff grew chains to depth 7 (source 4), and the 18-person sales branch (one vice-president, three managers, 14 representatives) vanished in most draws with one decision near the root; with it went the quotas of the sales staff (finding F38: 6 quota rows instead of 163). Level by level, three seeds give 17--18 sales representatives (source 17), and the scorecard of the employee table passes 58--59 of 59.

%=====================================================================
\section{Normalised and denormalised warehouses (v0.6)}\label{sec:v06}
The fourth real database was the Contoso Data Generator V2 release (100k customers; 93{,}470 orders, 223{,}974 order lines; a denormalised sales table next to the normalised orders and order lines; exchange rates by date and currency pair with no links; geography stored on the customer row; prices that are list prices times a factor). Its test (from v0.5, unedited drafts) exposed findings K01--K18 (report: \texttt{realdata/contoso/README.md}). This section states the model changes that remove their causes and the properties they rely on. Scenario S11 (a retail warehouse with known truth) reproduces each structure; checks C47--C56 test them. Regressions on AdventureWorks and S11 exposed seven more (R1--R7), also fixed here.

\subsection{Identifying the data model}
\begin{proposition}[Functional dependencies tested on rows]\label{prop:fdrows}
For columns $x,y$ on $n$ sampled rows let $s=\frac1n\sum_g\max_v n_{gv}$, the share of rows whose $y$ is the most common $y$ of their $x$-group $g$, and let $p_0$ be the mean of the same statistic over random permutations of $y$. Introspect accepts $x\to y$ iff $s\ge0.999$, the binomial test of Proposition~\ref{prop:signif} passes with $(s,n,p_0)$, and $\exp(-n(s-p_0)^2/2)\le\alpha$. Under independence of $x$ and $y$ the probability of acceptance is at most $\alpha$. Singleton groups agree under both hypotheses and add no evidence, so a near-unique $x$ is never accepted, and a code with few large groups (8 countries) is tested on its thousands of rows. The group-based test used before counted groups, so a dependency from 8 countries had $n=8<30$ and was never tested (the currency of the customer's country, K09).
\end{proposition}
\begin{proof}
Under independence every permutation of $y$ is equally likely given the data, so the permutation law is the exact null of $s$, with mean $p_0$. A transposition of two $y$ values changes $s$ by at most $2/n$; the bounded-differences inequality for random permutations (McDiarmid 1989) gives $P(s-p_0\ge t)\le\exp(-2t^2/(n(2/n)^2))=\exp(-nt^2/2)$. The binomial condition is kept as well (it is the stricter one when $p_0$ is close to 1).
\end{proof}

\begin{proposition}[Implied dimensions]\label{prop:dimension}
Let an entity $E$ satisfy $k\to Y$ for a code column $k$ and a set of columns $Y$. Then $E=\pi_{k,Y}(E)\bowtie\pi_{k,R}(E)$ (Heath 1971). Generating $D=\pi_{k,Y}(E)$ as a copied parent, $k$ as a link of $E$ to $D$ and $Y$ as lookups $k.y$ gives synthetic rows with $k\to Y$ exactly and every $(k,Y)$ combination an observed one, at any number of levels; the expected number of synthetic rows per $k$ follows the source degrees (Proposition~\ref{prop:shares}). If the dependency holds on a share $1-\varepsilon$ of the rows, using the first $Y$ of each $k$ changes the $Y$ values of at most a share $\varepsilon$ of the rows, so every marginal of $Y$ moves by at most $\varepsilon$ in total variation.
\end{proposition}
\begin{proof}
Heath's theorem as in Proposition~\ref{prop:header}; the lookup reads $D$'s row of each synthetic $k$, so $Y$ is a function of $k$. A row whose $Y$ differs from its group's first values is the only row whose value changes.
\end{proof}
\noindent\textbf{Anchors.} A code $k$ drawn as a link carries the dependence of the other attributes on $Y$ only through the segments. When an attribute depends on a coarse label of $k$ (latitude on the country, 8 countries against 565 areas) the segments, fitted to the whole entity, need not resolve it: Contoso's synthetic customers had the same mean latitude in every country (K16). An anchor $a\in Y$ (with $k\to a$) is kept as a modelled attribute and $k$ is drawn given it (the rule $k.a=a$ then holds on every row and is reported like any other).
\begin{proposition}[Anchored dimensions]\label{prop:anchor}
Let $a\in Y$ be the anchor and let the link $k$ be drawn from the link model restricted to $\{k: D[k].a=a\}$: a parent segment $t$ with probability $\propto\phi_{z t}\,W_t(a)$, where $W_t(a)$ is the share of segment $t$'s rows with anchor $a$ under the link's member law (degree weights, or uniform under a count model), then a member of $t$ with anchor $a$ by that law. Then (i) every synthetic row keeps $k\to Y$ and an observed $(k,Y)$ combination (Proposition~\ref{prop:dimension}); (ii) the joint law of $a$ with the other attributes is the attribute model's, as for any categorical; (iii) given $a$, $k$ has the link model's law restricted to the codes of $a$. If the attribute model reproduces the source law of $(a,X)$ and the link model reproduces $k\mid a$, the joint law of $(k,Y,X)$ is reproduced whenever $X\perp k\mid a$ in the source.
\end{proposition}
\begin{proof}
(i) $Y$ is still read from $D[k]$, so the anchor column is overwritten by $D[k].a$, which equals the modelled $a$ on every row. (ii) Only the link is drawn. (iii) $Q(k\mid z)=\sum_t\phi_{zt}\,\omega_t(k)$ with $\omega_t$ the member law of segment $t$; restricted to the codes of $a$ this is $\sum_t\phi_{zt}W_t(a)\,\omega_t(k)/W_t(a)$, which the two stages sample. When an attribute rule makes the row's attributes be redrawn, its link is drawn again from the new anchor. The last claim is the factorisation $P(k,X)=\sum_a P(a,X)P(k\mid a)$ under $X\perp k\mid a$, with $Y=D[k]$.
\end{proof}
\noindent The Metropolis--Hastings kernel used first left 289 of 105{,}000 Contoso customers outside their country's areas after 200 repair rounds, because a segment's link model can give the anchor's areas no weight; the two-stage draw has no rejection, and a row whose segment gives the anchor no weight takes its segment from the link's overall mixture (logged; none in Contoso). Introspect picks as anchor the finest categorical label among $Y$ with 2 to 50 levels (Country rather than Continent; SubCategoryName in Product), since a finer label keeps more of $X$'s dependence while the bounded number of levels keeps the attribute model's cells populated.

\noindent A joint categorical over the observed combinations (Proposition~\ref{prop:fd}) discloses the same combinations but pools those beyond 200 levels: with 565 states, 610 full state names and 8 countries the Contoso customers kept (Country, State)$\to$StateFull on 19\% of the synthetic rows (K13). Integer codes (SubCategoryKey) were not candidates of the joint categorical at all (subcategory$\to$category 69\%).

\begin{proposition}[Labels nested in a code]\label{prop:nestedlabels}
Let a categorical column $c$ have more levels than the marginal model keeps (200) and let a coarser column $g$ of the row (a state, or a dimension code) hold most of each label's rows: on the sampled rows whose label occurs at least twice, the share $s$ of rows in their label's most common $g$ is at least 0.8 and passes the test of Proposition~\ref{prop:fdrows}. Generate $c$ after $g$ from
\[p(c\mid g)=\frac{n_{gc}+\beta\,p_G(c)}{n_g+\beta},\qquad p_G(c)=\frac{n_c+1}{n+V},\]
over the $V$ labels held by at least $m$ source rows (\texttt{text.min\_token\_rows}, default 3). (i) This is the posterior predictive of a Dirichlet$(\beta p_G)$ prior on each group's label distribution, so as $n_g\to\infty$ it converges to the source's $P(c\mid g)$. (ii) No label held by fewer than $m$ rows is ever generated. (iii) A label is generated outside its source groups with probability at most $\beta/(n_g+\beta)$ per row.
\end{proposition}
\begin{proof}
(i) Dirichlet--multinomial conjugacy; the prior weight $\beta/(n_g+\beta)\to0$. (ii) Only vocabulary labels have mass. (iii) Outside $g$'s labels only the prior term $\beta p_G/(n_g+\beta)$ has mass.
\end{proof}
\noindent Contoso's 34{,}581 cities were a categorical pooled to 200 levels and drawn without their state: City$\to$State held on 0.92 of the source rows and 0.10 of the synthetic ones, and the Customer table's worst association (City with StateFull, 0.77 against 0.20) was the run's last scorecard failure (K17). The text model with \texttt{whole\_values} (each label one token, order 0) is this predictive; its end-of-text token takes part of the base measure, so the effective prior weight is $\beta$ times the labels' share of $p_G$, below $\beta$, and (i)--(iii) hold as stated. The labels lose their copula dependence on the row's other columns except through $g$, which is where the source keeps it. Check C53 tests the draft and the draws.

\subsection{Columns that are functions of other columns}
\begin{proposition}[Ratio columns]\label{prop:ratio}
Let $y=\mathrm{round}_d(x f)$ where $f$ takes values in a finite set $F$ of numbers with at most $k_0$ decimals and $d$ is the stored precision of $y$. (a) On every row with $|x|\ge 10^{k_0-d}$, $\mathrm{round}_{k_0}(y/x)=f$. Hence the coarsest $k\in\{2,\dots,8\}$ whose rounded ratios reproduce $y$ to the stored precision on at least 99.9\% of rows recovers $F$ when the condition holds on those rows. (b) Modelling $f$ as a categorical attribute and deriving $y=\mathrm{round}_d(xf)$ makes $y\in\mathrm{round}_d(xF)$ on every synthetic row, keeps the joint law of $f$ with the other attributes that the attribute model gives any categorical, and keeps two ratio columns that share their factor ($y_1=\mathrm{round}_d(x_1f)$, $y_2=\mathrm{round}_d(x_2f)$) in the same proportion on every row. (c) Factors seen in fewer than 5 rows are replaced by the nearest common factor; $y$ then changes by at most $|x|\,\Delta f$ on those rows, whose number is reported.
\end{proposition}
\begin{proof}
(a) $|y-xf|\le\frac12 10^{-d}$, so $|y/x-f|\le\frac12 10^{-d}/|x|\le\frac12 10^{-k_0}$ and rounding to $k_0$ decimals returns $f$ (on a tie the stored value also matches the other neighbour). Any $k<k_0$ rounds some $f\in F$ away from itself and fails to reproduce the rows that carry it, unless those are fewer than 0.1\%. (b) The derive rule is applied after the factor is generated; a categorical attribute is modelled like any other. (c) $|x f-x f'|=|x||f-f'|$.
\end{proof}
\noindent Discovery tests $x$ among the row's numeric columns and its parents' (one hop), with at most 50 factors; the chance level is the share of rows reproduced by some factor of $F$ when $x$ is permuted. A column is never its own base's base. Contoso: unit price $=$ list price $\times$ one of 16 factors, unit cost $=$ cost $\times$ the same factor, net price $=$ unit price $\times$ one of 15 discount factors, all to 4 decimals on 100\% of rows; v0.5 drew the three prices from a copula (K14, 1\% and 0.07\% of rows right).

\begin{proposition}[Column offsets]\label{prop:coloffset}
For two date (or number) columns $a,b$ of a row the map $(a,b)\mapsto(a,b-a)$ is a bijection; modelling the gap $b-a$ as an attribute and deriving $b$ keeps the gap distribution, its support (never negative) and its dependence on the other attributes as the attribute model gives them, and $b=a$ exactly on every row whose gap is 0. Discovery requires a gap that is never negative on 99.9\% of rows with at most 60 distinct values, beyond the chance share of permuted pairs whose gap falls in that set, and $a$ set on 99.9\% of the rows where $b$ is set ($b=a+\text{gap}$ is null wherever $a$ is; AdventureWorks' product StartDate is always set while EndDate mostly is not).
\end{proposition}
\begin{proof}
As Proposition~\ref{prop:offset} with the identity as the order column.
\end{proof}

\begin{proposition}[Views, reference tables, keyed and mapped lookups]\label{prop:view}
(a) Let $V$ and $B$ have the same key $U$ with the same key values, and let every column of $V$ be a column of $B$ (equal on the joined rows) or a lookup through $V$'s links. Then $V=B\bowtie(\text{its parents})$; rebuilding $V$ from the synthetic $B$ and parents by the same lookups gives $V_{\mathrm{syn}}=B_{\mathrm{syn}}\bowtie(\text{parents}_{\mathrm{syn}})$ exactly, with $|V_{\mathrm{syn}}|=|B_{\mathrm{syn}}|$. (b) A keyed lookup $y=R[v](e_1,\dots,e_m)$ into a copied reference table $R$ keyed by $(r_1,\dots,r_m)$ holds on every synthetic row whose key $(e_1,\dots,e_m)$ is a key of $R$; if $R$ is complete over the product of the key domains and each $e_j$ stays in the domain of $r_j$, it holds on every row. (c) A mapped lookup $y=m(p.x)$ holds on every row whose $p.x$ is in the map.
\end{proposition}
\begin{proof}
(a) Every column of $V$ is determined by $B$'s row and the parents; the rebuilt rows are those values. (b), (c) The lookups are evaluated on the synthetic rows after the columns they read.
\end{proof}
\noindent Contoso's Sales table is OrderRows joined to Orders with ExchangeRate $=$ rate(OrderDate, `USD', CurrencyCode); CurrencyExchange is the complete grid of 4{,}018 dates $\times$ 5 $\times$ 5 currencies. v0.5 modelled Sales as a second fact (K04: 0\% of its rows matched the order lines) and dropped CurrencyExchange (K01). Introspect tests a view on complete groups of keys read from both tables, and a keyed lookup over the combinations of the reference key's sources: row columns whose values lie in the key column's domain, and constants for key columns with at most 20 values.

\noindent\textbf{Files without declared keys.} The Parquet release declares no keys and names two columns differently from the SQL scripts (OrderDate is DT in orders, LineNumber is RowNumber in orderrows). A fact's pair (link $k$, integer column $c$) is taken as its line key when $c$ is distinct within $k$ and numbers the rows from 0 or 1 without gaps on complete groups of $k$ (at least 99.9\% of 30 or more groups with two or more rows; chance does not number thousands of rows this way). Proposition~\ref{prop:view}(a) needs only that $V$'s columns equal columns of $B$ on the joined rows, not that they share names: the key columns are paired by name with one left on each side paired with the other, and a column of $V$ equal on 99.9\% of the joined rows to exactly one unused column of $B$ is read from it (\texttt{base\_column}).

\noindent\textbf{Precedence among drafted structures.} A ratio $y=x\,f$ is not drafted when $x$ is computed from $y$ through derive rules or ratios already chosen (AdventureWorks: DealerPrice $=0.6\,$ListPrice with ListPrice $=$ StandardCost $\times f$ and StandardCost $=$ DealerPrice $\times f'$ was a cycle), and a column given by a derive rule is not a column of an implied dimension (ReorderPoint $=0.75\,$SafetyStockLevel keeps its formula, which holds for any level). The plan rejects any spec whose derived columns, ratio and offset columns included, depend on each other in a cycle.

\subsection{Rules across tables}
\begin{proposition}[Conditioning links on rules: target and bias]\label{prop:crossrules}
Let an entity's links $L$ be drawn given segments and attributes $A$ from $Q(L\mid z)$ and let $\mathcal R$ be rules over $(L,A)$ and the parents' columns. The independence Metropolis--Hastings kernel of Section~\ref{sec:v05} targets $Q(L\mid z,A,\mathcal R)$ and leaves the law of $A$ unchanged. (a) If $\mathcal R$ makes a class of a link a function of $A$ (the online store iff the delivery gap is positive), the synthetic share of that class equals the synthetic share of the attribute event, so it has the attribute model's accuracy and no conditioning bias. (b) In general, if the fitted joint law $Q$ of $(A,L)$ is within total variation $\delta$ of the source law $P$ and $P(\mathcal R)=1$, then $\mathrm{TV}(Q(\cdot\mid\mathcal R),P)\le2\delta$: conditioning on rules the source satisfies at most doubles the model's distance from the source.
\end{proposition}
\begin{proof}
The kernel moves links only; its acceptance ratio is 1 on valid states, so it is a Metropolis--Hastings kernel for the restricted law, and $A$ is never resampled. (a) On valid states $\mathbf 1[L\in C]=g(A)$. (b) $Q(\mathcal R)\ge P(\mathcal R)-\delta=1-\delta$; $\mathrm{TV}(Q(\cdot\mid\mathcal R),Q)=1-Q(\mathcal R)\le\delta$; the triangle inequality gives $2\delta$.
\end{proof}
\noindent Introspect drafts two kinds of cross-table rules, each tested against its permutation chance (Proposition~\ref{prop:signif}): a date of the row inside a parent's date window (\texttt{days(StoreKey.OpenDate) <= days(OrderDate)}; a null bound is open), and an equality between columns of two parents, or of the row, with one exception member of a copied parent (\texttt{CustomerKey.Country == StoreKey.CountryCode or StoreKey == 999999}), with its converse when the exception never satisfies the equality. Functions \texttt{days} and \texttt{add\_days} compare and shift dates whatever their storage (ISO text, datetimes, day numbers).

\noindent\textbf{Rules within a row.} Introspect also drafts order rules between two numeric or date columns of the same row ($\texttt{Cost < Price}$, $\texttt{days(StartDT) <= days(EndDT)}$), with the same test against the share of permuted pairs that satisfy the order. Ratio, offset and derived columns are skipped, since their order already follows from their construction. These rules read no link and are enforced by redrawing a row's attributes from its own conditional until they hold, so Proposition~\ref{prop:crossrules}(b) bounds their cost; the first v0.6 run (before these rules) left 1.5\% of Contoso's products priced below cost.

\begin{proposition}[Calibrated exact draws of links under rules]\label{prop:rake}
Let an entity have no unique key over its links, so that given segments its rows' links are independent with laws $Q_i$, and let $\mathcal R$ be its key rules. Let $c(L)$ be the row's cell (the joint of the categorical columns, at most 50 levels each, that the rules read through the links: the customer's country $\times$ the store's country) and $t_j(L)$ the segment of link $j$'s parent. For weights $w>0$ on cells and $u_j>0$ on parent segments let
\[\pi_i(L)\propto Q_i(L)\,\mathbf 1[\mathcal R]\,w(c(L))\prod_j u_j(t_j(L)).\]
(i) $\pi_i(\cdot\mid c,t)=Q_i(\cdot\mid c,t,\mathcal R)$: within a cell and parent segments the law is the model restricted to the rules. (ii) A row drawn as follows has law $\pi_i$: a cell $c$ with probability $\propto w(c)\prod_j\sum_t\phi_{j,z_it}u_j(t)W_j(t,c_j)$; for each link a segment $t\propto\phi_{j,z_it}u_j(t)W_j(t,c_j)$ and a parent of $t$ with value $c_j$ by the link's member law; the row is kept if $\mathcal R$ holds and drawn again otherwise. (iii) The weights solving $\pi(c)=P(c)$ and $\pi(t_j)=\bar Q(t_j)$ (the cell margin of the source, and the parent-segment margins the fitted link model gives before conditioning) give the I-projection of the restricted law onto those margins (Csisz\'ar 1975); iterative proportional fitting of $w$ and $u$ converges to it when the margins are attainable.
\end{proposition}
\begin{proof}
(i) The tilt is constant on each $(c,t)$. (ii) The draw before the check has probability $\propto w(c)\prod_j\phi_{j,z_it_j}u_j(t_j)\,\omega_{j,t_j}(L_j)=w(c)\prod_ju_j(t_j)\,Q_i(L)$ on the links with cell $c$, where $\omega_{j,t}$ is the member law of segment $t$ ($W_j(t,c_j)$ cancels between the two stages); rejection on $\mathcal R$ restricts it. (iii) The I-projection onto linear (margin) constraints is an exponential tilt by functions of the constrained variables, i.e.\ of this form (Csisz\'ar 1975, Theorem 3.1); IPF is its cyclic coordinate ascent.
\end{proof}
\noindent Without calibration Contoso's orders came 66\% from US customers against 51\% in the source, because the store-country rule rejected more proposals for the countries with fewer stores (K10); with the calibrated draws the cell margin is within 0.01 in total variation. The engine calibrates $w$ only and keeps $u_j=1$: adding the parent-segment margins $\bar Q(t_j)$ as a second set of constraints helped on S11 but oscillated on Contoso, where the rules make them unattainable together with the cells (Proposition~\ref{prop:rake}(iii) needs attainable margins); their distance before and after is logged. What remains is a shift in who buys: conditioning on a rule that the degree model already reflects (an order date inside the customer's active window; the source's counts per customer already grow with the window) concentrates the links on the customers that satisfy it most often. On S11 the share of customers with an order is 0.53--0.55 against 0.57 (0.58--0.60 without the rules); on Contoso 0.48 against 0.50, with 1.86 orders per buying customer against 1.79. Separating a customer's exposure (the orders it is eligible for) from its propensity in the degree model would remove it; it is listed as open. A row whose own segment mixture gives the rules no chance (one Contoso order whose customer segment never used the online store) is drawn after 100 attempts from the link's overall segment mixture $\sum_k\hat\pi_k\phi_{j,k\cdot}$, which stays inside the fitted model's support; rows still invalid after 100 more go to the Metropolis--Hastings kernel of Proposition~\ref{prop:crossrules}, which also serves entities with unique keys over their links. A cell coordinate that a stratified link fixes cannot be calibrated: its margin is that of the strata. A rule that reads only stratified links and the row's attributes (the order date before the store's closing date, with the store fixed by the stratum) cannot be repaired by moving links at all; it is a rule on the attributes given the links and is enforced by redrawing the attributes, together with the link rules that read the same attributes (Proposition~\ref{prop:crossrules}). On S11 five orders were left violating such rules after the repair rounds once the store became a stratifier.

\begin{proposition}[Keys in date order]\label{prop:keyorder}
Synthetic surrogate keys are labels the model never reads. Assigning the $j$-th smallest key to the row with the $j$-th smallest value of a column $c$ (ties in row order) is a permutation of labels: the joint law of all other values is unchanged and keys sorted by number are sorted by $c$ on every row, as in the source (Contoso order numbers, K05).
\end{proposition}
\begin{proof}
Children reference rows, not key values. The numbering is applied when the rows are drawn and again after the attribute redraws of the rules (which can move the ordering column), with the same permutation applied to every link that already points at the table.
\end{proof}

\begin{proposition}[Conditional inclusion on a fact's anchor]\label{prop:anchorincl}
Let a fact's anchor rows in segment $k$ have i.i.d.\ row counts with law $P_k$ ($P_k(0)=P0_k$; for an occupancy fact the Bernoulli activity with $P0_k=1-\pi_k$), and let the source satisfy a conditional inclusion dependency (Bravo, Fan and Ma 2007): anchor rows at an excluded level of a column (employees with \texttt{SalesPersonFlag = 0}) have no fact rows, a share $e_k$ of segment $k$. Then the eligible rows' counts follow
\[q_k(0)=\frac{P0_k-e_k}{1-e_k},\qquad q_k(c)=\frac{P_k(c)}{1-e_k}\ (c>0),\]
which keeps the segment's law of counts over all its rows. A draw from $P_k$ that is zero and kept with probability $r_k=1-e_k/P0_k$, and drawn again otherwise, has law $q_k$. For the Bernoulli law this is activity $\min(1,\pi_k/(1-e_k))$. When $e_k>P0_k$ no such law exists, $r_k=0$, and the shortfall is logged.
\end{proposition}
\begin{proof}
$P_k=e_k\delta_0+(1-e_k)q_k$ gives $q_k$. Under the rejection, $\Pr(\text{output}=0)\propto P0_kr_k$ and $\Pr(\text{output}=c)\propto P_k(c)$ for $c>0$; with $r_k=1-e_k/P0_k$ the normaliser is $1-e_k$ and the ratios are those of $q_k$. The DP count law is i.i.d.\ given its Bayesian-bootstrap weights, so these are drawn once and $P0_k$ is computed from them.
\end{proof}
\noindent The engine fits these rules on every link with at least 10 expected references to the excluded level and none observed, and masks the excluded rows when the link chooses a parent, but it skipped a fact's anchor, whose rows are not chosen but given counts. On AdventureWorks the employees' segments usually separate the 17 sales staff, so quotas went to sales staff by chance; after one segment merge 42 of 183 synthetic quotas named non-sales employees. The rule is now also fitted for the anchor. A Poisson anchor without the DP law has its row weights masked and rescaled within the segment instead. Check C52 forces a single employee segment: without the rule 165 quotas and 75 bonuses reference non-sales staff, and with it none do.

\noindent\textbf{The partner of an occupancy fact.} The same rule on the partner side (order lines never name a product without a price) was applied by masking the partner's row weights, which the Poisson block model uses but the occupancy model does not: it draws each (anchor, partner) pair of block $(k,l)$ with probability $p^*_{kl}$ over all partner rows of segment $l$. One AdventureWorks reseller line in 60{,}855 named a product without a price, and so had no unit price. The excluded partner rows had no pairs in the source, so the block's pairs belong to its eligible rows: they are drawn over the $E^*_l$ eligible synthetic rows only, with $p^*_{kl}=\min(1,p_{kl}m_l/E^*_l)$. This is the scale rule with $E^*_l$ in place of $m^*_l$, and keeps the expected pairs per anchor, $p_{kl}m_l$. Check C56 forces one product segment: before the change 2{,}251 of about 7{,}200 lines named never-sold products, and after it none do, with lines per order within 1\%.

\noindent\textbf{Mandatory participation under the DP count model.} When the counts per parent are modelled by the Dirichlet process (Section~\ref{sec:v05}), new atoms come from the Poisson base measure, which puts mass $e^{-\lambda}$ on 0 even when participation is mandatory (every order has a line); one synthetic Contoso order had no lines. With mandatory participation the base measure is $1+\mathrm{Poisson}(\lambda-1)$, the same shift the parametric model uses: the support is that of the source and the mean is unchanged. Both the unconditional draws and the draws given a segment's total use it (check C54).

\noindent\textbf{Self-anchored dimensions.} A dimension of fewer than 200 levels with no categorical label of at most 50 levels to anchor it (AdventureWorks resellers: the number of employees determines annual sales) was drafted as a link. The code's dependence on the other columns then rested on the segments alone, and business type with number of employees fell from 0.91 to 0.06. When the marginal model keeps the code exactly (a category, or an integer with at most 30 values) the code now anchors itself: it stays a modelled attribute, and the columns it determines are keyed lookups into the copied dimension by its value. The dependency then holds on every row, and the copula keeps the dependence. Otherwise no dimension is drafted. On AdventureWorks the prospective buyers' state code (38 values, which determines one column) is self-anchored, and the resellers' 76 employee counts are left to the copula (business type with number of employees 0.84 against 0.91). Check C55 tests both properties.

\subsection{Cost of stratification}
A Gibbs sweep scores every row of an entity against every segment: $O(nK)$ per attribute and link. Strata (rows of copied parents) nest segments, $K=\sum_s\mathrm{cap}_s$, so a 555-level geography with 105{,}000 customers gave $K=1{,}757$ and 21\,s per sweep (0.6\,s before). Stratifiers are now dropped, largest first, while $nK$ exceeds \texttt{segments.stratify\_max\_cells} ($5\cdot10^7$); a dropped stratifier's link is modelled by the degree-corrected block model like any other link (Proposition~\ref{prop:shares}), and the decision is logged. AdventureWorks ($18{,}484\times$ about 700) stays stratified. When stratifiers must be dropped, implied dimensions go first, whatever their size: their anchor keeps them as modelled attributes, while a real copied parent's link has no other carrier. In the first v0.6 AdventureWorks run the customers' education and occupation dimensions (5 levels each) were kept and the geography dropped, and the share of synthetic customer combinations not seen in the source rose to 0.12; excluding implied dimensions altogether fixed that but took the subcategory strata from Contoso's products (Weight with SubCategoryName 0.85 in the source, 0.52 synthetic).

%=====================================================================
\section{Bayesian synthesis and the whole data model (v0.7)}\label{sec:v07}
Version 0.7 makes every synthetic dataset a draw from the posterior predictive distribution, closes the open items of the Contoso and AdventureWorks reports (T5, T11, T14, the participation share and the quota segment), and adds the warehouse structures that the four databases did not all contain but that are standard practice (Kimball and Ross 2013): type 2 slowly changing dimensions, unknown members and bridge tables. Each structure is detected by \texttt{introspect}, written to the specification and can also be declared by the user (\texttt{docs/spec-reference.md}). Checks C57--C68 test this section on fake data with known truth.

\subsection{Posterior predictive synthesis}
\begin{proposition}[One parameter draw per dataset]\label{prop:postpred}
Let $\theta$ collect every parameter block of the model (segment rates, Dirichlet weights, copula correlation, null shares, the shares below) with posterior $p(\theta\mid D)$. Dataset $i$ draws $\theta^{(i)}\sim p(\theta\mid D)$ once from its own parameter stream and then every row from $p(\cdot\mid\theta^{(i)})$ with a separate data stream. (i) $D^{*(i)}$ is an exact draw from the posterior predictive $\int p(D^*\mid\theta)\,p(\theta\mid D)\,d\theta$ whenever each block's draw is exact (conjugate blocks) and approximately so for the MCMC blocks. (ii) For a statistic $q$, $\mathrm{Var}(q(D^*))=\E[\mathrm{Var}(q\mid\theta)]+\mathrm{Var}(\E[q\mid\theta])$; drawing $\theta$ afresh per row instead removes the second term asymptotically, so the between-dataset variance $b_m$ of the combining rules for fully synthetic data, $T_m=(1+1/m)b_m-\bar u_m$ (Raghunathan, Reiter and Rubin 2003), would be too small and intervals would under-cover. (iii) Datasets $i\ne j$ are independent given $D$.
\end{proposition}
\begin{proof}
(i) Composition sampling. (ii) Law of total variance; with per-row draws the rows are i.i.d.\ from the predictive $\int p(x\mid\theta)p(\theta\mid D)d\theta$, so $\mathrm{Var}(\E[q\mid\cdot])$ of a mean of $n$ rows is $O(1/n)$ of the per-row variance instead of the posterior variance of $\E[q\mid\theta]$. (iii) The streams are spawned from independent children of one \texttt{SeedSequence}.
\end{proof}
\noindent \texttt{posterior.mode: mean} replaces each draw by its posterior mean (each spike-and-slab indicator by its posterior median): the plug-in synthesis of Reiter and Kinney (2012), stable from run to run and meant for test data. Kept as point values on purpose: deterministic identities (lookups, derives, aggregates, ratios, offsets), copied tables, storage precision, privacy thresholds, user-set sizes, and model-structure choices that a Bayes factor settles (segment counts after merging, structural zeros above the prior threshold, count laws, the frailty switch). Averaging over structure would change the table layouts between datasets. Two blocks keep an empirical Bayes estimate: the Gamma frailty shape $\alpha$ (a moment estimate from the dispersion index) and the copula sampler's hyper-parameters; Proposition~\ref{prop:hurdle} corrects the one consequence of $\hat\alpha$ that the checks exposed (the share of parents without children). \emph{Check C57}: a dataset is reproducible from its index, datasets differ by the posterior spread, and \texttt{mean} mode is stable.

\subsection{Rules that hold on most rows}
\begin{proposition}[Soft rules]\label{prop:soft}
For a rule $y=f(x)$ that holds on $h$ of $n$ source rows, draw $\lambda\sim\mathrm{Beta}(1+h,1+n-h)$ and offset weights $w\sim\Dir(c+\frac12)$ over the observed offsets $d_j=y-f(x)$ (counts $c_j$) once per dataset; set $y=f(x)$ with probability $\lambda$, otherwise $y=f(x)+d_J$ with $J\sim w$ (or the modelled value of $y$ when there are no offsets or more than 50). (i) The expected share of synthetic rows on which the rule holds is $\E\lambda=(1+h)/(2+n)$, and its posterior predictive spread is that of the Beta-binomial. (ii) Every synthetic violation is an observed offset, so a rule that fails by one year (an age before the birthday) fails by one year. (iii) As $n\to\infty$ the share converges to $h/n$.
\end{proposition}
\begin{proof}
Beta--binomial and Dirichlet--multinomial conjugacy; (ii) by construction.
\end{proof}
\noindent \texttt{introspect} finds such rules for an age or a number of whole years between two dates: the reference date is swept over the interval that each row's value implies and kept when the rule holds on at least 80\% of rows, far above the share that permuting $y$ gives. Contoso's customer age holds on 90.5\% of source rows; v0.6 gave 42\%, v0.7 gives 90.4\% (T14). \emph{Check C58}.

\subsection{Values shared by a group}
\begin{proposition}[Group effects]\label{prop:groupfx}
Let rows be grouped by a key $g$ (columns of the row or of its parents, such as the product category and the order date), and let each row draw $U_r\sim\mathrm{Unif}(0,1)$ once. The group's value is $u_g=y_{r^*}$ with $r^*=\arg\min_{r\in g}U_r$, the value that the row model drew for that row, and every row of $g$ takes $u_g$. (i) The rule ``one value per group'' holds on every synthetic group. (ii) $r^*$ is uniform on the rows of $g$, so $u_g$ has the law $\frac1{|g|}\sum_{r\in g}p(y\mid x_r,\theta)$: a shared random effect whose distribution is the row model's averaged over the group (Hu, Reiter and Wang 2018 use the same construction for households). (iii) Redrawing some rows of $g$ (key rules) does not change $r^*$ unless $r^*$ itself is redrawn, so the recomputation after redraws is stable. (iv) For a ratio column $y=\mathrm{round}(xf)$ the effect is placed on the factor $f$, so rows of a group share the factor and keep their own base $x$.
\end{proposition}
\begin{proof}
(i) By assignment. (ii) The minimum of i.i.d.\ continuous uniforms is attained at a uniformly random index. (iii) The priorities are stored with the rows. (iv) Proposition~\ref{prop:ratio}(b).
\end{proof}
\noindent \emph{Detection.} With $s$ the share of groups of two or more rows that hold one value (the $g_3$-type measure; Parciak et al.\ 2024) and $p_0$ its permutation mean, a grouping is proposed when $s\ge0.999$ and the test of Proposition~\ref{prop:fdrows} passes; among groupings that pass, the coarsest is kept (fewest groups). Contoso's price factors are one per category and day on 100\% of source groups; v0.6 drew them per row (30\% of groups with one factor), v0.7 gives 100\% (T11). \emph{Check C59}.

\subsection{Two links that agree more often than chance}
\begin{proposition}[Link affinity]\label{prop:affinity}
Let a row have links $a$ and $b$ whose parents carry an attribute $v$ (the customer's area and the store's area). Under the link model $Q$ the share of rows with $v(a)=v(b)$ is $\pi_Q$; in the source it is $h/n$. Model
\[P(b\mid\text{row})=(1-\rho)\,Q(b\mid\text{row})+\rho\,Q\bigl(b\mid\text{row},\,v(b)=v(a)\bigr),\]
the diagonal-inflated quasi-independence table of Goodman (1968), with $\rho$ set so that the share of matching rows equals $s\sim\mathrm{Beta}(1+h,1+n-h)$, drawn per dataset. The engine realises it by moving $\lceil sn\rceil-m$ non-matching rows ($m$ already match) to a parent with value $v(a)$, chosen with probability proportional to the parent's degree $+\frac12$, keeping a move only when every key rule of the row still holds. (i) Whenever enough non-matching rows have a matching parent, the synthetic matching share is $\mathrm{round}(sn)/n$ exactly, with posterior predictive spread $\mathrm{Beta}(1+h,1+n-h)$. (ii) Moved rows keep their other columns; within the matching parents the choice follows the link's popularity (a Dirichlet$(\deg+\frac12)$ posterior mean). (iii) Rows are only moved toward matching parents, so the share of rows matching on a coarser attribute (country when $v$ is the area) does not decrease.
\end{proposition}
\begin{proof}
(i) Each accepted move turns one non-matching row into a matching one and no move does the reverse. (ii) By construction. (iii) $v(a)=v(b)$ implies equality of every function of $v$.
\end{proof}
\noindent \emph{Detection.} For each pair of links and each pair of attributes of their parents, the lift $s/\pi_0$ over the analytic permutation mean $\pi_0=\sum_v p_a(v)p_b(v)$ (the expected equality share when one column is permuted) must be at least 2 with the binomial test of Proposition~\ref{prop:signif}, for attributes with the same name or sharing at least half their values; the finest such attribute (most values) is kept, and the link to the smaller parent is the one moved. When the equality holds on at least 99.9\% of rows it is a rule $a.v=b.v$ (a \texttt{check}), not an affinity. Contoso: customer and physical store in the same geography area on 12.6\% of source lines; v0.6 gave 1.6\%, v0.7 gives 12.1\% (T5). \emph{Check C60}.

\subsection{Parents without children and segments that mix eligible rows}
\begin{proposition}[Participation hurdle]\label{prop:hurdle}
For a link from an entity to a parent with segments $t$, let $Z_t$ of the $N_t$ source parents of segment $t$ have no child; draw $s_{0t}\sim\mathrm{Beta}(1+Z_t,1+N_t-Z_t)$ per dataset (Mullahy 1986). After the links are drawn, in each segment with more empty eligible parents than $\mathrm{round}(s_{0t}N'_t)$, children are moved from parents of the same segment with two or more children (chosen with probability $\propto$ degree $-1$) to empty eligible parents, each move kept only when the child's key rules hold. Moves stay within a segment because the child's other columns were drawn given its parent's segment (what a customer of that segment buys); codes implied inside a table (Proposition~\ref{prop:anchor}) are drawn with their anchor and are not moved. (i) The number of children is unchanged, so the mean degree is unchanged. (ii) No move creates an empty parent. (iii) When enough movable children exist, the share of empty parents in segment $t$ is $\mathrm{round}(s_{0t}N'_t)/N'_t$.
\end{proposition}
\begin{proof}
(i) Moves change only the parent. (ii) The source parent keeps at least one child. (iii) Each move fills one empty parent.
\end{proof}
\noindent A Gamma frailty with shape $\alpha$ and rate $\lambda$ gives a zero share $(1+\lambda/\alpha)^{-\alpha}$; with the moment estimate of $\alpha$ this missed Contoso's share of customers with an order (0.48 against 0.50). The hurdle fixes the zero share without changing the number of children. \emph{Check C61}.

\begin{proposition}[Cannot-link strata]\label{prop:cannotlink}
Let a categorical column $c$ of an entity have a level $\ell$ whose $n_\ell$ rows are never referenced by a child relation that references $D$ of the $n$ rows, with $D n_\ell/n\ge10$. Under independence of reference and level, the probability of no reference is at most $\exp(-Dn_\ell/n)\le e^{-10}$, so $\ell$ is a structural zero. Stratifying the entity's segments by the indicator $\mathbf 1[c=\ell]$ (a cannot-link constraint, Shental et al.\ 2003) gives segments that never mix the two kinds of rows, and the child's link model gives the excluded segments rate zero, so no synthetic child references an excluded row.
\end{proposition}
\begin{proof}
Sampling $D$ referenced rows without replacement: $P(\text{none in }\ell)=\binom{n-n_\ell}{D}/\binom nD\le(1-n_\ell/n)^D\le e^{-Dn_\ell/n}$. Within strata the segment rates are fitted separately, and a stratum with no source children has posterior rate mass only from the prior, which the structural-zero test of Section~\ref{sec:v03} sets to zero.
\end{proof}
\noindent AdventureWorks: sales quotas exist only for sales staff; the employee segment that mixed sales staff with others made the quota count move between runs (138--173 against 163). Small tables (fewer than two segments' worth of rows) keep the earlier attribute strata and get exact stratum counts. \emph{Check C62}.

\subsection{Warehouse structures: slowly changing dimensions, unknown members, bridges}
\begin{proposition}[SCD type 2]\label{prop:scd2}
A type 2 dimension stores versions of a natural key $k$ with validity intervals $[f_v,t_v]$, the latest open ($t=\text{null}$) or closed. Fit: versions per key $V\sim\Dir(c+\frac12)$ over $1..M$; the share of keys whose latest version is open $\sim\mathrm{Beta}(1+o,1+n_k-o)$ (used when there is no valid-to column); the gap $t_v-f_{v+1}$ in days $\sim\Dir(c+\frac12)$ over the observed gaps; stable columns (equal across a key's versions on $\ge99\%$ of keys with two or more versions). Synthesis groups the synthetic rows into keys with sizes drawn from $V$ (rows of one segment together). A row whose modelled $t$ is null (the row model draws the null with the row's segment, under the null-share posterior) becomes the latest version of its key: open rows are spread one per key by swapping with closed rows of keys that have none (same segment first; keys never span two strata, so a stratifying column stays constant per key), and within a key the sorted start dates go to the closed rows first and to the open row last. The key and the stable columns are copied from each key's first row and $t_v=f_{v+1}+\text{gap}$ for every earlier version. (i) Every key has at most one open version and its versions do not overlap when the gaps are non-positive (0 or $-1$ day, as in every clean history). (ii) The number of keys is $n/\E V$ in expectation and the version-count law is the posterior predictive. (iii) Stable columns are constant within each key. (iv) The current flag, when declared, equals its value exactly on the open versions. (v) When the open rows can be spread (at most one per key), every row keeps the open or closed status the row model gave it, so children linked to a row by the segment model keep the source's share of links to open versions.
\end{proposition}
\begin{proof}
(i) Only the latest version may be left open, and $t_v=f_{v+1}+\text{gap}\le f_{v+1}$ with $f_{v}<f_{v+1}$. (ii) Renewal argument: groups are drawn until $n$ rows are covered. (iii), (iv) By assignment after grouping. (v) Swaps move rows between keys, not statuses; only the extra open rows of a key that could not be swapped are closed.
\end{proof}
\noindent \emph{Detection} needs a valid-from and a valid-to date, a candidate natural key repeated on some rows, distinct start dates within a key and at most one open version per key (the latest), on at least 99.9\% of keys, with at least 5 keys holding two or more versions. AdventureWorks: \texttt{DimProduct} (504 product codes, 77 with two or more versions) and \texttt{DimEmployee} (290 national IDs, 5 with two or more); \texttt{DimCustomer} is correctly not one. Facts keep pointing at surrogate keys, which is how a type 2 dimension is joined. \emph{Check C63}.

\begin{proposition}[Unknown members]\label{prop:unknown}
Let a dimension contain placeholder rows $U$ (key $-1$, 0 or a label such as ``Unknown'') that $r$ child rows reference while $m$ child rows have a null key. The placeholders are removed before fitting and copied to the output, and their references become nulls. After synthesis each null reference is set to a placeholder with probability $q\sim\mathrm{Beta}(1+r,1+m)$ (the member chosen by $\Dir(c+\frac12)$ over the placeholders). (i) Placeholders never enter the dimension's segments or the facts' rates, so the one ``Unknown'' row cannot become a segment of its own or pull the attribute model toward its nulls. (ii) The expected share of references to placeholders among null-or-placeholder references is $(1+r)/(2+r+m)$. (iii) A fact drops rows whose dimension key is null from its block model; for facts the placeholder rows are therefore restored by cloning $\mathrm{round}(q'n^*)$ synthetic rows chosen uniformly ($q'\sim\mathrm{Beta}(1+r,1+n-r)$) and pointing them at a placeholder. This assumes the source rows that reference a placeholder are missing at random with respect to the other columns; the user declares otherwise by modelling the placeholder as a member (not listing it in \texttt{unknown\_members}). \emph{Check C64}.
\end{proposition}
\begin{proof}
(ii) Beta--binomial. (iii) Under missing at random the rows with an unknown member have the same conditional law as the other rows, which a uniform clone reproduces in expectation.
\end{proof}

\begin{proposition}[Bridge weights]\label{prop:bridge}
In a bridge table (accounts of a customer, diagnoses of a stay) the weights $w_i$ of a group sum to a total $T$ (1 by default). The modelled weights are normalised, $w_i\leftarrow Tw_i/\sum_{j\in g}w_j$, rounded to $d$ decimals and the rounding remainder added to the group's largest weight. (i) Each group's weights sum to $T$ to the stored precision. (ii) Ratios of weights within a group are those of the model. (iii) If the modelled weights are independent $\mathrm{Gamma}(a_i,1)$ the normalised weights are $\Dir(a)$ (the standard construction), so a group's allocation has the Dirichlet law with the model's shapes.
\end{proposition}
\begin{proof}
(i) After rounding, the remainder is a multiple of $10^{-d}$ smaller than the largest weight. (ii) Scaling. (iii) The Gamma representation of the Dirichlet.
\end{proof}
\noindent \emph{Detection}: a fact whose float column sums to the same value (1 or 100) on at least 99.9\% of the groups (with two or more rows) of a foreign key. \emph{Check C65}.

\subsection{Identifying the whole model}
\begin{proposition}[Equality through two links]\label{prop:eqlink}
Let $a$ and $b$ be columns (possibly through links) of the $n$ rows of a table, with level shares $p_a(v)$ and $p_b(v)$. Under a uniform random permutation of $b$ the expected share of rows with $a=b$ is $\pi_0=\sum_v p_a(v)\,p_b(v)$. A rule $a=b$ is reported when the observed share is at least $0.999$, above $\pi_0$, and passes the test of Proposition~\ref{prop:fdrows}.
\end{proposition}
\begin{proof}
$\E\frac1n\sum_i\mathbf 1[a_i=b_{\pi(i)}]=\frac1n\sum_i P(b_{\pi(i)}=a_i)=\frac1n\sum_i n_b(a_i)/n=\sum_v p_a(v)p_b(v)$.
\end{proof}
\noindent The closed form replaces a Monte Carlo estimate, and constant groups are screened in $O(n\log n)$ by sorting, which made \texttt{introspect} on AdventureWorks (27 tables) feasible with the new detections. \texttt{relsynth model} writes the identified model with Kimball's vocabulary: facts (transaction, periodic snapshot, accumulating snapshot, factless), dimensions (conformed, role-playing, degenerate, outrigger, SCD type 2, bridge), hierarchies with their levels, rules, and the join graph as a Mermaid entity-relationship diagram.

\subsection{The specification}
Every key the specification accepts has a type, a default and a description in one schema table (\texttt{relsynth/schema.py}); unknown keys are errors with the closest valid key as a hint (\emph{check C67}), and the key reference in \texttt{docs/spec-reference.md} is generated from that table. Column laws can be stated instead of learned (\texttt{distribution}: fixed, categorical probabilities, uniform, normal, lognormal, Poisson, date range, sequence, empirical); a stated law replaces the learned one and is drawn independently of the other columns (once per row, so redraws of other columns keep it); \texttt{empirical} resamples the source values with Bayesian-bootstrap weights. Structures the data cannot reveal are listed as user guidance (Section 4 of the reference); \emph{check C66} runs its recipes. \emph{Check C68} runs the reference's annotated example, which uses every v0.7 option at once.

%=====================================================================
\section{Identification from the data, exact speed-ups and the exported model (v0.8)}\label{sec:v08}
Version 0.8 identifies keys and links from the data when the catalog declares none, makes the engine faster without changing a single output bit, corrects two run-length rules that the Contoso data showed to be wrong, and exports the fitted model as a self-contained mathematical object (\texttt{relsynth export-model}). Checks C69--C74 test this section.

\subsection{Keys and links from the data}
No schema template, naming list or known database is used: the evidence is the values, and optionally the column and table names weighted by how rare each word is in the database at hand (inverse token frequency, so words every table carries count for nothing; \texttt{--evidence data} switches names off). Primary keys are unique column combinations scored as in holistic key detection (Jiang and Naumann 2020). A link is a unary inclusion dependency $c\subseteq K$ with inclusion share $\sigma\ge0.95$; it counts as evidence only when it is unlikely by chance.

\begin{proposition}[Inclusion by chance]\label{prop:inclchance}
Let the $v\ge3$ distinct integer values of $c$ that lie in $K$ have extremes $a<b$, let $w=b-a+1$ and let $k=|K\cap[a,b]|$. Under the null that, given the extremes, the other $v-2$ distinct values are a uniform draw without replacement from the $w-2$ interior integers, the probability that all of them fall in $K$ is
\[P=\binom{k-2}{v-2}\Big/\binom{w-2}{v-2}\;\le\;d^{\,v-2},\qquad d=k/w.\]
A key that fills its range ($d=1$, a surrogate $1..N$) gives $P=1$: inclusion in it is no evidence at all, and the decision falls to the other features (randomness, coverage, names) or to a question.
\end{proposition}
\begin{proof}
The extremes are keys because every value counted lies in $K$. A uniform $(v-2)$-subset of the $w-2$ interior points lies in the $k-2$ interior keys with the hypergeometric probability shown. Each factor $(k-2-i)/(w-2-i)\le k/w$ for $k\le w$, which gives the bound.
\end{proof}

\begin{proposition}[Randomness bound]\label{prop:randomness}
Let $u_1,\dots,u_m\in[0,1]$ be the quantile positions of the child values in the sorted parent key and $F_m$ their empirical distribution. Then $W_1(F_m,U(0,1))=\int_0^1|F_m(x)-x|\,dx\le\frac12$, with equality only for a point mass at 0 or 1, so $R=1-2W_1\in[0,1]$. If the child values are a random sample of the key, $R\to1$ as $m\to\infty$.
\end{proposition}
\begin{proof}
For a point mass, $W_1(\delta_a,U)=\int_0^1|a-x|dx=\frac{a^2+(1-a)^2}2\le\frac12$, with equality at $a\in\{0,1\}$. Any $F$ on $[0,1]$ is a mixture $\int\delta_a\,d\mu(a)$, and $|F(x)-x|=|\int(\mathbf 1[a\le x]-x)\,d\mu(a)|\le\int|\mathbf 1[a\le x]-x|\,d\mu(a)$; integrating over $x$ gives $W_1(F,U)\le\int W_1(\delta_a,U)\,d\mu\le\frac12$. For a random sample, $W_1(F_m,U)\to0$ almost surely (Glivenko--Cantelli on a bounded interval). The implementation integrates $|F_m(x)-x|$ exactly piece by piece.
\end{proof}

\begin{proposition}[Chance of a forest]\label{prop:forest}
A table of $n$ rows with $r\ge1$ null pointers, whose other $m=n-r$ pointers are drawn independently and uniformly from the $n$ keys (a self-pointer counting as a top), has pointers that form a forest with probability
\[p(n,r)=\frac{r+1}{n}\Bigl(1+\frac1n\Bigr)^{m-1}\;\le\;\frac{e\,(r+1)}{n}.\]
In v0.8 a forest was accepted as evidence of a self-reference when $p(n,r)\le0.01$; otherwise the link needed other evidence (names, Proposition~\ref{prop:inclchance}) or was asked. \emph{v0.9:} $p(n,r)$ is the chance under \emph{uniform} pointers, which a code column (values $1..3$ in a table of 2{,}000 rows, its few targets pointing at themselves) does not follow; the decision now uses the permutation test of Proposition~\ref{prop:forestperm}, and $p(n,r)$ remains as reported context.
\end{proposition}
\begin{proof}
A rooted forest on $n$ labelled vertices with a given set of $k$ roots is a map from the other $n-k$ vertices to all $n$ vertices without cycles; there are $k\,n^{n-k-1}$ of them (Cayley's formula for forests). Let $s$ of the $m$ non-null pointers point to their own row; the other $m-s$ then form a forest rooted at the $r+s$ tops. Summing over $s$ and dividing by the $n^m$ maps,
\[p=\sum_{s=0}^m\binom ms\frac{(r+s)\,n^{m-s-1}}{n^m}=\frac1n\Bigl[r\Bigl(1+\tfrac1n\Bigr)^m+\frac mn\Bigl(1+\tfrac1n\Bigr)^{m-1}\Bigr]=\frac{r+1}n\Bigl(1+\tfrac1n\Bigr)^{m-1},\]
using $\sum_s\binom ms x^s=(1+x)^m$ and $\sum_s s\binom ms x^s=mx(1+x)^{m-1}$ with $x=1/n$, and $r(1+\frac1n)+\frac mn=r+1$. Finally $(1+\frac1n)^{m-1}\le e^{(m-1)/n}\le e$.
\end{proof}
\noindent The earlier documentation said that a random mapping ``almost surely has a cycle''; that holds only without null pointers. With $r$ tops the chance is about $e(r+1)/n$: 0.51 for Chinook's 8 employees with one top, 0.018 for AdventureWorks' 296 employees with one top. From values alone both are now asked instead of being accepted on weak evidence; with names, AdventureWorks' is found (its key words recur across the database) and Chinook's is still asked, and the declared catalog settles both when it has them. \emph{Check C71} verifies the three statistics against simulation and exact integrals; \emph{check C70} hides every declared key and link of the retail database and finds 8 of 8 links with names and 0 false links in both modes, with a question for each link that values alone cannot establish.

\noindent A column whose name shares a word with at least half of this database's primary-key names (learnt from the keys found, such as \emph{id} or \emph{key}) and whose values all lie in another table's key is asked about when its values do not establish the parent; with \texttt{--evidence data} names are not read and only the value evidence above applies.

\subsection{Exact speed-ups}
\begin{proposition}[Exact kernels]\label{prop:kernels}
Each of the following replaces an expression by code that performs the same IEEE-754 operations on each element in the same order and consumes the same random numbers, so outputs and the state of every random stream are bit-identical: (a) the Gaussian log-likelihood $L_{ik}\mathrel{+}=a_{kj}-\frac12(x_{ij}-\mu_{kj})^2/s^2_{kj}$ computed column by column (compiled when Numba is installed); (b) count tables by one \texttt{bincount} instead of \texttt{add.at} (integer counts are exact in floating point below $2^{53}$); (c) Dirichlet rows: NumPy draws $\Dir(\alpha)$ for $\max_j\alpha_j\ge0.1$ (corrected in v0.9: the v0.8 text said $\min$; below that NumPy switches to stick-breaking) as standard gammas $g_j$ in order, multiplied by $1/\sum_jg_j$, which is reproduced in one call; (d) categorical draws $\texttt{choice}(n,p)$ as \texttt{searchsorted} of the normalised cumulative sum at one uniform per draw; (e) the Gumbel-max draw $\arg\max_k(\ell_k-\log(-\log u_k))$ in one buffer; (f) the occupancy likelihood of a row computed only from the terms that row uses (a present pair or the absence term), the others never being read; (g) frequency counts by value instead of by row (Kish effective size, canonical spellings).
\end{proposition}
\begin{proof}
(a), (e) and (g) evaluate identical expressions on identical operands; (b) integer arithmetic; (c) and (d) follow NumPy's \texttt{Generator.dirichlet} and \texttt{Generator.choice} algorithms (uniforms drawn one per call equal those drawn as a vector, since \texttt{random(n)} consumes the stream in the same order); (f) the unread terms do not enter any output.
\end{proof}
\noindent \emph{Check C69} compares every kernel with the expression it replaces (and with a frozen copy of the previous occupancy code). On Contoso the fit is 2.5 times faster (Gibbs 3.7 times), sampling 2.1 times and evaluation 2.0 times, with every output file identical before the run-length changes below. Stan, Cython and C extensions were weighed: the remaining time is in vectorised NumPy calls already running compiled code, and the one loop that dominated (the Gaussian log-likelihood) is compiled by Numba when present; Stan is used for verification (Proposition~\ref{prop:stan}), not for production sampling, where its general-purpose NUTS sampler would be slower than the conjugate updates.

\subsection{Run-length control, corrected}
The Contoso fit with v0.7 settings left 3 of 56 copula chains and the final Gibbs phase with split $\hat R>1.1$. Two causes were found and fixed.

\begin{proposition}[Fixed-width stopping for correlation chains]\label{prop:fixedwidth}
Let a chain of correlation matrices have kept draws $C^{(1)},\dots,C^{(T)}$, $T\ge100$, and let $\hat\sigma_{ij}$ be the batch-means standard error (batch size $\lfloor\sqrt T\rfloor$) of the mean of $z_{ij}=\operatorname{atanh}C_{ij}$. Stop when every entry has split $\hat R\le1.1$ or $\hat\sigma_{ij}\le\varepsilon=0.01$, and $\log\det C$ has split $\hat R\le1.1$ or standard error $\le2p\varepsilon$. (i) For a polynomially ergodic chain the batch-means estimator is consistent, so $\bar z_{ij}$ is within $2\varepsilon$ of its posterior mean with probability about 0.95 (Jones, Haran, Caffo and Neath 2006; Flegal, Haran and Jones 2008). (ii) An error $\delta$ on the $z$ scale changes the conditional variance $1-\rho^2$ by a relative amount of at most $2|\rho|\delta\le2\delta$, uniformly in $\rho$, which is why the tolerance is set on the Fisher $z$ scale: near $\rho=\pm1$ a fixed tolerance on $\rho$ itself would be too loose.
\end{proposition}
\begin{proof}
(i) Markov chain central limit theorem with a consistent variance estimator. (ii) $d(1-\rho^2)/(1-\rho^2)=-2\rho\,d\rho/(1-\rho^2)=-2\rho\,dz$.
\end{proof}
\noindent Split $\hat R$ is scale-free: a nearly deterministic pair (Contoso's birthday and age, $\rho=-0.9989$, posterior s.d.\ $10^{-4}$) drifts by $10^{-4}$ and shows $\hat R=1.6$ although no synthetic value can reflect that drift; its standard error on the $z$ scale is 0.003. The fixed-width branch is what accepts such a chain; it needs at least 100 kept draws so that the batch means are not computed on a chain shorter than its autocorrelation.

\begin{proposition}[Rare columns per estimation sample]\label{prop:raresample}
Proposition~\ref{prop:rare} applied to the sample a copula is fitted on: a column with fewer than $m$ non-modal rows in a segment's sample ($m=$\texttt{marginals.min\_dependence\_rows}, default 5), which are also at most 1\% of that sample, is kept independent of the other columns in that segment. Within the segment the joint law of that column with any other changes by at most 2\% in total variation.
\end{proposition}
\begin{proof}
Proposition~\ref{prop:rare} with the segment's sample as the population and $q\le0.01$.
\end{proof}
\noindent A column can be nearly constant within a segment although it is not in the whole table: Contoso's customer country took a second value on 3 of the 1{,}500 sampled rows of one segment, and those 3 rows sat at the extreme longitudes. The correlations of that column are then a function of 3 rows (quasi-separation in a probit model), the posterior is weakly identified and the chain wandered ($\hat R=2.6$ after 3{,}000 iterations, chain means of one correlation from $-0.12$ to $-0.65$ over 4 seeds). The rule fixes it (3 rows are 0.2\% of the sample). The 1\% bound keeps the rule away from small segments, where a handful of rows is a sizeable share and independence would cost more than it saves; one segment of 28 still has a column (country, 16 and 2 rows away from the mode) whose correlations mix slowly, and the export reports it as a warning with the column pairs and the setting that removes it.

\begin{proposition}[Burn-in after the merge]\label{prop:mergeburn}
The merge step moves the segment chain to a new partition (a deterministic map of the state), after which the blocked Gibbs sampler (layer 1, Algorithm 3) needs a burn-in. The final phase is therefore monitored on the second half of its sweeps and extended in blocks while the split $\hat R$ of the attribute log-likelihood exceeds 1.1 (up to \texttt{gibbs.max\_final\_iterations}); the partition used is the last state, so the window changes only when the chain stops, not the kernel.
\end{proposition}
\begin{proof}
Discarding an initial segment of a chain does not change its stationary law; the stopping rule reads only past states.
\end{proof}
\noindent On Contoso the log-likelihood falls for about 25 sweeps after the merge and is then flat: monitoring the whole final phase gave $\hat R=1.23$ and a Geweke $z$ of 8 even after 120 sweeps; monitoring its second half gives $\hat R=1.00$ and $|z|<2$ after 60 sweeps.

\subsection{The exported model}
\texttt{relsynth export-model spec.yaml} writes \texttt{model.json} (structure, posterior mean and $M$ posterior draws of every parameter block, strict JSON), \texttt{model.md} and \texttt{model.pdf} (the generative model with this database's numbers and the verification table), \texttt{replicate.py} with \texttt{relsynth\_expr.py} (a standalone sampler needing NumPy, pandas and SciPy) and \texttt{run\_model.py} (the engine's own calls).

\begin{proposition}[Exported draw $m$ is the parameter of dataset $m$]\label{prop:exportdraw}
Draw $m$ of \texttt{model.json} is computed by the engine's parameter draw with the parameter stream of dataset $m$ (Proposition~\ref{prop:postpred}), so it equals the $\theta^{(m)}$ from which \texttt{relsynth run} generates dataset $m$; \texttt{run\_model.py} performs the calls of \texttt{relsynth run} and reproduces its output byte for byte.
\end{proposition}
\begin{proof}
Same function, same seed sequence child, same order of draws; JSON writes floats by their shortest round-trip representation, which parses back to the same double.
\end{proof}
\noindent \emph{Checks C72 and C73}: draws equal the engine's parameters for every dataset index tested; 18 of 18 files of two datasets are identical.

\begin{proposition}[Verification conditions]\label{prop:verify}
Before the model is written, each parameter set (the mean and every draw) is checked against the conditions under which the generative process is defined and the propositions apply: (i) every probability vector on the simplex, every probability in $[0,1]$, every rate non-negative and every quantile function non-decreasing (proper conjugate posteriors, Propositions~\ref{prop:eta}, \ref{prop:zoib}, \ref{prop:sparsedir}, \ref{prop:dp}); (ii) every copula correlation matrix symmetric with unit diagonal and positive semi-definite, the condition of Sklar's theorem for a Gaussian copula (Propositions~\ref{prop:pit}, \ref{prop:copula}); (iii) every block with a positive rate has source partner rows, the condition of the scale rule (Proposition~\ref{prop:scale}); (iv) hierarchies drawn from level counts, so their size is finite whatever the spectral radius of the offspring matrix (Propositions~\ref{prop:treesize}, \ref{prop:treelevels}); (v) the Markov chains stationary by Propositions~\ref{prop:mergeburn} and \ref{prop:fixedwidth}. A failure of (i)--(iv) is reported at the top of \texttt{model.md} and makes \texttt{export-model} exit with an error status; (v) is reported as a warning, because the draws are still states of the right chains.
\end{proposition}
\begin{proof}
Each condition is a hypothesis of the proposition cited.
\end{proof}

\begin{proposition}[The replica samples the core generative process]\label{prop:replica}
Given $\theta$, \texttt{replicate.py} draws: entity segment sizes by largest remainder of $n\pi$; latent segments; attributes from the segment's Gaussian copula and marginals (given parent attributes, from the conditional normal with the given latent values uniform on their bins); hierarchies from the level-by-segment counts; links by segment and degree weight or DP counts; fact rows by the Poisson block or occupancy model; time columns from the anchor segment's period profile; then sequences, offset keys, lookups, derived columns and aggregates with the engine's own expression evaluator. Each step draws from the same conditional law as the engine, so the replica's output has the engine's law for the parts it covers; it is not bit-identical (different random streams). The parts it does not cover are listed in \texttt{replica\_scope} (rules imposed by rejection or moves, slowly changing versions, affinity and participation moves, group effects, text and identifier models, nested facts and virtual dimensions).
\end{proposition}
\begin{proof}
Step by step, the laws are those of Propositions~\ref{prop:pit}, \ref{prop:copula}, \ref{prop:eta}, \ref{prop:treelevels}, \ref{prop:time} and \ref{prop:computed} given $\theta$.
\end{proof}
\noindent \emph{Check C73}: no orphan keys and entity sizes exact on S9, S10 and S11; replica scorecards 104/105, 139/143 and 85/91 (the misses are the engine-only parts in the scope list).

\begin{proposition}[Independent check with Stan]\label{prop:stan}
For the rate, occupancy, category, segment-attribute and copula blocks, the engine's posterior draws and those of a Stan program written in textbook form (prior and likelihood only) agree in mean within four Monte Carlo standard errors and in standard deviation within 20\%.
\end{proposition}
\begin{proof}
By computation (\texttt{relsynth.stancheck}); the two implementations share no code.
\end{proof}
\noindent \emph{Check C74}: 11 of 11 parameters agree (the check is skipped when Stan is not installed).

%=====================================================================
\section{Review: pooling, exact conditioning and calibrated limits (v0.9)}\label{sec:v09}
Version 0.9 is a review of the whole engine. Its common thread is a pattern found in the issues of earlier runs: \emph{many units, each with its own fixed prior, threshold or test, and no pooling or multiplicity control}, give a steady trickle of small errors that no single unit shows. Examples are structural blocks drawn free in complete relations (bi B7), small segments given empty parents (AdventureWorks customers without a sale), a maximum over 48 cells compared with a fixed $z=3$ (Contoso), a forest test whose null does not fit code columns, and per-check limits whose family produces false failures. The propositions below give the corrected models, samplers and limits. Checks \texttt{v09\_core.py} (14), \texttt{v09\_eval.py} (14), \texttt{v09\_ident.py} (14) and \texttt{v09\_io.py} (28), run as C75--C78, test this section.

\subsection{Structural blocks pooled across a relation}
\begin{proposition}[Learned structural weights]\label{prop:structpool}
Let the rows $k$ of a relation fall in groups $g$ (rows that share their structure; an ungrouped row is a group of its own), and let block $(k,l)$ have $S_{kl}$ successes in $T_{kl}$ trials. Column $l$ of group $g$ has an indicator $G_{gl}\in\{0,1,s\}$ with weights $w$: structural zero (every block of the group-column has $p=0$), structural one ($p=1$), or free; $w\sim\Dir(\kappa\pi)$, $\pi=(\pi_s/2,\pi_s/2,1-\pi_s)$, exchangeable over the relation's group-columns. In a free group-column of two or more rows each block has its own indicator $H_{kl}$ with the fixed weights $\pi$ (as in v0.7--v0.8); a one-row group-column is its block ($H=G$). A free block has $p_{kl}\sim\mathrm{Beta}(a_k,b_k)$. With $m_0=\mathbf 1[S=0]$, $m_1=\mathbf 1[S=T]$, $m_s=B(a_k+S,b_k+T-S)/B(a_k,b_k)$ for one block and $\bar m=\sum_h\pi_h m_h$:
(a) the blocked Gibbs sampler that draws, given $w$, each $G_{gl}$ with
\[P(G_{gl}=0)\propto w_0\prod_{k\in g}m_0(S_{kl},T_{kl}),\quad P(G_{gl}=1)\propto w_1\prod_{k\in g}m_1,\quad P(G_{gl}=s)\propto w_s\prod_{k\in g}\bar m(S_{kl},T_{kl}),\]
and then $w\sim\Dir(\kappa\pi+n(G))$, leaves the posterior of $(w,G)$ invariant, with the $H$ integrated out;
(b) an indicator without exposure ($T=0$ throughout) has all marginals equal to 1, so it is independent of the data and integrates out exactly: running the chain on exposed group-columns only has the same stationary law for $w$, and mixes faster;
(c) the draw of a dataset decides each group-column with the last $w$ of the chain and, inside a free group-column, each block with $\pi$ (\texttt{grouped\_zoib\_draw}); $\kappa=3$ by default (\texttt{structure.structural\_pooling}); $\kappa=0$ restores the fixed weights $\pi$ of v0.7--v0.8 everywhere.
\end{proposition}
\begin{proof}
(a) Given $w$ the likelihood of group-column $(g,l)$ under $G=0$ or $G=1$ is the product of its blocks' spike likelihoods; under $G=s$ the blocks' indicators are independent with weights $\pi$, and summing them out gives $\prod_k\bar m$, which does not depend on $w$. Given the $G$, $w$ has the Dirichlet--multinomial conjugate posterior. Both are full conditionals of the model with $H$ summed out, so the Gibbs sampler leaves that posterior invariant. (b) With $T=0$ every marginal is 1, so the indicator's conditional is $w$ itself and its factor in the posterior of $w$ sums to 1 over its values. (c) is a definition; given $G=s$, drawing $H$ with $\pi$ is its exact conditional.
\end{proof}
\noindent \emph{Why only the group level is learned.} The group-columns are the units the draw decides on, and they are exchangeable within a relation; the blocks inside a free group-column are not exchangeable with them, because their indicators also absorb how the segments mix rows of different structure. An earlier version learned one $w$ from both levels. On bi (real data) the many partial blocks inside free group-columns drove $w_1$ to 0.14 (0 in one relation), so complete blocks with one or two trials were drawn from the slab with probability up to 0.6, and 4--7\% of the reported values were lost (cardinality KS 0.26--0.46 against a limit of 0.12). With the row level kept at $\pi$, the expected number of present pairs equals that of fixed weights (62{,}498 against 62{,}481 of 62{,}761 observed in bi's largest relation), and $w_1$ learned on grouped thin rows equals the share of complete group-columns ($0.290$ against $0.292$, \texttt{t\_zoib\_pooling}). In a relation whose 112 observed blocks are all complete (bi-like, ungrouped), fixed weights drew a free block somewhere in 99.7\% of datasets (5.56 free blocks per dataset); learned weights draw 0.00. Data drawn from the model itself recover its weights within 3 posterior standard deviations. In the reporting scenario a filed header form is complete in 99.8\% of datasets and 0.01\% of item values carry another version (C26, C27); \texttt{t\_fixed\_weights\_option} shows that $\kappa=0$ gives back v0.8's exact structure.

\begin{proposition}[The slab is estimated from partial entries]\label{prop:partialslab}
The free component of Prop.~\ref{prop:structpool} must put little mass where the spikes are for the weights to be identified (George and McCulloch 1993). The empirical-Bayes slab of row $k$ (Kleinman's moment estimator) is therefore estimated from the entries with $0<S<T$ only, the entries that only a free block explains; a row with fewer than three such entries takes the slab estimated from the partial entries of the whole relation (rows exchangeable a level up), and $\mathrm{Beta}(a_0,b_0)$ is kept only when the relation itself has fewer than three. Up to v0.8 complete entries ($S=T$) entered the estimate, so a row of complete entries had a slab concentrated near 1, indistinguishable from the structural one: with weights learned from the data, complete blocks moved to the slab and were drawn with $p\approx0.99$ (a header form not complete in 2\% of filings).
\end{proposition}
\begin{proof}
The moment estimator is that of Prop.~\ref{prop:structpool}'s slab restricted to entries whose likelihood under both spikes is zero, so the estimate is a function of the free blocks only. If the slab has mass $\epsilon$ within $\delta$ of 1, a complete block of $T$ trials has slab marginal at least $\epsilon(1-\delta)^T$, and the posterior odds of one against free are $w_1/(w_s\,m_s)$, which tend to $w_1/w_s$ as the slab concentrates at 1: the data then cannot move the weights between the two components.
\end{proof}

\subsection{Counts, inclusion and participation}
\begin{proposition}[P\'olya-urn marginal in closed form]\label{prop:dpclosed}
For counts $t_1..t_n$ under $\mathrm{DP}(a,F_0)$ with $F_0=\mathrm{Poisson}(\lambda)$, $f_0$ its pmf, let $m_x$ be the multiplicity of value $x$ and $b_x=a f_0(x)$. Then
\[\log p(t)=\sum_x\Bigl[\log b_x+\log\Gamma(b_x+m_x)-\log\Gamma(b_x+1)\Bigr]-\log\frac{\Gamma(a+n)}{\Gamma(a)}.\]
\end{proposition}
\begin{proof}
By the urn, draw $i$ equals $x$ with probability $(b_x+c)/(a+i)$, $c$ the earlier copies of $x$. The numerators for the $m_x$ copies of $x$ multiply to $\prod_{c=0}^{m_x-1}(b_x+c)=b_x\Gamma(b_x+m_x)/\Gamma(b_x+1)$ and the denominators to $\Gamma(a+n)/\Gamma(a)$. Computing $\log b_x$ directly avoids the underflow of $a f_0(x)$ far in the Poisson tail, which made v0.8 return $-\infty$ (and a spurious Bayes factor) for a single large count.
\end{proof}

\begin{proposition}[Inclusion grid for the sparse Dirichlet]\label{prop:sdgrid}
In Prop.~\ref{prop:sparsedir}'s representation the auxiliary $u\in(0,1)$ has density $\propto u^{A_o-1}(1-u)^{n-1}\prod_j(1+\rho u^{\alpha_j})$. With $y=\log(-\log u)$ the density of $y$ is
\[g(y)\propto\exp\bigl(A_o\ell+(n-1)\log(1-e^{\ell})+y\bigr)\prod_j\bigl(1+\rho e^{\alpha_j\ell}\bigr),\qquad\ell=-e^{y},\]
which decays like $\exp(-A_oe^{y})$ as $y\to\infty$ and like $e^{ny}$ as $y\to-\infty$; a grid on $[-60,40]$ holds all but a negligible mass for every $A_o\ge10^{-12}$ and $n\ge1$.
\end{proposition}
\begin{proof}
$u=\exp(-e^y)$, $|du/dy|=u\,e^y$, which gives $g$. As $y\to\infty$, $\ell\to-\infty$ and the factor $u^{A_o}=e^{A_o\ell}$ dominates; as $y\to-\infty$, $1-u=1-e^{\ell}\sim e^{y}$, so $(1-u)^{n-1}e^{y}\sim e^{ny}$. The v0.8 grid on $\mathrm{logit}\,u$ decays only like $e^{A_oy}$ at $-\infty$ and lost most of the mass for small $A_o$.
\end{proof}
\noindent Check \texttt{t\_sparse\_dirichlet} compares the sampler's expected number of included unobserved levels with the exact posterior $P(|J|=j)\propto\binom mj\rho^j\Gamma(A_o+j\alpha)/\Gamma(A_o+j\alpha+n)$ in three regimes.

\begin{proposition}[Pooled participation prior]\label{prop:particippool}
The shares $s_{0t}$ of parents without children in the segments $t$ of a link are exchangeable: $s_{0t}\sim\mathrm{Beta}(cm,c(1-m))$ with $m=(\sum Z+\frac12)/(\sum N+1)$ and $c$ by the Beta--binomial moment estimator in $[1,10^4]$; the per-dataset draw is $s_{0t}\sim\mathrm{Beta}(cm+Z_t,c(1-m)+N_t-Z_t)$. When no source parent is empty, $E[s_{0t}]\le(\frac12/(\sum N+1))\,c/(c+N_t)\cdot(1+o(1))$, instead of $1/(N_t+2)$ under the flat per-segment prior of v0.7--v0.8.
\end{proposition}
\begin{proof}
Conjugacy gives the posterior. With $Z=0$ and $\tau^2\le0$ the estimator gives $c=10^4$, and the posterior mean is $cm/(c+N_t)$.
\end{proof}
\noindent The flat prior gave every small segment about $1/(N_t+2)$ of empty parents, which is the 3\% of AdventureWorks customers without a sale reported since v0.5; the fake order warehouse now has every customer with an order in every dataset (\texttt{t\_mandatory\_participation}).

\begin{proposition}[Covering mandatory parents]\label{prop:cover}
When every parent needs at least $\mu$ children, the parent segments of the child rows are drawn i.i.d.\ from $\phi_{z}$ and accepted when every segment $t$ receives $N_t\ge\mu M_t$ rows (rejection, up to \texttt{rules.max\_rejection\_rounds} times): an accepted draw has exactly the conditional law given the constraint. If no draw is accepted, rows are moved from segments with a surplus, the row moving from $s$ to $t$ chosen with probability proportional to $\phi_{z,t}/\phi_{z,s}$, until every deficit is filled; this moves exactly the total deficit and changes nothing when the constraint already holds.
\end{proposition}
\begin{proof}
Rejection sampling from a product law restricted to an event yields the restricted law. The repair stops when every $N_t\ge\mu M_t$; each move lowers one deficit by one and never takes a source below its own need, so the number of moves is $\sum_t(\mu M_t-N_t)^+$.
\end{proof}

\subsection{Exact conditioning on rules}
\begin{proposition}[Rules are checked on the draws that are kept]\label{prop:hidden}
Let a fact row be $X=(Y,D)$ with $Y$ the modelled columns and $D$ the independent per-row draws (stated distributions, soft-rule indicators and offsets, group-effect priorities). Rejection per row that redraws $(Y,D)$ of invalid rows, keeps $(Y,D)$ of valid ones and evaluates the rules on that same $(Y,D)$ yields rows with law $Q(X\mid\text{valid})$. v0.8 drew $D$ afresh on a copy for the check and again for the output, so the output's $D$ was independent of the check: rows with $\mathrm{qty}<2$ under the rule $\mathrm{qty}\ge2$ (415 of 1{,}978 in the probe) were output.
\end{proposition}
\begin{proof}
Each row's accepted value is the first valid one in an i.i.d.\ sequence from $Q$, which has law $Q(\cdot\mid\text{valid})$. The proof needs the checked and the returned value to be the same draw. The engine stores $D$ in hidden columns with a per-row mask, draws only rows whose mask is unset (new or redrawn rows), and drops the hidden columns at the end of the dataset; recomputing a frame (as the SCD step does) reuses them.
\end{proof}

\begin{proposition}[Unique keys over a link and attributes]\label{prop:mixedkeys}
The attributes of an entity's rows are drawn before its links. A unique key made of links and modelled attributes is therefore a constraint on the links given the attributes, and the Metropolis--Hastings kernel of the link moves (independence proposals from each row's own factor, accepted iff the new state is valid) targets the product law restricted to it, as for keys over links only. Attribute redraws in the finishing step keep it as part of their rejection.
\end{proposition}
\begin{proof}
Given the attributes, the key's attribute part is fixed, so the set of valid link assignments is a fixed set; the acceptance ratio of an independence proposal from $Q_i$ restricted to a fixed set is 1 on valid states, as in the keys-over-links case.
\end{proof}
\noindent v0.8 kept keys made only of links (link moves) or only of attributes (finishing), and silently ignored mixed keys of entities (26 duplicates of (product, balance) in \texttt{t\_unique\_moves\_fast}).

\subsection{Copula sample size}
\begin{proposition}[Pooled copula sample]\label{prop:pooledcells}
The posterior standard deviation of a latent correlation fitted on $n_{\rm fit}$ rows is $\approx(1-\rho^2)/\sqrt{n_{\rm fit}}$ (measured 0.0171, 0.009, 0.0047 for $n=3$k, 12k, 48k). The pooled copula of a table of $n$ rows and $p$ columns is now fitted on $\max(2\,\mathrm{cap},\lfloor C/p^2\rfloor)$ rows, $C=$\,\texttt{copula.pooled\_cells} ($10^6$), within the $O(np^2)$ cost of Prop.~\ref{prop:fixedwidth}'s sampler, so that the dataset-to-dataset spread of a weak correlation is not several times its sampling error.
\end{proposition}
\begin{proof}
The rank likelihood's Fisher information for $\rho$ is $n/(1-\rho^2)^2$ to first order (Hoff 2007), and each sweep costs $O(np^2)$, so $C/p^2$ rows fix the cost per sweep.
\end{proof}
\noindent Contoso's join association $z=3.57$ (Quantity against order date) was a single dataset's posterior draw of $+0.021$ for a correlation whose posterior mean is $-0.022$, fitted on 3{,}000 of 224{,}000 rows.

\subsection{Identification}
\begin{proposition}[Permutation test of a forest]\label{prop:forestperm}
Let $T(\pi)=1$ when the column's pointers permuted by $\pi$ (nulls included) form a forest. Under the null that the column's values are exchangeable over the rows, $p=(1+\#\{b\le B:T(\pi_b)=1\})/(B+1)$ with $\pi_1..\pi_B$ uniform permutations satisfies $P(p\le\alpha)\le\alpha$ for every $\alpha$. The engine uses $B=199$, $\alpha=0.01$ (smallest $p=0.005$), stops early once $p>\alpha$ is certain, and accepts the self-reference when $p\le\alpha$; otherwise it asks when no row points to itself.
\end{proposition}
\begin{proof}
Under exchangeability the observed column is one more draw from the permutation distribution, so its rank among the $B+1$ values is uniform (ties broken in its favour), which gives the stated $p$-value its validity (Phipson and Smyth 2010). Stopping when $1+\#>\alpha(B+1)$ only skips draws that cannot change the decision.
\end{proof}
\noindent A code column with values $1..3$ in a 2{,}000-row table was accepted as a hierarchy in 124 of 200 samples under Prop.~\ref{prop:forest}'s uniform null; with the permutation test, 0 of 60.

\begin{proposition}[Inclusion of strings by chance]\label{prop:stringchance}
Let the $v\ge3$ distinct child strings in $K$ use the length set $\mathcal L$ and character classes $\Sigma$ (capitals, small letters, digits, each counted if used, plus every other character seen), $A=|\Sigma|$, $w=\sum_{L\in\mathcal L}A^L$ and $k=|K\cap\Sigma^{\mathcal L}|$. Under the null that the $v$ values are a uniform draw without replacement from the code space, $P=\binom kv/\binom wv\le(k/w)^v$ ($P=1$ when $k\ge w$; $P=1$ for $v\le2$). Inclusion is evidence when $P\le10^{-6}$.
\end{proposition}
\begin{proof}
Hypergeometric, as Prop.~\ref{prop:inclchance}, with the bound from $(k-i)/(w-i)\le k/w$. v0.8 returned $P=0$ for any three strings, so three section codes $\{A,B,C\}$ inside grades $A..F$ were a link; now $\log_{10}P=-2.11$.
\end{proof}

\noindent \textbf{Names need values.} For a link between two tables, names count as evidence only with randomness $R\ge R_{\rm crit}(v)=\min(0.5,1-2/\sqrt v)$ over the $v$ included values, or $P\le10^{-6}$; the same $R_{\rm crit}$ replaces the fixed 0.6 of the question rule. $R_{\rm crit}(v)$ lies below the simulated 0.1\% quantile of $R$ for random samples of $v=5,10,20,50$ values. Composite keys and self-references are now confirmed on the whole table, not on a sample.

\subsection{Calibrated limits of the scorecard}
\begin{proposition}[Posterior-predictive null of a distribution difference]\label{prop:ppnull}
Let the synthetic sample of size $m$ be drawn from a Bayesian bootstrap of the real sample of size $n$ (the posterior predictive of a saturated model). For each $x$, $E[F_{\rm syn}(x)\mid\text{real}]=F_n(x)$ and
\[\mathrm{Var}\bigl(F_{\rm syn}(x)-F_n(x)\mid\text{real}\bigr)=F_n(x)\bigl(1-F_n(x)\bigr)\Bigl(\frac1{n+1}+\frac1m\Bigr).\]
The two-sample effective size $1/(2/n+1/m)$ of v0.8 therefore overstated the null spread by $\sqrt{(2/n+1/m)/(1/(n+1)+1/m)}$ ($\approx1.22$ at $n=m$); KS, TVD, null-share, association and join limits now use $1/(n+1)+1/m$ with the effective (Kish) sizes.
\end{proposition}
\begin{proof}
With Dirichlet$(1,\dots,1)$ weights $W$, $G(x)=\sum_iW_i\mathbf 1[x_i\le x]$ has mean $F_n(x)$ and variance $F_n(1-F_n)/(n+1)$ (the variance of a sum of $\Dir(1^n)$ components over a subset of size $nF_n$). Given $G$, $F_{\rm syn}(x)$ is a binomial share with variance $G(1-G)/m$, whose expectation is $F_n(1-F_n)(1-\frac1{n+1})/m$; to first order the total is as stated.
\end{proof}

\begin{proposition}[Maximum over a family]\label{prop:sidak}
For $K$ centred Gaussian statistics $Z_1..Z_K$ with unit variances and any correlations, $P(\max_k|Z_k|\le z)\ge\prod_kP(|Z_k|\le z)$ (\v{S}id\'ak 1967). The limit $z_K=\Phi^{-1}\bigl(1-\frac12(1-(1-\alpha)^{1/K})\bigr)$ therefore keeps the family's false-failure rate at most $\alpha$; for $K=48$ and $\alpha=0.01$, $z_K=3.71$ (v0.8 used 3 for the maximum over up to 48 cells).
\end{proposition}
\begin{proof}
\v{S}id\'ak's inequality, and $\prod_kP(|Z_k|\le z_K)=(1-(1-(1-\alpha)^{1/K}))^K=1-\alpha$.
\end{proof}
\noindent The same correction is used for the worst association pair and across tables for the discriminator, whose null standard deviation is bounded by that of one fold's Mann--Whitney AUC, $\sqrt{(n+1)/(3n^2)}$ (within a fold the scores come from a model fitted on the other fold, so the Mann--Whitney null law is exact; the mean of two correlated folds has at most that variance). Novel combinations use the Good--Turing missing mass $N_1/n$ with Esty's variance (Good 1953; Esty 1983), and unequal expected cardinalities compare binomially thinned counts (thinning a Poisson count keeps it Poisson, R\'enyi 1956). The scorecard reports the expected number of false failures, $\sum\alpha_i$, so a run of 1{,}303 checks is read against about $0.01\times1303\approx13$ expected chance failures.

%=====================================================================
\section{Fitting inside the database (v1.0)}\label{sec:v10}
Version 1.0 packages the engine for production use and fits fact tables larger than memory from statistics computed where the data lives (\emph{pushdown}). A fact table is \emph{pushed down} when it is a fact that no table references, has no type~2 history, and is not the source of an aggregate column that the source does not store (\texttt{relsynth/pushdown/planner.py}). Its rows never leave the database as a whole: the engine reads (i) exact aggregates computed by \texttt{GROUP BY} and window functions, and (ii) a uniform sample of at most \texttt{synthesis.pushdown.sample\_rows} rows. Every other table is loaded as before. The propositions below show that this changes nothing in the relational model, and that the attribute models fitted on the sample are an exact posterior given the sampled rows. The checks are in \texttt{tests/unit/test\_sqlfuncs\_stats.py}, \texttt{tests/unit/test\_pushdown.py} and \texttt{tests/integration/test\_servers.py}; they run on SQLite, PostgreSQL~16, SQL Server~2022 and Oracle Database~26ai Free (23.26).

\subsection{What the relational model needs from a fact}
Let a pushed fact $F$ have $N$ rows, key columns $k_1..k_m$ (links to dimensions or to a parent fact, the parts of a composite key) and continuous columns $c_1..c_p$. For a non-null value $x_i$ of column $c$ with average rank $r_i$ among the $N_c$ non-null values let $z_i=\Phi^{-1}\bigl(r_i/(N_c+1)\bigr)$ (its \emph{normal score}), and $z_i=0$ for a null. For a key tuple $u$ write
\[n(u)=\#\{i:k(i)=u\},\qquad S_c(u)=\sum_{i:k(i)=u}z_i^{(c)},\qquad Q_c(u)=\sum_{i:k(i)=u}\bigl(z_i^{(c)}\bigr)^2,\]
with a null key value forming a group of its own.

\begin{proposition}[Sufficiency of the in-database statistics]\label{prop:pdsuff}
Given the loaded tables (entities and non-pushed facts), the full conditional of every latent quantity of the relational model of Sections~\ref{sec:star}--\ref{sec:v09} that involves a pushed fact $F$---segment assignments of the entities it links, block and link rates (Props.~\ref{prop:eta}, \ref{prop:split}), count and multiplicity laws (Props.~\ref{prop:mult}, \ref{prop:countbf}, \ref{prop:dpclosed}), participation and inclusion (Props.~\ref{prop:mandatory}, \ref{prop:particippool}, \ref{prop:sparsedir}), nested refinement (Prop.~\ref{prop:nested}), order headers (Prop.~\ref{prop:header}) and running balances (Prop.~\ref{prop:cumulative})---depends on the rows of $F$ only through
\begin{enumerate}
\item the counts $n(u)$ over the key tuples the model uses (pairs of links, a composite key with its links, a header key with its header columns);
\item the score sums $S_c(u)$ and $Q_c(u)$ for the (at most three) continuous columns that segment discovery uses;
\item for an order header, the first row of every document key in primary-key order and the number of cells that differ from it; for a running balance, the first row of every group in the rule's order, the number of non-null values and the minimum;
\item $N$, the number of non-null values of each key, and the distinct referenced values of each key.
\end{enumerate}
\end{proposition}
\begin{proof}
Each listed model block is a product over the rows of $F$ of a factor that depends on a row only through its key tuple $u$ and, for segment discovery, through the rank features $z^{(c)}$ of the row. Block-rate and multiplicity likelihoods are products of Poisson or multinomial terms whose exponents are the counts of rows per pair of linked rows, so they are functions of $n(u)$ (the conjugate posteriors of Props.~\ref{prop:eta} and \ref{prop:mult} use sums of $n(u)$ over segment pairs). The rank-feature factor of segment discovery is Gaussian in the scores, with a mean per pair of segments and a common variance per column (\texttt{PairRel}, \texttt{OccRel}), so its log-likelihood over the rows of a key tuple is $a\,n(u)+b\,S_c(u)+c\,Q_c(u)$ for coefficients that do not depend on the rows, and the conjugate updates of the means and variances use the same three sums. Participation and inclusion use which keys occur (the distinct values) and how often ($n(u)$). A nested fact is linked to its parent fact through its composite key, so the counts over (composite key, links) give the parent row and the dimension rows of every group, with the multiplicity $n(u)$ as the group's weight. An order header is the first row of its document key (Prop.~\ref{prop:header}); its split is lossless exactly when no cell differs from that row, which is what (3) counts. A running balance needs its opening value per group (the first row in the rule's order) and its floor (the minimum); the rest of the column is a deterministic function of the flows (Prop.~\ref{prop:cumulative}). Nothing else of $F$ enters these factors.
\end{proof}

\begin{proposition}[Agreement with in-memory fitting]\label{prop:pdexact}
Let the loaded rows of every table be put in a canonical order: by primary key, else by the first declared unique key that holds, and by a hash of the row's values when there is neither (\texttt{canonical\_order}); and let order headers split from a fact be put in the order of their document key. If the sample of a pushed fact holds all its rows, the fitted model and every synthetic dataset are identical, value for value, to those of fitting the whole table in memory with the same seed; and the model fitted from a database does not depend on the order in which the database returns rows.
\end{proposition}
\begin{proof}
By Prop.~\ref{prop:pdsuff} the relational blocks see the same numbers, computed either by \texttt{FrameStats} (pandas) or by \texttt{SqlStats} (SQL); the attribute models see the same rows. The fit is a deterministic function of its inputs, the seed and the order in which rows are visited (the Gibbs sweeps visit rows in order). The canonical order (by primary key, else by the first declared unique key, else by a hash of the row; order headers split from a fact by their document key) is a function of the set of rows only, so two backends, or two scans of one table, give the same input sequence. The only remaining difference is floating-point: SQL sums are accumulated in another order than in pandas, and servers compute $\Phi^{-1}$ with Prop.~\ref{prop:acklam}; both perturb $S_c,Q_c$ by relative errors below $10^{-8}$, far under the resolution at which a segment draw changes. The checks compare the synthetic tables cell by cell: identical for all four pushed fact shapes on SQLite (\texttt{test\_full\_sample\_pushdown\_equals\_in\_memory}) and for a pushed star schema on PostgreSQL, SQL Server and Oracle against SQLite (\texttt{test\_full\_sample\_pushdown\_matches\_sqlite\_in\_memory}).
\end{proof}
\noindent Without the canonical order the same database gave different models on different servers (SQL Server and Oracle return rows in a different order than SQLite and PostgreSQL), and a server using parallel scans could give different models on two runs. v1.0 orders every loaded table; this changes in-memory results relative to v0.9 by a permutation of the input rows, which leaves the posterior unchanged and changes individual draws only.

\subsection{Normal scores on database servers}
\begin{proposition}[Acklam's quantile approximation]\label{prop:acklam}
PostgreSQL, SQL Server and Oracle have no standard-normal quantile function. relsynth computes $\Phi^{-1}(p)$ in SQL with Acklam's rational approximation (a rational function of $q=p-\tfrac12$ on $[0.02425, 0.97575]$ and of $\sqrt{-2\log p}$ in the tails), whose relative error is below $\delta=1.15\times10^{-9}$ on $(0,1)$. Hence $|\hat S_c(u)-S_c(u)|\le \delta\sum_{i:k(i)=u}|z_i|$ and $|\hat Q_c(u)-Q_c(u)|\le (2\delta+\delta^2)\,Q_c(u)$. SQLite registers the exact function (\texttt{scipy.special.ndtri}) as a user function.
\end{proposition}
\begin{proof}
The error bound of the approximation is Acklam's (2003), and \texttt{test\_acklam\_error\_bound} confirms it in double precision on 4{,}081 points spanning $[10^{-12},1-10^{-12}]$ (relative error where $|\Phi^{-1}(p)|>10^{-3}$, absolute error below $10^{-8}$ everywhere). For a sum, $|\sum\hat z_i-\sum z_i|\le\sum|\hat z_i-z_i|\le\delta\sum|z_i|$ with $\delta=1.15\times10^{-9}$; for squares $|\hat z^2-z^2|=|z|^2\,|(1+\epsilon)^2-1|\le(2\delta+\delta^2)z^2$.
\end{proof}

\subsection{The sample}
\begin{proposition}[Seeded hash sampling is a Bernoulli sample]\label{prop:hashsample}
The sample keeps row $i$ when $h(s,\kappa_i)<\lfloor q\,2^{32}\rfloor$, where $\kappa_i$ is the row's primary key (all its columns when there is none), $s$ the seed, $q=\min(1,n/N)$ for the requested size $n$, and $h$ the database's hash reduced to $[0,2^{32})$ (\texttt{hashtextextended} on PostgreSQL, SHA-256 by \texttt{HASHBYTES} on SQL Server, \texttt{ORA\_HASH} on Oracle, BLAKE2b on SQLite and files). If $h(s,\cdot)$ behaves as a uniform random function independent of the data, then
(a) the inclusion indicators are i.i.d.\ Bernoulli$(q)$ and independent of every column value, so the sample is a uniform random subset with $|S|\sim\mathrm{Bin}(N,q)$;
(b) the same seed gives the same sample on every run and every engine with the same hash;
(c) $P\bigl(\bigl||S|-Nq\bigr|>8\sqrt{Nq(1-q)}+10\bigr)<1.3\times10^{-15}$, so the guard that stops a sample of that size (a broken hash, such as one returning NULL) does not stop a valid sample in practice.
\end{proposition}
\begin{proof}
(a) Under the assumption, $h(s,\kappa_i)$ are i.i.d.\ uniform on $[0,2^{32})$ because the keys are distinct, so each comparison is Bernoulli with probability $\lfloor q2^{32}\rfloor/2^{32}$, which differs from $q$ by less than $2^{-32}$. (b) The comparison is a deterministic function of $(s,\kappa_i)$. (c) With $\sigma^2=Nq(1-q)$ and $t=8\sigma+10$, Bernstein's inequality gives $P(||S|-Nq|>t)\le 2\exp\bigl(-t^2/(2\sigma^2+2t/3)\bigr)\to2e^{-32}<2.6\times10^{-14}$ as $\sigma\to\infty$; the exact binomial tails, computed for $N\in\{10,\dots,10^9\}$ and $q\in[10^{-6},1-10^{-6}]$, are at most $1.24\times10^{-15}$, the normal value $2\Phi(-8)$ approached for large $\sigma$ (for small $\sigma$ the additive 10 makes the tail far smaller). \texttt{test\_python\_hash\_is\_uniform\_and\_seeded} tests uniformity, and \texttt{test\_sample\_is\_uniform\_reproducible\_and\_sized} size and reproducibility on each engine.
\end{proof}

\begin{proposition}[Attribute models from the sample are an exact posterior]\label{prop:pdsample}
Let the sample be drawn as in Prop.~\ref{prop:hashsample}. The posterior of the attribute-model parameters $\theta$ of the pushed fact given the sampled rows,
\[p(\theta\mid x_S,S)\propto p(\theta)\prod_{i\in S}p(x_i\mid\theta,\text{keys}_i),\]
is the exact posterior given the observed data: the sampling mechanism is ignorable. Synthetic datasets are therefore posterior-predictive draws (Prop.~\ref{prop:postpred}) given the sampled rows for the attributes and given all rows for the relational part; their sizes follow the exact $N$, not $|S|$. The posterior is wider than with all rows (its variance scales with $1/|S|$ for regular parameters), so the scorecard's real side uses exact statistics where they exist (row counts, null keys, cardinalities, duplicate rates, pair multiplicities) and the sample elsewhere, with limits computed for the sample's size.
\end{proposition}
\begin{proof}
The inclusion indicators are independent of the values and of $\theta$ (Prop.~\ref{prop:hashsample}a), so $p(S\mid x,\theta)=p(S)$ and the likelihood of the observed data factors as $p(S)\prod_{i\in S}p(x_i\mid\theta)$; the factor $p(S)$ does not involve $\theta$ (Rubin 1976: data missing completely at random). Prop.~\ref{prop:pdsuff} makes the relational blocks depend on all rows through exact statistics. The sizes use $N$ because the scale rule (Prop.~\ref{prop:scale}) is applied to the exact counts.
\end{proof}

\subsection{Why entities are always loaded}
\begin{proposition}[Eligibility]\label{prop:pdeligible}
A table cannot be fitted from aggregates alone when any of the following holds: it is an entity (each of its rows carries its own latent segment, whose full conditional needs that row's attributes and links, Prop.~\ref{prop:perrow}); another table references it (it is then an entity of its own rows for that table's links, including the parent of a nested fact); it has type~2 history (Prop.~\ref{prop:scd2} draws versions per business key and row); or an aggregate column of another table is computed from its rows and the source does not store that column (the column must be recomputed from the rows). Every other fact is eligible, including nested facts, facts split into order headers, bases of views and facts with running balances.
\end{proposition}
\begin{proof}
For each excluded case the named proposition's full conditional has a factor per row of the table that is not a function of the statistics of Prop.~\ref{prop:pdsuff}. For the eligible shapes, Prop.~\ref{prop:pdsuff} lists the statistics that replace the rows.
\end{proof}

\noindent\emph{Memory and time.} A pushed fact costs memory proportional to the number of distinct key tuples plus the sample, instead of rows times columns: on AdventureWorks (all ten eligible facts pushed, including the 776{,}286-row inventory fact with its running balance) the peak resident memory fell from 2.1\,GiB to 1.6\,GiB at a sample of 30{,}000 rows, and the run took 668\,s instead of 803\,s on SQLite. The database does the grouping; the number of round trips is a constant per statistic and does not grow with $N$.

\subsection{Sizes of child entities}
An entity $C$ whose rows link to a parent entity $P$ (orders to customers) has its own size factor $s_C$; the link model draws each of its rows' parents from the segments of $P$. If only $P$ is resized, the expected number of $C$ rows per $P$ row changes by $s_C/s_P$, which is rarely what a user who asks for ``ten times the customers'' wants. From v1.0 an entity without a factor of its own takes the factor of the first parent it links to that has one (\texttt{inherit\_scale}, applied in topological order, so a chain of entities follows its root), and the scorecard expects the same sizes. With $s_C=s_P$ the expected rows per parent are those of the source, and facts anchored on $C$ follow by Prop.~\ref{prop:scale}. A table given its own factor keeps it. \emph{Test} \texttt{test\_child\_entities\_follow\_a\_scaled\_parent}: customers $\times3$ gives exactly $3\times$ the invoices and $3\times$ the lines within 10\%, with every invoice-per-customer check passing.

%=====================================================================
\section{The canonical intermediate representation (v2.0)}\label{sec:v20}
Version 2.0 replaces the relational spec as relsynth's internal description of the data by one modality-neutral
document, the \emph{IR} (\texttt{relsynth/ir}; reference \texttt{docs/ir-reference.md}, design record
\texttt{docs/adr/0001-canonical-ir.md}). An IR document $E$ (an \emph{estate}) lists entities with a shape (record,
document, node, edge, event), attributes (nested objects, arrays, polymorphic variants), keys, references, time,
physical layout, rules and privacy annotations. The engine of this release synthesises the relational part: it runs
the 1.0 spec $\lambda(E)$ that $E$ \emph{lowers} to. A 1.0 spec $s$ is \emph{lifted} to an IR document $\ell(s)$.
Write $\nu$ for the normalisation of a 1.0 spec (\texttt{spec.normalise}: defaults, shorthands, named
references), so that the engine is a function of $\nu(s)$, the data and the seed. The checks are in
\texttt{tests/unit/test\_ir.py}.

\subsection{Converting 1.0 specs}
\begin{proposition}[Loss-free conversion and identical output]\label{prop:irround}
For every 1.0 spec $s$ that the 1.0 key schema accepts, $\nu(\lambda(\ell(s)))=\nu(s)$. Hence, for the same data
and seed, the IR document $\ell(s)$ and the spec $s$ give the same fitted model and the same synthetic datasets.
Moreover $\ell(\lambda(\ell(s)))=\ell(s)$ up to the canonical form $\kappa$ (Definition~\ref{def:canon}), so a lifted
document is a fixed point of lowering and lifting.
\end{proposition}
\begin{definition}[Canonical form]\label{def:canon}
$\kappa(E)$ makes explicit what follows from other fields (an attribute's modelling kind from its logical type,
a reference's target attributes from the parent's primary key, the privacy role of a \texttt{pii} attribute) and
drops what is implied (a logical type that the kind fixes).
\end{definition}
\begin{proof}
$\ell$ and $\lambda$ are defined field by field by a fixed correspondence (module \texttt{ir/v1.py}, table in its
docstring). Every key of the 1.0 key schema has an image: either an IR field holding the same value
(\texttt{columns}$\to$\texttt{attributes}, \texttt{primary\_key}/\texttt{unique}/\texttt{partial\_unique}$\to$
\texttt{keys}, \texttt{scd}$\to$\texttt{time.valid\_time}+\texttt{time.versioning},
\texttt{fact.time\_column}$\to$\texttt{time.event\_time}, \texttt{dimensions}$\to$\texttt{dependencies}, \ldots) or
a value that $\nu$ itself would produce (a default relationship name, kind or parent key; a string column
shorthand; a soft rule written as a derive rule with \texttt{holds}$<1$). $\lambda$ inverts the field images and
emits defaults that $\nu$ fills in the same way. Order is kept: entities, attributes, keys, references and rules
are ordered maps or lists in $E$, and $\nu$ does not reorder them (the engine breaks ties in generation order by
these orders, so order is part of the meaning). Therefore $\nu$ applied to both sides sees the same values at the
same keys in the same order. The engine reads only $\nu(s)$, the data and the seed, so its output is the same.
For the fixed point: $\ell(\lambda(\ell(s)))$ reads back the values $\lambda$ wrote, which are those of $\ell(s)$
except where $\ell$ derives one from another; $\kappa$ normalises exactly those derivations.
\end{proof}
\noindent\emph{Checks.} \texttt{test\_real\_specs\_round\_trip}: the eleven 1.0 specs of the four golden databases
(Chinook, Contoso SQLite and Parquet, AdventureWorks DW, bi/COREP\_LR, including an introspect draft that does not
validate) and the two shipped examples convert both ways with $\nu(\lambda(\ell(s)))=\nu(s)$, as do their normalised
forms; \texttt{test\_fake\_scenario\_specs\_round\_trip} does the same for S1--S11. \texttt{test\_ir\_document\_gives\_identical\_synthetic\_data}
fits S2, S9 and S10 from $s$ and from $\ell(s)$ and compares the synthetic frames cell by cell;
\texttt{test\_cli\_runs\_an\_ir\_file} compares the written files byte by byte. The real databases were run the same way
(report \texttt{roadmap/reports/T02.md}). A lowering $\lambda(E)$ of a document that uses what a 1.0 spec cannot
hold returns the list of those paths, computed as the difference between $\kappa(E)$ and $\kappa(\ell(\lambda(E)))$, so
nothing is dropped in silence (\texttt{test\_lower\_reports\_annotations\_it\_drops}).

\subsection{The privacy unit}
Differential privacy on a database protects a \emph{unit}---a customer, a patient---rather than a row: neighbouring
databases differ in all the records of one unit (add/remove adjacency at the unit level; Ponomareva et al., JAIR 2026;
for foreign keys, Cai, Xiao and Cormode, PrivLava, SIGMOD 2023). Let $U$ be the unit entity. Let $G$ be the graph
whose vertices are the entities and whose arcs are the references $\rho:A\to B$ ($A\neq B$, many-to-one or
one-to-one) that are not hierarchies and not listed in \texttt{privacy.exclude\_paths}. For a path
$P=\rho_1\cdots\rho_k$ from entity $A$ to $U$ in $G$ write $o_P:\mathrm{rows}(A)\to\mathrm{rows}(U)\cup\{\bot\}$ for the
composition of the reference maps (a null or dangling reference gives $\bot$).

\begin{proposition}[Ownership along reference paths]\label{prop:irowner}
Under referential integrity, each $o_P$ is a function: a record of $A$ has at most one owner along each path. If $G$
is acyclic, the set $\mathcal P(A)$ of simple paths from $A$ to $U$ that \texttt{privacy.trace} returns is the set of
all paths, and the records of $A$ that belong to a unit $u$ are $C_u(A)=\bigcup_{P\in\mathcal P(A)}o_P^{-1}(u)$; a
record has at most $|\mathcal P(A)|$ owners. The status of every other entity is \emph{context} (reachable from an
owned entity: the unit's records reference it) or \emph{unrelated}.
\end{proposition}
\begin{proof}
A many-to-one or one-to-one reference maps each record to at most one parent record (the referenced attributes are a
key of the parent; validation requires this), so it is a partial function, and a composition of partial functions is
a partial function. In an acyclic graph every path is simple, and the depth-first search of \texttt{trace} enumerates
all simple paths from $A$ to $U$ (it extends every path by every outgoing arc to a vertex not yet on it). A record
belongs to $u$ when some chain of its references leads to $u$, which is the union over paths. The bound on owners
follows because each path contributes at most one. Context and unrelated entities have no path to $U$ by
construction. If $G$ has a cycle, a non-simple path can lead to a different owner than the simple ones; \texttt{trace}
reports the cycle as a warning (a 1.0 spec cannot have one: its validator requires a topological order).
\end{proof}

\begin{proposition}[Removing a unit keeps referential integrity]\label{prop:irremove}
Let $G$ be acyclic and suppose no reference outside $G$ (a hierarchy or an excluded reference) points into a record of
an owned entity or of $U$. Removing the units of a set $S$ and every record in $\bigcup_{u\in S}C_u(\cdot)$ leaves a
database in which every non-null reference resolves.
\end{proposition}
\begin{proof}
Suppose a kept record $r$ of $A$ references a removed record $t$ of $B$ through $\rho$. By the assumption $\rho$ is an
arc of $G$. As $t$ was removed, $o_P(t)\in S$ for a path $P$ from $B$ to $U$ (or $t$ is a unit in $S$, $P$ empty).
Then $\rho P$ is a path from $A$ to $U$; it is simple because $G$ is acyclic, so $\rho P\in\mathcal P(A)$ and
$o_{\rho P}(r)=o_P(t)\in S$, hence $r$ was removed: a contradiction.
\end{proof}
\noindent When the assumption fails (a customer who referred another, an employee's manager when employees are the
unit) the reference links two units; \texttt{trace} lists it as a cross-unit link and warns that removal leaves it
dangling, so the privacy design (T07) must null or bound it. \emph{Checks.}
\texttt{test\_privacy\_unit\_paths\_of\_the\_real\_databases} states the unit of each golden database and checks its
paths (an invoice line of Chinook belongs to its invoice's customer; a Contoso sale has two paths, directly and
through its order; a COREP reported value has two paths to its reporting entity);
\texttt{test\_contribution\_matches\_the\_joins} recomputes the owners of the S9 shop by joins along every path
(Prop.~\ref{prop:irowner}); \texttt{test\_removing\_a\_unit\_keeps\_referential\_integrity} removes every seventh
customer with its records and checks every reference (Prop.~\ref{prop:irremove});
\texttt{test\_hierarchy\_on\_the\_unit\_is\_a\_cross\_unit\_link} checks the warning.

%=====================================================================
\section{Generator plug-ins (v2.0)}\label{sec:v20gen}
% numbered from 700 so that the parallel 2.x threads (106-108 for the IR, 600-605 for privacy) do not collide
\newcounter{savedpropgen}\setcounter{savedpropgen}{\value{proposition}}
\setcounter{proposition}{699}
Version 2.0 separates the \emph{relational skeleton} of a synthetic database from the \emph{attribute columns} of
each table (\texttt{relsynth/generators}, guide \texttt{docs/generators.md}, design record
\texttt{docs/adr/0002-generator-plugins.md}). The skeleton $\mathcal S$---which rows exist, their keys, the links
between them, the latent segment of every row, counts and scale---is drawn by the engine of
Sections~\ref{sec:star}--\ref{sec:v10}. The attribute columns $x_t$ of table $t$ are drawn by a \emph{generator}
$G_t$ chosen in the spec: a pair of maps $\mathrm{fit}$ and $\mathrm{sample}$, where $\mathrm{sample}$ returns, for
rows with contexts $c_i$, values $x_i\sim q_{G_t}(\cdot\mid c_i;\theta_t)$. The context $c_i=(g_i,\bar g_i,a_i)$ is
the row's latent group $g_i$ (its segment, or its segment block for a fact), the coarser group $\bar g_i$ the engine
backs off to, and---for a generator that declares \emph{parent} conditioning---the attributes $a_i$ of the row's
synthetic parents. Write $\mathcal R$ for the rules of the spec, split into \emph{computed} rules (derive,
aggregate, cumulative, computed columns) and \emph{check} rules and unique attribute keys. The checks are in
\texttt{tests/unit/test\_generators.py}, and the identity runs and scorecards of every generator on every suite are
in the T03 report.

\subsection{What the engine guarantees with any generator}
\begin{proposition}[Skeleton and rules do not depend on the generator]\label{prop:genskeleton}
Let every $G_t$ satisfy the contract: $\mathrm{sample}(n,c,\rho)$ returns $n$ rows with exactly the table's own
attribute columns. Then, whatever the generators,
\begin{enumerate}
\item every synthetic dataset has unique, non-null primary keys; every foreign key is null or names an existing
parent row; trees are acyclic; and the row counts, scaling and link structure follow the laws of
Sections~\ref{sec:star}--\ref{sec:v10};
\item every computed rule holds exactly on every row;
\item a row that breaks a check rule or a unique attribute key after generation is redrawn from
$q_{G_t}(\cdot\mid c_i;\theta_t)$ until it holds or the round limit is reached. If the acceptance probability
$\pi_i=q_{G_t}(\mathcal R_i\mid c_i)$ is positive, each accepted row has law $q_{G_t}(\cdot\mid c_i,\mathcal R_i)$,
and the probability that a row still breaks a rule after $r$ rounds is $(1-\pi_i)^{r+1}$; such rows are reported.
\end{enumerate}
Parts of the skeleton that read attributes (rule-constrained links, key order by an attribute, attribute strata,
group effects, type 2 history) are functions of the generated values. Their law is then the 1.0 law given those
values.
\end{proposition}
\begin{proof}
(1) Keys and links are produced by the key and link samplers (Props.~\ref{prop:eta}--\ref{prop:cumulative} and the
v0.4--v1.0 sections). They read the fitted relational model and the entity encodings used for segment discovery,
which are built before and independently of any generator (\texttt{TableModel} encodes the columns; the generator is
attached afterwards). The generator's output is written into attribute columns only. The engine checks that it
returns exactly those columns (\texttt{TableModel.sample}), so no key or link column can be overwritten. (2) Computed
columns are evaluated by the engine after generation in the topological order of their inputs (Prop. on computed
columns, v0.7), from the final values, so they hold whatever the inputs. (3) The rejection loop
(\texttt{Synthesizer.\_finish}) redraws the offending rows by calling the same generator with the same context. The
rounds are independent draws from $q(\cdot\mid c_i)$. By the standard rejection-sampling argument (von Neumann
1951), the first accepted draw has the restricted law, and the number of failed rounds is geometric with success
probability $\pi_i$. The last sentence holds because these blocks are evaluated by the 1.0 code on the generated
columns.
\end{proof}

\begin{proposition}[Seed determinism]\label{prop:genseed}
If each generator draws every random number from the generator $\rho$ it is handed (or from generators seeded from
$\rho$), and its fit and sample are deterministic functions of their inputs and the state of $\rho$, then the fitted
model and every synthetic dataset are deterministic functions of the normalised spec, the data and the seed.
\end{proposition}
\begin{proof}
The engine's fitting stream is seeded by the spec seed, and the stream of dataset $m$ by (seed, sample seed, $m$)
(\texttt{posterior.run\_streams}). It hands $\rho$ to every generator call in a fixed order: tables in spec order,
facts in fact order, rejection rounds in row order. A composition of deterministic maps of
$(\text{inputs},\rho\text{-state})$ is deterministic, and so is the state of $\rho$ after each call, by induction on
the calls. The hypothesis is the plug-in's responsibility. The built-in generators satisfy it (numpy only, no global
state), and on a GPU it requires deterministic kernels.
\end{proof}

\begin{proposition}[The copula plug-in is relsynth 1.0]\label{prop:gencopula}
With every table on the \texttt{copula} generator (the default), the fitted model and every synthetic dataset are
identical, byte for byte, to relsynth 1.0's for the same spec, data and seed.
\end{proposition}
\begin{proof}
The copula generator consists of 1.0's \texttt{TableModel} methods \texttt{fit}, \texttt{\_fit\_marginals},
\texttt{redraw} and \texttt{sample}, moved without changing any operation. It reads the same encoded arrays (the
encoding stays in \texttt{TableModel} and runs before the generator is attached). It calls the same random
functions on the same $\rho$ in the same order. The order of the per-dataset redraw is unchanged: marginals, then
copula matrices, then the weights of pooled tail levels. The decoding loop is the same code. The engine adds no
random call when it attaches a generator, and the registry and compute resolution draw no random numbers. Hence
every intermediate value, and the output, is the same. The identity was also checked by running unmodified
\texttt{main} and this version on the eleven fake scenarios (two posterior datasets each), on Chinook (in memory, and
pushed down with two datasets), on Contoso (SQLite and Parquet), on bi/COREP and on AdventureWorks DW, and comparing
the SHA-256 of every output file.
\end{proof}

\begin{proposition}[Posterior predictive datasets]\label{prop:genpost}
Suppose that each generator's \texttt{posterior\_draw} replaces $\theta_t$ by an independent draw from
$p(\theta_t\mid x_t,\hat{\mathcal S})$, its posterior given the table's data and the fitted partition $\hat{\mathcal
S}$. Then dataset $m$ is a draw from the posterior predictive distribution of the joint model (skeleton and
attribute blocks) given the fitted partition, as in 1.0 (Prop.~\ref{prop:ppc}). If
\texttt{posterior\_draw} draws from an approximation $\tilde p(\theta_t)$, the dataset is a draw from the
corresponding approximate predictive. If it does nothing (\texttt{posterior=False}), datasets share $\theta_t$ and
differ only by sampling noise and by the relational parameters.
\end{proposition}
\begin{proof}
For each dataset, the engine draws the relational parameters from their posterior given $\hat{\mathcal S}$ and then
calls \texttt{posterior\_draw} on every fitted table with the dataset's parameter stream
(\texttt{Synthesizer.redraw}). The attribute blocks are conditionally independent of the relational parameters
given the partition and the contexts: the joint posterior factorises over the blocks, because each block's
likelihood involves only its own parameters and the contexts. A joint draw is therefore obtained by drawing each
block independently. Rows are then sampled given the drawn parameters, which is the definition of a posterior
predictive draw (Rubin 1993). The two weaker cases replace the block's exact posterior by $\tilde p$, or by a point
mass at the fitted value.
\end{proof}
\noindent\emph{The built-in generators.}
\begin{itemize}
\item \texttt{copula} is exact in this sense (1.0).
\item \texttt{independent} uses the copula's marginal posteriors and is exact for its own model.
\item \texttt{cart} redraws Bayesian-bootstrap leaf weights given trees fitted once. That is the exact posterior of
the leaf distributions under a Dirichlet process with vanishing concentration given the tree (Rubin 1981, Prop.~\ref{prop:cartbb}),
and an approximation of the posterior over trees. synthpop's ``proper'' synthesis refits the trees on a bootstrap
sample instead.
\end{itemize}

\subsection{Conditioning}
\begin{proposition}[What group-only conditioning keeps and loses]\label{prop:gencond}
Let a fact row $i$ with group $g_i$ have parent attributes $a_i$ and own attributes $x_i$. If the generator
conditions on the group only, then in the synthetic data $x_i\perp a_i\mid g_i$, and
\[\operatorname{Cov}(x_i,a_i)=\operatorname{Cov}_{g}\bigl(\mathbb E[x\mid g],\mathbb E[a\mid g]\bigr),\]
the covariance between the group means. With parent conditioning, $q(x\mid g,a)$ can also carry the within-group
dependence on $a$.
\end{proposition}
\begin{proof}
The engine draws the skeleton (and with it $g_i$ and the parent rows, hence $a_i$) first, then $x_i$ from
$q(\cdot\mid g_i)$, which does not read $a_i$; so $x_i$ is conditionally independent of $a_i$ given $g_i$. The law of
total covariance gives $\operatorname{Cov}(x,a)=\mathbb E_g\operatorname{Cov}(x,a\mid g)+\operatorname{Cov}_g(\mathbb
E[x\mid g],\mathbb E[a\mid g])$, and the first term is zero by conditional independence.
\end{proof}

\begin{proposition}[Independent marginals are a latent class model]\label{prop:genind}
The \texttt{independent} generator has the copula generator's marginal posteriors per group and the identity
correlation. Within a group its columns are independent, and between columns $j,k$ of a table
$\operatorname{Cov}(x_j,x_k)=\operatorname{Cov}_g(\mathbb E[x_j\mid g],\mathbb E[x_k\mid g])$. It is therefore the
latent class model of Goodman (1974), with the latent segments as classes and the hierarchical Dirichlet marginals of
Section~\ref{sec:star} (Bayesian form: Dunson and Xing 2009).
\end{proposition}
\begin{proof}
It fits the marginals with the copula generator's code (\texttt{SegmentMarginals}). It draws each column as
$F^{-1}_{g,j}(U_j)$ with independent uniforms $U_j$. That is the copula sampler with $C=I$, because $\Phi(Z_j)$ with
$Z\sim N(0,I)$ are independent uniforms. Independence within a group follows, and the covariance identity is
Prop.~\ref{prop:gencond} applied to two columns of the same row.
\end{proof}

\subsection{Sequential CART}
\begin{proposition}[Leaf donors with Bayesian-bootstrap weights]\label{prop:cartbb}
Let leaf $\ell$ of the tree of column $j$ hold the observed values $y_1,\dots,y_{n_\ell}$, with
$n_\ell\ge$\,\texttt{min\_leaf}. Draw $w\sim\mathrm{Dirichlet}(1,\dots,1)$ once per dataset and draw a synthetic
value $Y=y_I$ with $P(I=i\mid w)=w_i$. Then:
\begin{enumerate}
\item marginally over $w$, $Y$ has the empirical law of the leaf, $P(Y=v)=\#\{i:y_i=v\}/n_\ell$;
\item given $w$, the probability of a value $v$ with count $c$ in the leaf has mean $c/n_\ell$ and variance
$c(n_\ell-c)/\bigl(n_\ell^2(n_\ell+1)\bigr)$;
\item $w$ is the posterior of the leaf's distribution under a Dirichlet process prior whose concentration tends to 0,
given the $n_\ell$ observations (Rubin 1981).
\end{enumerate}
In mean mode, $w_i=1/n_\ell$ and the dataset uses the empirical leaf law.
\end{proposition}
\begin{proof}
(1) $P(Y=v)=\mathbb E[\sum_{i:y_i=v}w_i]=c\,\mathbb E[w_1]=c/n_\ell$, because the Dirichlet$(1,\dots,1)$ is
exchangeable with mean $1/n_\ell$. (2) $W_v=\sum_{i:y_i=v}w_i\sim\mathrm{Beta}(c,n_\ell-c)$ by the aggregation
property, which has the stated mean and variance. (3) The posterior of $\mathrm{DP}(\alpha,H)$ given the
observations is $\mathrm{DP}(\alpha H+\sum_i\delta_{y_i})$. As $\alpha\to0$ its weights on the atoms $y_i$ tend to
$\mathrm{Dirichlet}(1,\dots,1)$ (Ferguson 1973; Rubin 1981).
\end{proof}
\noindent The tree is fitted once: synthetic rows walk it with their own synthetic predictors (group, coarse group,
parent attributes and the columns before $j$), which is the sequential synthesis of Reiter (2005) and Nowok, Raab and
Dibben (2016). A row whose group the fit never saw (a sparse fact block passed as its coarse group) goes down the
larger branch at a split on the group, and the coarse-group splits then apply to it as they do to the fitted rows.
Without \texttt{smoothing}, every synthetic value of a column is an observed value of that column, as in synthpop's
default; \texttt{smoothing: density} adds Gaussian kernel noise with Silverman's bandwidth to numbers.
\setcounter{proposition}{\value{savedpropgen}}

%=====================================================================
\section{The canonical intermediate representation (v2.0)}\label{sec:v20}
Version 2.0 replaces the relational spec as relsynth's internal description of the data by one modality-neutral
document, the \emph{IR} (\texttt{relsynth/ir}; reference \texttt{docs/ir-reference.md}, design record
\texttt{docs/adr/0001-canonical-ir.md}). An IR document $E$ (an \emph{estate}) lists entities with a shape (record,
document, node, edge, event), attributes (nested objects, arrays, polymorphic variants), keys, references, time,
physical layout, rules and privacy annotations. The engine of this release synthesises the relational part: it runs
the 1.0 spec $\lambda(E)$ that $E$ \emph{lowers} to. A 1.0 spec $s$ is \emph{lifted} to an IR document $\ell(s)$.
Write $\nu$ for the normalisation of a 1.0 spec (\texttt{spec.normalise}: defaults, shorthands, named
references), so that the engine is a function of $\nu(s)$, the data and the seed. The checks are in
\texttt{tests/unit/test\_ir.py}.

\subsection{Converting 1.0 specs}
\begin{proposition}[Loss-free conversion and identical output]\label{prop:irround}
For every 1.0 spec $s$ that the 1.0 key schema accepts, $\nu(\lambda(\ell(s)))=\nu(s)$. Hence, for the same data
and seed, the IR document $\ell(s)$ and the spec $s$ give the same fitted model and the same synthetic datasets.
Moreover $\ell(\lambda(\ell(s)))=\ell(s)$ up to the canonical form $\kappa$ (Definition~\ref{def:canon}), so a lifted
document is a fixed point of lowering and lifting.
\end{proposition}
\begin{definition}[Canonical form]\label{def:canon}
$\kappa(E)$ makes explicit what follows from other fields (an attribute's modelling kind from its logical type,
a reference's target attributes from the parent's primary key, the privacy role of a \texttt{pii} attribute) and
drops what is implied (a logical type that the kind fixes).
\end{definition}
\begin{proof}
$\ell$ and $\lambda$ are defined field by field by a fixed correspondence (module \texttt{ir/v1.py}, table in its
docstring). Every key of the 1.0 key schema has an image: either an IR field holding the same value
(\texttt{columns}$\to$\texttt{attributes}, \texttt{primary\_key}/\texttt{unique}/\texttt{partial\_unique}$\to$
\texttt{keys}, \texttt{scd}$\to$\texttt{time.valid\_time}+\texttt{time.versioning},
\texttt{fact.time\_column}$\to$\texttt{time.event\_time}, \texttt{dimensions}$\to$\texttt{dependencies}, \ldots) or
a value that $\nu$ itself would produce (a default relationship name, kind or parent key; a string column
shorthand; a soft rule written as a derive rule with \texttt{holds}$<1$). $\lambda$ inverts the field images and
emits defaults that $\nu$ fills in the same way. Order is kept: entities, attributes, keys, references and rules
are ordered maps or lists in $E$, and $\nu$ does not reorder them (the engine breaks ties in generation order by
these orders, so order is part of the meaning). Therefore $\nu$ applied to both sides sees the same values at the
same keys in the same order. The engine reads only $\nu(s)$, the data and the seed, so its output is the same.
For the fixed point: $\ell(\lambda(\ell(s)))$ reads back the values $\lambda$ wrote, which are those of $\ell(s)$
except where $\ell$ derives one from another; $\kappa$ normalises exactly those derivations.
\end{proof}
\noindent\emph{Checks.} \texttt{test\_real\_specs\_round\_trip}: the eleven 1.0 specs of the four golden databases
(Chinook, Contoso SQLite and Parquet, AdventureWorks DW, bi/COREP\_LR, including an introspect draft that does not
validate) and the two shipped examples convert both ways with $\nu(\lambda(\ell(s)))=\nu(s)$, as do their normalised
forms; \texttt{test\_fake\_scenario\_specs\_round\_trip} does the same for S1--S11. \texttt{test\_ir\_document\_gives\_identical\_synthetic\_data}
fits S2, S9 and S10 from $s$ and from $\ell(s)$ and compares the synthetic frames cell by cell;
\texttt{test\_cli\_runs\_an\_ir\_file} compares the written files byte by byte. The real databases were run the same way
(report \texttt{roadmap/reports/T02.md}). A lowering $\lambda(E)$ of a document that uses what a 1.0 spec cannot
hold returns the list of those paths, computed as the difference between $\kappa(E)$ and $\kappa(\ell(\lambda(E)))$, so
nothing is dropped in silence (\texttt{test\_lower\_reports\_annotations\_it\_drops}).

\subsection{The privacy unit}
Differential privacy on a database protects a \emph{unit}---a customer, a patient---rather than a row: neighbouring
databases differ in all the records of one unit (add/remove adjacency at the unit level; Ponomareva et al., JAIR 2026;
for foreign keys, Cai, Xiao and Cormode, PrivLava, SIGMOD 2023). Let $U$ be the unit entity. Let $G$ be the graph
whose vertices are the entities and whose arcs are the references $\rho:A\to B$ ($A\neq B$, many-to-one or
one-to-one) that are not hierarchies and not listed in \texttt{privacy.exclude\_paths}. For a path
$P=\rho_1\cdots\rho_k$ from entity $A$ to $U$ in $G$ write $o_P:\mathrm{rows}(A)\to\mathrm{rows}(U)\cup\{\bot\}$ for the
composition of the reference maps (a null or dangling reference gives $\bot$).

\begin{proposition}[Ownership along reference paths]\label{prop:irowner}
Under referential integrity, each $o_P$ is a function: a record of $A$ has at most one owner along each path. If $G$
is acyclic, the set $\mathcal P(A)$ of simple paths from $A$ to $U$ that \texttt{privacy.trace} returns is the set of
all paths, and the records of $A$ that belong to a unit $u$ are $C_u(A)=\bigcup_{P\in\mathcal P(A)}o_P^{-1}(u)$; a
record has at most $|\mathcal P(A)|$ owners. The status of every other entity is \emph{context} (reachable from an
owned entity: the unit's records reference it) or \emph{unrelated}.
\end{proposition}
\begin{proof}
A many-to-one or one-to-one reference maps each record to at most one parent record (the referenced attributes are a
key of the parent; validation requires this), so it is a partial function, and a composition of partial functions is
a partial function. In an acyclic graph every path is simple, and the depth-first search of \texttt{trace} enumerates
all simple paths from $A$ to $U$ (it extends every path by every outgoing arc to a vertex not yet on it). A record
belongs to $u$ when some chain of its references leads to $u$, which is the union over paths. The bound on owners
follows because each path contributes at most one. Context and unrelated entities have no path to $U$ by
construction. If $G$ has a cycle, a non-simple path can lead to a different owner than the simple ones; \texttt{trace}
reports the cycle as a warning (a 1.0 spec cannot have one: its validator requires a topological order).
\end{proof}

\begin{proposition}[Removing a unit keeps referential integrity]\label{prop:irremove}
Let $G$ be acyclic and suppose no reference outside $G$ (a hierarchy or an excluded reference) points into a record of
an owned entity or of $U$. Removing the units of a set $S$ and every record in $\bigcup_{u\in S}C_u(\cdot)$ leaves a
database in which every non-null reference resolves.
\end{proposition}
\begin{proof}
Suppose a kept record $r$ of $A$ references a removed record $t$ of $B$ through $\rho$. By the assumption $\rho$ is an
arc of $G$. As $t$ was removed, $o_P(t)\in S$ for a path $P$ from $B$ to $U$ (or $t$ is a unit in $S$, $P$ empty).
Then $\rho P$ is a path from $A$ to $U$; it is simple because $G$ is acyclic, so $\rho P\in\mathcal P(A)$ and
$o_{\rho P}(r)=o_P(t)\in S$, hence $r$ was removed: a contradiction.
\end{proof}
\noindent When the assumption fails (a customer who referred another, an employee's manager when employees are the
unit) the reference links two units; \texttt{trace} lists it as a cross-unit link and warns that removal leaves it
dangling, so the privacy design (T07) must null or bound it. \emph{Checks.}
\texttt{test\_privacy\_unit\_paths\_of\_the\_real\_databases} states the unit of each golden database and checks its
paths (an invoice line of Chinook belongs to its invoice's customer; a Contoso sale has two paths, directly and
through its order; a COREP reported value has two paths to its reporting entity);
\texttt{test\_contribution\_matches\_the\_joins} recomputes the owners of the S9 shop by joins along every path
(Prop.~\ref{prop:irowner}); \texttt{test\_removing\_a\_unit\_keeps\_referential\_integrity} removes every seventh
customer with its records and checks every reference (Prop.~\ref{prop:irremove});
\texttt{test\_hierarchy\_on\_the\_unit\_is\_a\_cross\_unit\_link} checks the warning.

%=====================================================================
\section{Privacy risk measurement (v2.5)}\label{sec:v25}
% numbered to match checks C600--C606; the shared counter is restored at the end
\newcounter{savedprop}\setcounter{savedprop}{\value{proposition}}
Release 2.5 measures the privacy risk of a synthetic release with attacks (\texttt{relsynth/privacy}). The propositions of this section are numbered from 600 so that they do not collide with those of other work streams. Each is tested by a check in \texttt{tests/unit/test\_privacy.py} (fast) or \texttt{tests/science/test\_privacy\_checks.py} (checks C600--C606, coverage by simulation).

\paragraph{Setting.} A database $D$ holds the rows of a finite set of people $I$. The \emph{privacy unit} is an entity table $U$ with primary key $k$; a table $T$ is \emph{owned} by $U$ when a chain of single-column foreign keys $T=T_0\to T_1\to\dots\to T_m=U$ exists that does not pass through another person table. Every row $r$ of an owned table has an owner $o(r)\in I\cup\{\bot\}$, the key reached by following the chain ($\bot$ when a key on the chain is null). Tables that are not owned are \emph{shared}. For $S\subseteq I$, $D[S]$ keeps the rows of owned tables whose owner is in $S$, the rows of $U$ with key in $S$, and every shared table whole. A \emph{generator} $G$ maps a database and a seed to a synthetic database. The people are split at random into training $A$, control $C$ and reference $R$; the release is $Y=G(D[A])$.

\setcounter{proposition}{599}
\begin{proposition}[Splits by person]\label{prop:pusplit}
(a) For disjoint $S_1,S_2$, $D[S_1]$ and $D[S_2]$ share no row of an owned table, and $D[S_1\cup S_2]$ is their union on owned tables. (b) Every foreign key of $D[S]$ along an ownership chain resolves inside $D[S]$; every foreign key into a shared table resolves because shared tables are whole. Hence $D[S]$ is a valid database for $G$, and any function of the rows of person $i$ is a function of $D[\{i\}]$.
\end{proposition}
\begin{proof}
(a) A row of an owned table has exactly one owner, because the chain is a sequence of single-column foreign keys, each of which selects one parent row; it is kept in $D[S]$ iff its owner is in $S$. (b) Let $r$ be a kept row of $T_j$ with parent $r'$ in $T_{j+1}$ on the chain. Both have the same owner (the owner of $r$ is defined through $r'$), so $r'$ is kept. Foreign keys into shared tables resolve since those tables are unchanged. Rows whose owner is $\bot$ belong to nobody and are kept in no $D[S]$; they are not personal rows of anyone. \texttt{test\_unit\_paths\_and\_split\_keep\_people\_whole} checks that no row is lost or shared and that shared tables stay whole; \texttt{test\_unit\_chain\_stops\_at\_other\_person\_tables} checks the chain rule.
\end{proof}

\setcounter{proposition}{604}
\begin{proposition}[Profiles do not see keys or order]\label{prop:puprofile}
Let $\pi$ be a profile map that sends a database to one row per person (relsynth: the person's attributes, attributes of their parents up to two hops, and for each owned table the number of rows and the mean or the mode, ties broken to the smallest value, of each attribute of the rows and of their parents). Then $\pi(D)$ is invariant, up to the matching relabelling of its index, under any bijective relabelling of key values (applied consistently to primary and foreign keys) and any permutation of the rows of any table.
\end{proposition}
\begin{proof}
No key column enters a profile: keys are only used to join and to group. Joins and group memberships are preserved by a consistent bijection of key values. Counts and means are symmetric functions of the multiset of values in a group, and the mode with ties broken by value order is a function of the multiset as well, so row order does not matter. \texttt{test\_profiles\_ignore\_keys\_and\_row\_order} relabels every key and shuffles every table of a star schema.
\end{proof}
\noindent This is why the attacks are not fooled by synthetic keys being new labels: a key match between real and synthetic data carries no information about a person.

\subsection{Record isolation, linkage and inference}
The three criteria of EDPB Guidelines 02/2026 \cite{edpb2026} (sec.~3.4: no record isolation, no linkage, no inference) are tested by Anonymeter's attacks \cite{giomi2023}. An attack builds guesses $g_1,\dots,g_n$ from $Y$ (and the attacker's knowledge of a target), each judged a success or not on a target set. Let $k_A$ be the number of successes on training people and $k_C$ on control people with the same number of guesses. The \emph{risk} is $\rho=(r_A-r_C)/(1-r_C)$ with Wilson estimates $r_A,r_C$ \cite{wilson1927}, its interval propagated to first order as in Anonymeter.

\setcounter{proposition}{600}
\begin{proposition}[Control exchangeability]\label{prop:puctrl}
Suppose $Y$ is independent of which people are in $A$ given $|A|$ (the generator ignores individuals: it uses only population-level information that does not depend on the split). For record isolation, let $A'\subseteq A$ be a uniform subset with $|A'|=|C|$ ($A'=A$ when the sizes are equal). Then $E[k_{A'}]=E[k_C]$ for record isolation, and $E[k_A]=E[k_C]$ for linkage and inference when targets are drawn uniformly from $A$ and $C$. Hence $E[r_A]-E[r_C]=O(1/n)$ (the Wilson shrinkage only) and the risk is centred at zero up to this term.
\end{proposition}
\begin{proof}
The split is a uniform random partition of $I$, so given $Y$ the sets $A'$ and $C$ are exchangeable: $(Y,A',C)$ and $(Y,C,A')$ have the same law, because $Y$ does not depend on the split and both are uniform subsets of equal size of $I$ disjoint from each other. The queries are functions of $Y$ only, and a query's success is a function of the query and the target set; so the number of successes has the same law on $A'$ and on $C$. For linkage and inference each guess is about one target drawn uniformly from its set, and the success depends on $Y$ and that target's profile; by the same exchangeability the target from $A$ and the target from $C$ have the same law given $Y$. The Wilson estimate $(k+z^2/2)/(n+z^2)$ is a fixed affine map of $k$, so equal means of $k$ give equal means of $r$; the residual is a smooth function of $(r_A,r_C)$ whose value at $r_A=r_C$ is 0, and a first-order expansion gives the $O(1/n)$ bias.
\end{proof}
\noindent Without the subsample, record isolation on a training set larger than the control is biased upwards (a query that isolates in the control may match several training people), which is what Anonymeter corrects with a fitted curve. Check C601 shows that the fitted correction is not exact: with $|A|=2.5|C|$ and a generator that never saw the people the risk was $0.10$--$0.14$ with intervals excluding $0$; the subsample gives $0$ with intervals covering $0$. \texttt{assess} therefore splits training and control in equal numbers (40\%/40\%, 20\% reference), so the subsample is the whole training set and no power is lost.

\subsection{Membership inference}
Three attacks read membership from a release: shadow modelling \cite{stadler2022,houssiau2022} with a logistic classifier on target-specific counting queries and distances, a LiRA-style Gaussian likelihood ratio on the distance from the target to its closest synthetic profile \cite{carlini2022} on the same shadow fits, and DOMIAS \cite{vanbreugel2023}, the ratio of the synthetic density to the population density at the person, on the release itself. Targets are the people with the largest mean cosine distance to their five nearest neighbours \cite{meeus2023}, and as many random people.
\setcounter{proposition}{601}
\begin{proposition}[Shared shadow fits]\label{prop:pushadow}
Let targets $t_1..t_m$ be fixed and, for shadows $s=1..S$, draw $B_s$ uniformly from a population pool $P$ that contains no target, and membership indicators $M_{sj}\sim\mathrm{Bernoulli}(1/2)$ independent of each other and of $B_s$; fit $Y_s=G(D[B_s\cup\{t_j:M_{sj}=1\}])$. For each $j$ the pairs $(Y_s,M_{sj})_{s}$ are i.i.d., and $M_{sj}$ is independent of the background $(B_s,(M_{sl})_{l\ne j})$. So the test-shadow AUC of an attack trained on the training shadows is an unbiased estimate of the AUC of that attack in the game ``target $t_j$ in or out, over a random background of $|B|+\mathrm{Bin}(m-1,1/2)$ people''.
\end{proposition}
\begin{proof}
Independence of the $M_{sj}$ across $j$ and of $B_s$ makes $M_{sj}$ independent of everything else that enters $Y_s$ except through $t_j$'s inclusion. Shadows are independent across $s$ (independent draws and seeds), and identically distributed. The empirical AUC on i.i.d. test pairs is a two-sample U-statistic whose expectation is $P(\text{score}_{in}>\text{score}_{out})+\tfrac12 P(=)$, the AUC of the game, given the attack fitted on the training shadows (which are independent of the test shadows). The background of each game is the pool sample plus the other targets that happen to be in, which is the stated distribution.
\end{proof}
\noindent One set of $S$ fits thus serves all $m$ targets, instead of $mS$ fits \cite{houssiau2022,meeus2023}. The check \texttt{test\_shadow\_attack\_raw\_release\_vs\_population\_generator} and \texttt{benchmarks/privacy/reproduce\_mia.py} find AUC~$1.0$ for a generator that releases its training rows and an interval covering $0.5$ for one that samples the population.

\setcounter{proposition}{603}
\begin{proposition}[Exact intervals for TPR at a fixed FPR]\label{prop:putpr}
Let a threshold $\tau$ be the $\lceil(1-\alpha)N_0\rceil$-th order statistic of $N_0\ge1/\alpha$ calibration non-member scores, and let the test scores be independent of the calibration scores. Conditionally on $\tau$, the number of test members above $\tau$ is $\mathrm{Bin}(P,\mathrm{TPR}(\tau))$ and the number of test non-members above $\tau$ is $\mathrm{Bin}(N,\mathrm{FPR}(\tau))$, so the Clopper--Pearson intervals of both have coverage at least the nominal level for $\mathrm{TPR}(\tau)$ and $\mathrm{FPR}(\tau)$. Moreover $E[\mathrm{FPR}(\tau)]\le\alpha+1/(N_0+1)$ when non-member scores are continuous.
\end{proposition}
\begin{proof}
Given $\tau$ the test scores are i.i.d. within each class and independent of $\tau$, so the counts are binomial with the stated success probabilities; the Clopper--Pearson interval is exact for a binomial proportion \cite{clopper1934}. For the last claim, $\mathrm{FPR}(\tau)=1-F_0(\tau)$, and $F_0(\tau)$ is the $j$-th order statistic of $N_0$ uniforms with $j=\lceil(1-\alpha)N_0\rceil\ge(1-\alpha)N_0$, whose mean is $j/(N_0+1)$. Hence $E[\mathrm{FPR}(\tau)]=1-j/(N_0+1)\le(\alpha N_0+1)/(N_0+1)\le\alpha+1/(N_0+1)$, and the report also prints the realised FPR with its interval. A threshold set on the scores that are then counted would have no such guarantee, which is why the report never does it; with $N_0<1/\alpha$ the entry is marked not identified (C604).
\end{proof}

\setcounter{proposition}{602}
\begin{proposition}[Differential privacy bounds every membership test]\label{prop:pudp}
If the release mechanism is $(\varepsilon,\delta)$-differentially private with respect to adding or removing one privacy unit, then any membership test of one person satisfies $\mathrm{TPR}\le\min\{e^{\varepsilon}\mathrm{FPR}+\delta,\;1-e^{-\varepsilon}(1-\delta-\mathrm{FPR})\}$.
\end{proposition}
\begin{proof}
This is the hypothesis-testing characterisation of differential privacy \cite[Thm.~2.1]{kairouz2015}: for neighbouring databases and any rejection region, $\mathrm{FPR}+e^{\varepsilon}\mathrm{FNR}\ge1-\delta$ and $e^{\varepsilon}\mathrm{FPR}+\mathrm{FNR}\ge1-\delta$. Substituting $\mathrm{FNR}=1-\mathrm{TPR}$ gives both bounds. The neighbouring relation must be the one of the privacy unit (all rows of one person, Prop.~\ref{prop:pusplit}); a guarantee per row does not bound a test of a person with several rows.
\end{proof}
\noindent The report prints this bound next to the empirical TPR when the release carries a DP guarantee (filled by the DP engine). C603 checks it on a person count released with Laplace noise at $\varepsilon\in\{0.5,1,2\}$: exactly on a grid of thresholds (where the bound is attained in the tail), and on simulated releases with thresholds set on independent draws.

\subsection{Inference beyond the population}
\setcounter{proposition}{605}
\begin{proposition}[Specific leakage is centred at zero for a generator that ignores individuals]\label{prop:puattr}
Let an attacker model $f_Y$ be trained on $Y$ and a baseline $f_R$ on the reference people, both predicting a secret $s$ from known attributes; let $c^Y_i,c^R_i$ indicate correct guesses on person $i$. With $\Delta_A=\overline{c^Y-c^R}$ over training people and $\Delta_C$ over control people, $E[\Delta_A-\Delta_C]=0$ under the condition of Prop.~\ref{prop:puctrl}, and the normal interval built from the per-person paired differences of each group has asymptotic coverage at the nominal level.
\end{proposition}
\begin{proof}
$f_Y$ depends on $Y$ and $f_R$ on $R$; neither depends on which of the remaining people are in $A$ or $C$ under the condition. Training and control people are exchangeable given $(Y,R)$, so the expected per-person differences are equal in the two groups. Within a group, the paired differences are i.i.d. given $(Y,R)$ with finite variance (they take values in $\{-1,0,1\}$), so the central limit theorem gives the coverage of each group mean; the two groups are independent given $(Y,R)$, so the variance of the difference is the sum. EDPB Guidelines 02/2026 (para.~67) call an inference meaningful only when it is not available from population data: $\Delta_C$ is what $Y$ teaches about everyone, $\Delta_A-\Delta_C$ what it reveals about the people in the training data.
\end{proof}
\setcounter{proposition}{\value{savedprop}}

%=====================================================================
\section{Governance: privacy budget ledger and tamper evidence (v2.5)}\label{sec:v25gov}
% numbered to match checks C800--C803; the shared counter is restored at the end
\newcounter{savedpropgov}\setcounter{savedpropgov}{\value{proposition}}
Release 2.5 records every run in a signed, append-only registry and charges every differentially private release to the privacy budget of the people it is about (design: \texttt{docs/governance/design.md}). The propositions of this section are numbered from 800 so that they do not collide with those of other work streams. Each is tested by a check in \texttt{tests/science/test\_governance\_checks.py} (C800--C803). In the design phase the checks test small reference implementations; the build phase keeps the checks and points them at \texttt{relsynth.governance}.

\paragraph{Setting.} A \emph{privacy domain} $\mathcal{D}$ is a population of privacy units (people) together with every database that holds their data, under one adjacency relation: two collections of databases are neighbours when they differ in all the data of one unit (add or remove), across all databases of the domain at once (Ponomareva et al.\ \cite{ponomareva2023}; NIST SP 800-226 \cite{nist800226}, sec.~2.4). A \emph{charge} is the privacy guarantee of one mechanism run on data of $\mathcal{D}$, as reported by the DP engine (T07) with respect to that adjacency. Charges are of three kinds: pure $\varepsilon$-DP, $\rho$-zCDP, and approximate $(\varepsilon,\delta)$-DP. The \emph{ledger} of $\mathcal{D}$ is the sequence of charges in the order the mechanisms ran; each mechanism may be chosen adaptively from the outputs of earlier ones.

\setcounter{proposition}{799}
\begin{proposition}[The ledger bounds the privacy loss of the domain]\label{prop:ledger}
Map each charge to a pair $(\rho_i,\delta_i)$: pure $\varepsilon$ to $(\varepsilon^2/2,0)$, $\rho$-zCDP to $(\rho,0)$, and $(\varepsilon,\delta)$ to $(\varepsilon^2/2,\delta)$. Let $\rho=\sum_i\rho_i$ and $\delta_0=\sum_i\delta_i$.
(a) The adaptive composition of all mechanisms in the ledger is $\delta_0$-approximately $\rho$-zCDP.
(b) For every $\delta>\delta_0$ it is $(\varepsilon_\rho(\delta-\delta_0),\delta)$-DP, where $\varepsilon_\rho(\delta')$ is the smallest $\varepsilon\ge0$ with
\[\inf_{\alpha>1}\frac{\exp\big((\alpha-1)(\alpha\rho-\varepsilon)\big)}{\alpha-1}\Big(1-\frac1\alpha\Big)^{\alpha}\le\delta',\]
and $\varepsilon_\rho(\delta')\le\rho+2\sqrt{\rho\log(1/\delta')}$.
(c) If every charge is pure or approximate, it is also $(\sum_i\varepsilon_i,\delta_0)$-DP; the ledger states the smaller of the two $\varepsilon$ at the requested $\delta$.
\end{proposition}
\begin{proof}
A pure $\varepsilon$-DP mechanism is $\varepsilon^2/2$-zCDP (Bun and Steinke \cite{bun2016}, Prop.~1.4). An $(\varepsilon,\delta)$-DP mechanism is $\delta$-approximately $\varepsilon^2/2$-zCDP: its output law on neighbours can be written as a mixture with weight $1-\delta$ on a pair of laws that satisfy $\varepsilon$-DP, hence $\varepsilon^2/2$-zCDP, and weight $\delta$ on an arbitrary event (\cite{bun2016}, approximate zCDP). Approximate zCDP composes adaptively with $\rho$ and $\delta$ adding (\cite{bun2016}, composition of approximate zCDP), which is (a). For (b), conditioning on the complement of the $\delta_0$-event leaves a $\rho$-zCDP mechanism, whose $(\varepsilon,\delta')$ profile is bounded by the stated infimum (Canonne, Kamath and Steinke \cite{cks2020}, Prop.~12); adding the $\delta_0$ mass gives total $\delta=\delta'+\delta_0$. Taking $\alpha=1+\sqrt{\log(1/\delta')/\rho}$ in the R\'enyi conversion (Mironov \cite{mironov2017}, Prop.~3) gives the classic bound $\rho+2\sqrt{\rho\log(1/\delta')}$ (\cite{bun2016}, Prop.~1.3), and the stated infimum is never larger (\cite{cks2020}). (c) is basic adaptive composition of approximate DP (Dwork and Roth \cite{dworkroth2014}, Thm.~3.16 and its adaptive form). Both statements hold, so the minimum $\varepsilon$ at the same $\delta$ holds too.
\end{proof}
\noindent Check C800 composes up to 20 Gaussian mechanisms and up to 50 randomised responses, whose exact privacy profiles are known (Balle and Wang \cite{ballewang2018}; Kairouz, Oh and Viswanath \cite{kairouz2015}). The ledger's $\varepsilon$ is never below the exact one and never above the classic conversion; it is at most 1.6 times the exact $\varepsilon$ for Gaussians and never above basic composition for pure charges. For a sense of scale at $\delta=10^{-8}$: $\varepsilon=1$ corresponds to $\rho=0.017$, $\varepsilon=4$ to $\rho=0.23$ and $\varepsilon=8$ to $\rho=0.82$.

The privacy domain, not the source table, is the key of the ledger. Two sources that hold data of the same people (a warehouse and its operational source, or two extracts) are one domain: releases from both compose. The design therefore asks the organisation to name its domains in the policy and maps each source to one.

\setcounter{proposition}{800}
\begin{proposition}[A per-period budget is a per-period guarantee]\label{prop:period}
Suppose a unit's data is used in releases of $P$ periods, each charged $\rho_p$. Then the guarantee for that unit across the periods is $\sum_p\rho_p$-zCDP and no better in general: for Gaussian mechanisms the exact privacy profile is that of a Gaussian with $\mu=\sqrt{2\sum_p\rho_p}$. A cap of $\rho^\ast$ that is reset each period therefore bounds the loss of a unit only if its data enter one period alone (disjoint data, parallel composition).
\end{proposition}
\begin{proof}
The releases of all periods are mechanisms run on data of the same domain, so Prop.~\ref{prop:ledger} applies to their concatenation. For Gaussian mechanisms with sensitivities $\Delta_p$ and scales $\sigma_p$ the composition of the privacy loss random variables is Gaussian with $\mu^2=\sum_p\Delta_p^2/\sigma_p^2=2\sum_p\rho_p$ (Dong, Roth and Su \cite{dong2022}, Cor.~3.3), which is attained by the worst-case neighbours. If a unit's data enter only period $p$, the other periods' mechanisms do not depend on it and the bound is $\rho_p$ (parallel composition, \cite{dworkroth2014}).
\end{proof}
\noindent This is NIST's warning about user-day privacy (\cite{nist800226}, sec.~2.4.2: $\varepsilon=1$ per day can give $\varepsilon=365$ over a year). Check C801: with $\rho=0.05$ per month for 12 months at $\delta=10^{-8}$, the per-period statement is $\varepsilon=1.76$ while the exact loss over the year is $\varepsilon=6.39$; the lifetime ledger states $6.74$. The design therefore makes the lifetime cap mandatory and the period cap an optional throttle.

\setcounter{proposition}{801}
\begin{proposition}[Reserve before reading]\label{prop:reserve}
Let every run (i) append a reservation $r$ to the ledger under a lock, refused when the charged total plus $r$ would exceed the cap; (ii) append a \emph{touch} record before its first read of real data; (iii) on success, append a settlement $s\le r$, the accountant's figure for the mechanisms it ran; and (iv) receive a refund only if no touch record exists. Count a run as charged $r$ from its reservation until a settlement or refund replaces it. Then at every time, for any interleaving of runs and any crashes, the charged total is at least the true privacy loss of the domain and at most the cap.
\end{proposition}
\begin{proof}
Consider any run. If it never touched data, its true loss is $0$ and its charge is $\ge0$. If it touched data and crashed, its charge stays $r$ and the mechanisms it ran are a prefix of a plan whose total is at most $r$, because T07 plans the run's budget before reading and refuses mechanisms that exceed it; so its true loss is $\le r$. If it settled, its loss is the accountant's $s$. Summing over runs, and using that adding a charge in zCDP terms is monotone (Prop.~\ref{prop:ledger}), the charged total bounds the true loss. The cap holds because each reservation is checked against the charged total under the same lock that appends it, and settlements and refunds never increase the total.
\end{proof}
\noindent Check C802 runs 40 concurrent runs in 20 trials with random give-ups before reading, crashes mid-run and settlements; the charged total was never below the true loss nor above the cap, and refunding a run that touched data is refused.

\setcounter{proposition}{802}
\begin{proposition}[Tamper evidence]\label{prop:chain}
Let the registry log be entries $e_0,\dots,e_n$ with $h_{-1}=0^{256}$, $h_i=\mathrm{SHA256}(h_{i-1}\,\|\,\mathrm{JCS}(e_i))$, and let a checkpoint $(n,h_n)$ be signed with Ed25519 by a key the verifier trusts. If an adversary without the key produces a log $e'_0,\dots,e'_{m}$ that verifies against the checkpoint, then $m=n$ and $e'_i=e_i$ for all $i$, unless it has found a SHA-256 collision or forged an Ed25519 signature.
\end{proposition}
\begin{proof}
Verification recomputes the chain and requires $m=n$ and $h'_n=h_n$. We show by downward induction on $j$ that $h'_j=h_j$ implies $e'_j=e_j$ and $h'_{j-1}=h_{j-1}$, unless a collision is found. The inputs of the two hashes are $h'_{j-1}\|\mathrm{JCS}(e'_j)$ and $h_{j-1}\|\mathrm{JCS}(e_j)$. If they differ, they are a SHA-256 collision. If they are equal, the 32-byte prefixes are equal, so $h'_{j-1}=h_{j-1}$, and the rest is equal, so $e'_j=e_j$ because JCS is injective on JSON values (RFC 8785 \cite{rfc8785}). Starting from $j=n$ this gives $e'_j=e_j$ for every $j$. A checkpoint with another $(n,h)$ needs an Ed25519 signature under the trusted key, which the adversary cannot produce (EUF-CMA security of Ed25519, RFC 8032 \cite{rfc8032}).
\end{proof}
\noindent Without a signed checkpoint, an operator who controls the store can rewrite the whole chain; the checkpoint, kept outside the store (in the release package, with the approver, or in a transparency log), is what makes the log tamper-evident to a third party. Check C803 applies 400 random edits, deletions, insertions, swaps and truncations to a log of 200 entries; all 400 are detected, and a forged or re-signed checkpoint is rejected. The check uses Python's JSON with sorted keys and no floats in hashed fields, which equals JCS on that subset; the build uses an RFC 8785 implementation.

\setcounter{proposition}{\value{savedpropgov}}

%=====================================================================
\section{Differential privacy mode (v2.0)}\label{sec:v25dp}
% numbered to match checks C900--C906; the shared counter is restored at the end
\newcounter{savedpropdp}\setcounter{savedpropdp}{\value{proposition}}
Release 2.0 adds a differentially private mode (\texttt{relsynth run --dp}; guide \texttt{docs/dp.md}). It works in two stages. Stage~1 reads the real tables once and releases noisy integer statistics: ranges and category levels, one-dimensional histograms, histograms of rows per parent row, and sums of outer products of normal scores. Stage~2 decodes them into a \emph{seed database} and fits the ordinary engine to it. The propositions of this section are numbered from 900. Each is tested by a check in \texttt{tests/science/test\_dp\_checks.py} (C900--C906) on the shipped code (\texttt{relsynth.dp}).

\paragraph{Setting.} The privacy unit is a table $U$ of the spec with one row per person. A table is \emph{owned} by $U$ when a chain of foreign keys leads from it to $U$; the plan computes the chains from the spec before reading data, and refuses a table with two chains to $U$, a self-reference on an owned table, and a foreign key from an owned table to another owned table off its chain. Every other table must be declared public; it is released as given and a foreign key to it is an ordinary categorical coordinate whose levels are the public keys. Two databases are neighbours when one is obtained from the other by removing one row of $U$ and every row of the owned tables whose chain ends in it (add/remove adjacency, user level: NIST SP 800-226 \cite{nist800226}, sec.~2.4). Every statistic is a vector of integers, perturbed by the discrete Gaussian $\mathcal{N}_\mathbb{Z}(0,\sigma^2)$ on each coordinate, drawn by OpenDP from the operating system's random source.

\setcounter{proposition}{899}
\begin{proposition}[Discrete Gaussian]\label{prop:dgauss}
Let $f$ map databases to $\mathbb{Z}^d$ with $\|f(x)-f(x')\|_2\le\Delta$ for neighbours. Then $M(x)=f(x)+Z$, with $Z$ having independent $\mathcal{N}_\mathbb{Z}(0,\sigma^2)$ coordinates, satisfies $\rho$-zCDP with $\rho=\Delta^2/(2\sigma^2)$. Each coordinate of $Z$ has variance below $\sigma^2$.
\end{proposition}
\begin{proof}
This is the multivariate discrete Gaussian mechanism of Canonne, Kamath and Steinke \cite{cks2020}. The R\'enyi divergence of order $\alpha$ between $\mathcal{N}_\mathbb{Z}(\mu,\sigma^2)$ and $\mathcal{N}_\mathbb{Z}(\mu',\sigma^2)$ for $\mu,\mu'\in\mathbb{Z}^d$ is at most $\alpha\|\mu-\mu'\|_2^2/(2\sigma^2)$, and the variance bound is also theirs. Integer statistics with integer noise avoid the floating-point attacks on continuous noise (Mironov \cite{mironov2012}).
\end{proof}
\noindent relsynth sets $\sigma=\lceil\Delta\rceil/\sqrt{2\rho}$ and charges the figure that OpenDP's privacy map returns for $\lceil\Delta\rceil$. Check C900 computes the R\'enyi divergence of orders $1.5$ to $16$ exactly on the support for three $(\Delta,\rho)$ pairs. It also checks that the accountant charges $\rho$ and refuses a mechanism that would exceed the run's plan.

\setcounter{proposition}{900}
\begin{proposition}[Contribution bounding is stable]\label{prop:bound}
Order the owned tables along their chains from $U$. For a child table $t$ with parent $p$ on its chain, keep at most $D_t$ rows of $t$ per row of $p$: the $D_t$ with the smallest keyed hash $h_k(\text{row})$, where $k$ is a secret key drawn once per run. Let $c_U=1$ and $c_t=c_p D_t$. Then for neighbours $x,x'$ that differ in unit $u$, the truncated databases $T(x),T(x')$ differ only in rows whose chain ends in $u$, and in each table $t$ there are at most $c_t$ such rows.
\end{proposition}
\begin{proof}
By induction along the chains. $U$ has one row for $u$. Suppose that in the kept rows of $p$, the rows of units other than $u$ are the same in $T(x)$ and $T(x')$ and that $u$ has at most $c_p$. A row of $t$ belongs to the unit of its parent row. The rows of $t$ that are kept for a parent row depend only on the rows of $t$ with that parent and on $k$, because each row's priority $h_k(\text{row})$ depends on nothing else. So the kept rows of other units are unchanged, and $u$ keeps at most $D_t$ rows for each of its at most $c_p$ parent rows. The chains are unique, so no row is reached twice. Truncation by data-independent priority is the stable transformation of Wilson et al.\ \cite{wilson2020}. A random priority has the same law for a fixed key, and the key is never released.
\end{proof}
\noindent The caps $D_t$ are chosen after the histogram of rows per parent row has been released (Prop.~\ref{prop:degree}), at a public quantile of it. Choosing them is therefore post-processing, and the sensitivities below hold given the chosen caps (adaptive composition). Check C901 removes the busiest customer of a star schema. The kept sales of every other customer are unchanged, and no customer keeps more than $D$.

\setcounter{proposition}{901}
\begin{proposition}[Sensitivity of the released statistics]\label{prop:sens}
After the truncation of Prop.~\ref{prop:bound}, for a table $t$ with $c=c_t$:
(a) a vector of $q$ histograms of single columns over cells fixed before the read, with one null cell each, plus the row count, has $L_2$ sensitivity at most $c\sqrt{q+1}$;
(b) let $z(r)\in\mathbb{R}^m$ be the normal scores of row $r$ under released marginals (attributes, truncated degrees and, for a child table, the parent row's attribute scores). Clip $z$ to norm $R=\kappa\sqrt m$ and let $v(r)$ be the upper triangle with diagonal of $\mathrm{round}(S\,zz^\top)$ for the fixed scale $S=1024$. Then $\sum_r v(r)$ has $L_2$ sensitivity at most $c\,(SR^2+\tfrac12\sqrt{m(m+1)/2})$.
\end{proposition}
\begin{proof}
(a) Each of the at most $c$ rows of $u$ adds one to one cell of each histogram and one to the count, so each of the $q+1$ blocks changes by at most $c$ in $L_1$, hence in $L_2$, and the blocks add in squares. (b) For one row, the upper triangle with diagonal of $zz^\top$ has norm at most $\|zz^\top\|_F=\|z\|_2^2\le R^2$. Rounding changes each of the $m(m+1)/2$ coordinates by at most $1/2$. A row's score depends only on that row, its parent row and released marginals, so rows of other units do not change. The triangle inequality over the at most $c$ rows of $u$ gives the bound.
\end{proof}
\noindent The normalised sum is a noisy second-moment matrix of the scores. relsynth projects it to the nearest correlation matrix, which is post-processing, and uses it as the Gaussian copula of the seed database. Check C902 is the central check. With public ranges, levels and caps fixed, so that the shapes do not depend on the data, it removes thirteen customers one at a time, including the busiest one. Each time it replays the noisy outputs of the full run into the neighbour's run and compares the exact statistics of every mechanism. Every difference is within its charged sensitivity, and for the marginals of both tables it comes within a factor of three of it.

\setcounter{proposition}{902}
\begin{proposition}[Category levels by partition selection]\label{prop:levels}
For a categorical column of $t$ that is not public, release the noisy count of each level present in the data and keep the levels whose noisy count exceeds $\tau=c+\tau_0$, where $\tau_0$ is the smallest integer with $k\,\Pr[Z>\tau_0]\le\delta_c$ for $Z\sim\mathcal{N}_\mathbb{Z}(0,\sigma^2)$ and $k=c$. This is $\delta_c$-approximately $c^2/(2\sigma^2)$-zCDP.
\end{proposition}
\begin{proof}
On the levels present in both neighbours the counts differ by at most $c$ in $L_1$, so the released vector restricted to them satisfies Prop.~\ref{prop:dgauss}. Unit $u$ alone can hold at most $c$ levels, each with count at most $c$. Each such level is released with probability $\Pr[\text{count}+Z>\tau]\le\Pr[Z>\tau_0]$, so by the union bound with probability at most $\delta_c$ in all. Outside that event the output law is that of the restricted mechanism, which is the definition of approximate zCDP (Bun and Steinke \cite{bun2016}). This is Gaussian partition selection (Desfontaines, Voss, Gipson and Mandayam \cite{desfontaines2022}).
\end{proof}
\noindent relsynth sums the tail of the discrete Gaussian exactly. Check C903 found that the continuous approximation $\tau_0=\sigma\,\Phi^{-1}(1-\delta_c/k)$ is not conservative: for $c=40$, $\rho=0.1$ it gives $1.03\,\delta_c$. The shipped threshold is the exact one. Numeric ranges are found the same way: the noisy counts on a fixed, data-free logarithmic grid, keeping the outermost cells above $3\sigma$ (pure zCDP, by (a) of Prop.~\ref{prop:sens}). A declared public range or level list replaces discovery and its budget share.

\setcounter{proposition}{903}
\begin{proposition}[Rows per parent row]\label{prop:degree}
For a parent table $t$ and a child $s$, the histogram over the rows of $t$ of their number of $s$-rows, on cells fixed before the read, has $L_2$ sensitivity at most $c_t$.
\end{proposition}
\begin{proof}
Each of the at most $c_t$ rows of $u$ in $t$ falls in one cell, and the rows of other units do not change their counts (their children are theirs, Prop.~\ref{prop:bound}).
\end{proof}

\setcounter{proposition}{904}
\begin{proposition}[Post-processing]\label{prop:post}
The seed database, the model of stage~2 fitted to it, every synthetic database sampled from that model and every report computed from the synthetic data alone satisfy the guarantee of stage~1.
\end{proposition}
\begin{proof}
Each is a randomised function of the released statistics, the spec and the public tables, and does not read the real data. Post-processing preserves zCDP and approximate DP (\cite{bun2016}; Dwork and Roth \cite{dworkroth2014}, Prop.~2.1). The spec is an input. Values in it that were drawn from the data (levels, ranges, rules written by \texttt{introspect}) are outside the guarantee unless the spec is declared public (\texttt{spec\_is\_public}), and the DP statement says so.
\end{proof}
\noindent Check C904 runs the whole pipeline with nothing declared public and plants a region held by one customer. The region does not appear in the synthetic data. The run spends at most its plan, and the stated $\varepsilon$ is at most the requested one. Check C905 confirms that every link of the synthetic database is valid and that the public table is used as given.

\setcounter{proposition}{905}
\begin{proposition}[The run's guarantee]\label{prop:run}
Let the run's plan be $(\rho,\delta)$ with $\delta\le10^{-5}$, and let a share $s$ of $\delta$ be reserved for level selection. The accountant charges every mechanism of stage~1 and refuses any that would exceed $\rho$ or $s\delta$. The whole run is then $s\delta$-approximately $\rho$-zCDP and $(\varepsilon_\rho((1-s)\delta),\delta)$-DP for one unit and all their rows. When the run is governed, its reservation is $(\rho,s\delta)$ in the ledger of its domain, at the domain's $\delta$, and Props.~\ref{prop:ledger}--\ref{prop:reserve} apply.
\end{proposition}
\begin{proof}
Stage~1 is an adaptive composition of the mechanisms of Props.~\ref{prop:dgauss}, \ref{prop:levels} and \ref{prop:degree}, each with the sensitivity of Props.~\ref{prop:sens} and \ref{prop:degree} given the earlier outputs. Their $\rho$ and $\delta$ add (\cite{bun2016}), and the accountant keeps the sums within the plan. Stage~2 is post-processing (Prop.~\ref{prop:post}). The conversion to $(\varepsilon,\delta)$ is (b) of Prop.~\ref{prop:ledger}. A run that is refused, because it names no budget, has a $\delta$ above $10^{-5}$, has a unit that is not a table or has an ownership the plan cannot trace, is refused before it reads data and before it reserves budget.
\end{proof}
\noindent Check C906 tests these refusals. relsynth ships no default $\varepsilon$: a DP run without \texttt{epsilon} or \texttt{rho} is refused.

\paragraph{What the guarantee does not cover.} The guarantee does not cover tables that hold no data of the unit (they are released as given), spec values not declared public, and releases of the same people made without DP. An empirical attack audit (\texttt{relsynth privacy}, which audits the DP pipeline itself) is an average-case check. Its true-positive rate is compared with the bound $\mathrm{TPR}\le e^\varepsilon\mathrm{FPR}+\delta$ (Kairouz, Oh and Viswanath \cite{kairouz2015}), but it does not bound outliers. Rare levels are suppressed and large contributions are truncated, which biases the data towards typical units, and the DP statement lists both.

\setcounter{proposition}{\value{savedpropdp}}

%=====================================================================
\section{Degree laws for the keys of link tables}\label{sec:linkdeg}
A link (many-to-many) table is a pair fact: each row joins an anchor row $i$ (a movie, the first atom of a bond) to a partner row $j$ (an actor, the second atom). The block model fixes the number $N_{kl}$ of rows joining anchor segment $k$ to partner segment $l$ (Props.~\ref{prop:eta}, \ref{prop:scale}). Within a segment, 1.0 spreads the rows uniformly (occupancy facts: every pair of a block present with the same probability) or by gamma frailties (Poisson facts), and only the anchor of a Poisson fact can have a degree law of its own (Prop.~\ref{prop:countbf}). The SyntheRela benchmark showed what this costs: the second key of \texttt{bond} (every atom the second atom of about one bond) and the keys of \texttt{movies2actors} and \texttt{movies2directors} (heavy-tailed degrees) were the weakest relationships, with cardinality similarity 0.68--0.80. With \texttt{synthesis.facts.link\_degree: auto} (per fact \texttt{fact.link\_degree}) each key of a pair fact may get a degree law; the default \texttt{segment} keeps the 1.0 sampler and its output unchanged. Code: \texttt{Synthesizer.\_fit\_link\_degree}, \texttt{\_link\_degrees}, \texttt{\_cap\_degrees}, \texttt{\_repair\_pairs}; \texttt{counts.nb\_size\_pooled}.

\paragraph{Choice of the law.} For one side, let $t_i$ be the source degree of row $i$ (rows of the fact for a Poisson fact, distinct pairs for an occupancy fact), grouped by the side's segments. The count-model evidence of Prop.~\ref{prop:countbf} is computed on these groups. The side takes the DP law when the DP beats Poisson and NB by the Bayes factor \texttt{facts.count\_model\_bf}; otherwise the NB law when NB beats Poisson and the side has no frailty of its own (occupancy facts, or frailty off); otherwise it keeps the segment law, since Poisson degrees within a segment are what the block model already gives. Under \texttt{auto} a side also keeps the segment law when its links are concentrated on a few rows: with $L_k$ links in segment $k$ and Kish's effective number of rows $n^{\mathrm{eff}}_k=(\sum_i t_i)^2/\sum_i t_i^2$ over the segment's rows, the link-weighted mean $\sum_k L_k n^{\mathrm{eff}}_k/\sum_k L_k$ must reach \texttt{facts.link\_degree\_min\_rows} (default 50). Any exchangeable law hands a heavy segment's links to a few randomly chosen synthetic rows, so an attribute joined through the key (the style of the products sold) becomes a mixture over about $n^{\mathrm{eff}}$ rows, whose category shares have standard deviation $\sqrt{p(1-p)/n^{\mathrm{eff}}}$ (0.07 at $n^{\mathrm{eff}}=50$, $p=\frac12$), while the segment law spreads them evenly. AdventureWorks showed it: products in internet sales ($n^{\mathrm{eff}}=8.5$, NB size 0.15) raised the product-style join TVD from 0.05--0.13 to 0.11--0.33. The SyntheRela keys have 88 to 23{,}642 effective rows, Chinook's track keys 83 and 152, Contoso's product key 67. A side keeps the segment law when its degrees are already modelled (DP counts per anchor, mandatory participation), when its rows are copied or fixed (each row keeps its own observed share, Prop.~\ref{prop:shares}), for virtual dimensions, and under eligibility groups (structural zeros between groups).

\begin{proposition}[Negative binomial given its total]\label{prop:nbdirmult}
Let $x_1,\dots,x_M$ be i.i.d.\ $\mathrm{NB}(r,p)$. Given $\sum_i x_i=R$, the vector $x$ is Dirichlet-multinomial $(R;\,r,\dots,r)$, whatever $p$. Equivalently: draw $w_i\sim\mathrm{Gamma}(r)$ and $x\mid w\sim\mathrm{Multinomial}(R,\,w/\sum w)$.
\end{proposition}
\begin{proof}
$P(x)=\prod_i \frac{\Gamma(x_i+r)}{\Gamma(r)\,x_i!}(1-p)^r p^{x_i}$, and on $\{\sum x_i=R\}$ the factor $(1-p)^{Mr}p^R$ is constant, so $P(x\mid R)\propto\prod_i \Gamma(x_i+r)/x_i!$, which is the Dirichlet-multinomial kernel $\frac{R!\,\Gamma(Mr)}{\Gamma(R+Mr)}\prod_i\frac{\Gamma(x_i+r)}{x_i!\,\Gamma(r)}$. The gamma-multinomial representation is the usual one: normalised gammas are $\mathrm{Dirichlet}(r,\dots,r)$.
\end{proof}
The size $r$ is shared by the side's segments (each with its own mean, so only $r$ is common): the maximum-likelihood value over all segments, so that small segments borrow strength (the pooling pattern of Section~\ref{sec:v09}). Like the frailty's method-of-moments shape it is a plug-in (empirical Bayes) value. When every source row of a segment has a link, the law is shifted by one: $x_i=1+y_i$ with $y\sim\mathrm{DirMult}(R-M;\,r,\dots,r)$.

\begin{proposition}[Block configuration sampler]\label{prop:linkdeg}
Fix the block of every row as drawn by the block model, so $N_{kl}$ are given. For a side with a degree law, segment $k$, its $M_k$ eligible rows and its $R_k=\sum_l N_{kl}$ rows of the fact, let $d\sim F_k^{\otimes M_k}\mid\sum d=R_k$ (DP: the posterior predictive of Prop.~\ref{prop:dp} as in Prop.~\ref{prop:condcounts}; NB: Prop.~\ref{prop:nbdirmult}), and deal the multiset of stubs $\{i^{(d_i)}\}$ to the $R_k$ rows in a uniformly random order; the other side is dealt independently. Then
\begin{enumerate}
\item[(a)] the block totals $N_{kl}$, hence the expected table size, the scale rule (Prop.~\ref{prop:scale}) and the attribute groups, are those of the block model;
\item[(b)] the degree vector of segment $k$ has the law $F_k^{\otimes M_k}$ conditioned on the segment total, so each row's degree follows the source's degree law of its segment up to that conditioning;
\item[(c)] given both degree sequences, row $i$'s degree is split over the blocks $(k,l)$ by the multivariate hypergeometric law, and within a block the anchor stubs are matched to the partner stubs by a uniformly random bijection: the degree-corrected (configuration) block model of Karrer and Newman (2011), conditioned on its degrees.
\end{enumerate}
\end{proposition}
\begin{proof}
(a) Only the row indices within the given blocks change. (b) is Prop.~\ref{prop:condcounts} for the DP law (the pairwise Gibbs sampler leaves the conditional law invariant) and Prop.~\ref{prop:nbdirmult} for the NB law. (c) A uniformly random order of the $R_k$ stubs assigned to rows already sorted into blocks of sizes $N_{kl}$ gives each block a uniformly random sub-multiset of size $N_{kl}$ (multivariate hypergeometric counts per row) in uniformly random order; dealing the other side independently pairs the two orders, which is a uniform bijection between the block's anchor and partner stubs.
\end{proof}

\paragraph{Unique pairs and self-links.} An occupancy fact has at most one row per pair (multiplicities are added after, Prop.~\ref{prop:mult}), and a link table to a single parent whose source has no self-links (\texttt{bond}, \texttt{cites}) should not get any. The configuration sampler can create both. Pairs of a complete block ($p_{kl}=1$ in a draw: every product on every date in AdventureWorks' inventory) are structural; they, and every segment that has one, keep the occupancy draw, so a complete grid stays complete. Degrees of an occupancy side are first capped at the number of partner rows the segment can reach (the excess is moved to rows below the cap, keeping the total). Then a bad row $(a,b)$ exchanges ends with a good row $(a',b')$, giving $(a,b')$ and $(a',b)$: partners when $j$ is in the same partner segment, anchors when it is in the same anchor segment, so every row keeps a block that matches its segments and every degree and $N_{kl}$ is unchanged. A swap is accepted only when both new pairs are new and not self-links, so each removes one bad row and creates none. This is the edge-switching repair used for configuration models (Milo et al.\ 2003; Fosdick et al.\ 2018); it makes the graph simple without changing (a)--(c)'s degree sequences, but it is not an exact draw from the configuration model conditioned on simplicity, and we do not claim it is. A row that finds no swap takes a free partner of its block's partner segment (one partner degree moves by one); a row with no free partner is dropped. Both are counted in the sample log (\texttt{link\_degree\_repair}). On the IMDB ratings (1.0~M rows, $3{,}832\times6{,}039$ possible pairs, degrees up to 3{,}427) 0.19\% of the rows are dropped; on the other databases none.

\paragraph{Checks.} \texttt{tests/unit/core/test\_link\_degree.py}: the pooled size recovers $r$ on simulated NB data; i.i.d.\ NB given its total and the gamma-multinomial draw both match the beta-binomial margin of the Dirichlet-multinomial (Prop.~\ref{prop:nbdirmult}); the repair keeps every degree, block and segment match while removing repeated pairs and self-links, and moves or drops only what a full block cannot hold; the cap keeps the total; on a small bonds table the second key's degree KS distance falls from above 0.2 (\texttt{segment}) to below 0.05 (\texttt{auto}) with no self-bonds, and \texttt{segment} adds no decision.

%=====================================================================
\section{Time series and event sequences (v3.2)}\label{sec:v32ts}
% numbered from 1200 so that the parallel 2.x/3.x threads (600 privacy, 700 generators, 800 governance) do not collide
\newcounter{savedpropts}\setcounter{savedpropts}{\value{proposition}}
\setcounter{proposition}{1199}
Version 3.2 adds two models of time (\texttt{relsynth/timeseries}, guide \texttt{docs/timeseries.md}, design record
\texttt{docs/adr/0004-time-series.md}). An \emph{event-sequence} table (\texttt{sequence:} in the spec) holds the
events of a parent entity, such as a customer's orders: the engine still decides how many events $n_i$ entity $i$
has and what each holds, and the sequence model redraws their times jointly. A \emph{series} table
(\texttt{series:}) holds regular readings of an entity, such as a meter's hourly load, and is drawn by a seasonal
autoregressive model. The checks are \texttt{tests/unit/test\_timeseries.py} and C1200--C1206 in
\texttt{tests/science/test\_timeseries\_checks.py}; numbers on real data are in the T12 report.

\paragraph{Event sequences.} Times are in days. Entity $i$ has segment $k_i$ (the engine's latent segment), a window
$[A_i,B_i)$ (from \texttt{sequence.window}, or from the check rules that bound the time by a column of the entity),
and events $t_{i1}\le\dots\le t_{in_i}$. Segment $k$ has a piecewise-constant calendar intensity $\lambda_k$ on day
(or hour) cells, $\lambda_k(\text{cell})\propto P_{kb}\,w_{\mathrm{dow}}\,h_{\mathrm{hour}}$, with
$\Lambda_k(t)=\int_{T_0}^t\lambda_k$ normalised to one unit per average day. In operational time $u=\Lambda_k(t)$,
with $U_i=\Lambda(B_i)-\Lambda(A_i)$, the history of $i$ is: the head share $h_i=(u_{i1}-\Lambda(A_i))/U_i$; for
$n_i\ge2$ the span share $q_i=(u_{in_i}-u_{i1})/(\Lambda(B_i)-u_{i1})$; for $n_i\ge3$ the gap shares
$g_{ij}/(u_{in_i}-u_{i1})$, $j=1,\dots,n_i-1$. The model is: $(h_i,q_i)\mid n_i$ from a Gaussian copula per count
stratum ($1,2,3$--$4,5$--$8,\dots$) with quantile-grid marginals; gap shares
$\sim\mathrm{Dirichlet}(\alpha_i\mathbf 1)$ with $\log\alpha_i\sim N(\mu_k,\sigma_k^2)$.

\begin{proposition}[Time change]\label{prop:tschange}
Let $N$ be a point process on $[T_0,T_1)$ with intensity $\lambda>0$ and compensator $\Lambda$. Then
$\{\Lambda(t_j)\}$ is a unit-rate Poisson process on $[0,\Lambda(T_1))$, and conversely
$\{\Lambda^{-1}(u_j)\}$ has intensity $\lambda$ when $\{u_j\}$ is unit-rate Poisson. For the piecewise-constant
$\Lambda$ above, $\Lambda$ and $\Lambda^{-1}$ are exact (linear on each cell), and the Gamma--Poisson and Dirichlet
posteriors of $(P,w,h)$ give $\mathbb E\,\lambda_k$ equal to the observed bucket, weekday and hour rates up to the
prior pseudo-counts. When the profile is read as a rate per entity at risk, multiplying it by
$(\text{real}_b+1)/(\text{simulated}_b+1)$ until the bucket counts of the fitted model equal the real ones is a
fixed-point (method of simulated moments) step whose fixed point reproduces the bucket totals in expectation.
\end{proposition}
\begin{proof}
The first part is the random time change theorem (Daley and Vere-Jones 2003, Thm.~7.4.I; Brown et al.\ 2002). For
the converse, $P(\text{no event of }\Lambda^{-1}(u)\text{ in }(s,t])=e^{-(\Lambda(t)-\Lambda(s))}$ and independent
increments carry over through the monotone map. Exactness: on a cell of width $\Delta$ with mass $m$,
$\Lambda(t)=C+m(t-t_c)/\Delta$, which \texttt{Intensity.forward}/\texttt{inverse} evaluate in closed form. The
posterior means are those of $\mathrm{Dirichlet}(n_{kb}+c\,p_b)$, $\mathrm{Gamma}(1+n_j,\mathrm{days}_j)$ and
$\mathrm{Dirichlet}(1+n_{\mathrm{hour}})$. For the last sentence: the expected simulated count of bucket $b$ is a
continuous, increasing function of its own profile weight at fixed others, so the multiplicative update is the
iterative proportional fitting map, whose fixed point equates expected and real counts (Deming and Stephan 1940).
\end{proof}

\begin{proposition}[Gamma renewal bridge]\label{prop:tsbridge}
Let the gaps $G_1,\dots,G_m$ of a renewal process be i.i.d.\ $\mathrm{Gamma}(\alpha,\beta)$. Then, given the first
and last event, i.e.\ given $S=\sum_jG_j$, the shares $G/S$ are $\mathrm{Dirichlet}(\alpha\mathbf 1)$ and independent
of $S$, whatever $\beta$. The sampler, which places the interior events at the cumulative Dirichlet shares of the
span in operational time, draws exactly the gamma renewal bridge in operational time.
\end{proposition}
\begin{proof}
Lukacs (1955): for independent gammas with a common scale, $S\sim\mathrm{Gamma}(m\alpha,\beta)$ and $G/S$ is
Dirichlet and independent of $S$ (the joint density factorises after the change of variables
$(G)\mapsto(S,G/S)$ with Jacobian $S^{m-1}$). The sampler draws $\gamma_j\sim\mathrm{Gamma}(\alpha_i)$ and sets the
shares $\gamma/\sum\gamma$, which is the same construction. $\alpha=1$ gives uniform order statistics (Poisson),
$\alpha<1$ bursts and $\alpha>1$ regularity. C1201 tests shapes 0.3, 1 and 3.
\end{proof}

\begin{proposition}[Posterior of the sequence parameters]\label{prop:tspost}
(1) With $\log\alpha_i\sim N(\mu,\sigma^2)$ integrated out by $24$-node Gauss--Hermite quadrature, the gap likelihood
of an entity is $\int\mathrm{Dir}(s_i;\alpha\mathbf 1)\,\phi(\log\alpha)\,d\log\alpha$ up to a quadrature error that
decays geometrically in the node count for analytic integrands; the MAP of $(\mu,\log\sigma)$ under a Normal prior
centred on the pooled fit and its Laplace approximation are consistent and asymptotically normal (Bernstein--von
Mises). (2) The life-time marginals at their posterior mean are the empirical quantile functions on the grid; their
posterior draw gives the $N_Q$ cells $\mathrm{Dirichlet}(m/N_Q)$ masses, the Bayesian bootstrap of the cells (Rubin
1981), and the copula correlation the Fisher-$z$ posterior $N(\operatorname{atanh}\hat r,1/(m-3))$.
(3) No entity value is stored: a marginal is $N_Q+1=201$ quantile knots.
\end{proposition}
\begin{proof}
(1) The Dirichlet log-density is analytic in $\log\alpha$, so the Gauss--Hermite error bound applies (Davis and
Rabinowitz 1984, Ch.~4); MAP consistency and the Laplace limit follow from the regularity of a finite-dimensional
exponential-family mixture with a positive continuous prior (van der Vaart 1998, Thm.~10.1). (2) With equal cell masses
the piecewise-linear CDF through the knots is the empirical quantile function interpolated on the grid. Cell masses
with exchangeable Dirichlet weights of total concentration $m$ have the moments of the Bayesian bootstrap of $m$
observations aggregated to $N_Q$ cells (aggregation property of the Dirichlet). The Fisher transform of a sample
correlation of $m$ bivariate normals is approximately $N(\operatorname{atanh}\rho,1/(m-3))$ (Fisher 1921). (3) By
construction of \texttt{\_Marginal}. C1202 checks recovery and coverage.
\end{proof}

\begin{proposition}[The sequence hook keeps the skeleton and the rules]\label{prop:tshook}
The hook changes only the time column (or the calendar foreign key) of the rows of a sequence table, and columns
computed from it. Row counts, keys, links and the other tables are unchanged, and draws use their own random
streams, so the other tables of a dataset are identical with and without \texttt{sequence}. A row that breaks a check
rule the engine's row satisfied is redrawn with its entity up to \texttt{redraws} times; entities still breaking a
rule get their engine times back. Hence no rule the 1.0 output satisfied is broken, and each accepted entity history
has the model's law restricted to the rules.
\end{proposition}
\begin{proof}
\texttt{apply\_sequence} writes only \texttt{tc} (or the calendar key) and then calls \texttt{\_computed} for the
derived columns. The parameter draw uses \texttt{SeedSequence([seed, 1202, run, k])} and the sampling draw
\texttt{SeedSequence([seed, 1201, dataset, t])}, both separate from the engine's streams, so the engine's draws are
unchanged (the earlier version took its draw from the parameter stream and shifted every later table; C1203 compares
the tables). The rejection argument is that of Prop.~\ref{prop:genskeleton}(3), per entity. Reverted entities have
the 1.0 times, which satisfied the rules.
\end{proof}

\begin{proposition}[Calibrated limits]\label{prop:tscalib}
Let $T$ be a distance between two populations whose null law for halves of the real entities is approximately
normal on the log scale with variance proportional to $1/N_a+1/N_b$. (1) The limit
$\exp(\bar\ell+t_{S-1,0.99}\,s_\ell\sqrt{1+1/S})$, rescaled by $\kappa$ (Prop.~\ref{prop:splitnull}), keeps level $1\%$ for an
independent draw of the real law. (2) For a posterior predictive draw, the parameter uncertainty adds about one
sampling variance at the real size $N$, so $\kappa^2=(2/N+1/M)/(1/N_a+1/N_b)$ keeps the level. (3) For series that
share common shocks, both halves of a split share the shock path, which the split cannot see; comparing half $A$ with
a moving-block bootstrap (weekly blocks, the same blocks for every series) of half $B$ and taking $\kappa=1$ gives a
null with two shock paths, as in a real-versus-synthetic comparison, and more entity noise, so the limit is
conservative.
\end{proposition}
\begin{proof}
(1) is Prop.~\ref{prop:splitnull} applied to the sequence distances. (2) A posterior predictive draw is
$y^*\sim p(\cdot\mid\theta^*)$, $\theta^*\sim p(\theta\mid y)$. By the law of total variance,
$\operatorname{Var}(T(y^*)-T(y))=\operatorname{Var}_{\theta}(\text{sampling at }M)+\operatorname{Var}(\mathbb E[T\mid\theta^*]-T(y))$,
and under Bernstein--von Mises the second term is the sampling variance at $N$ plus the posterior variance, which is
asymptotically also the sampling variance at $N$. (3) The moving-block bootstrap reproduces the sampling law of
smooth functionals of a weakly dependent series (K\"unsch 1989), so $B^*$ has an independent copy of the shock
path's sampling error; $\kappa=1$ uses the half sizes, larger noise than at full size. C1204 measures 4 false failures
in 330 independent checks (1.2\%) and 0 in 192 predictive checks, and the planted Poisson-for-bursty defect fails;
without (3), two independent populations of shared-shock meters fail the ACF and spectrum checks (T12
report), and with it they pass.
\end{proof}

\paragraph{Series.} For a value column of a series table, $z=\log(1+y)$ per entity is decomposed into a level, daily,
weekly and (spans of at least 300 days) yearly profiles (principal component scores), and a residual
$\sigma(\lambda f_t+\sqrt{1-\lambda^2}\,\varepsilon_t)$, where $f$ is a common factor and $\varepsilon$ an AR($p$)
process (parametrised by partial autocorrelations, Barndorff-Nielsen and Schou 1973) times a seasonal AR at lag one
day. Entity parameters are drawn from a Gaussian mixture over the standardised parameter vectors with
Normal--inverse-Wishart posteriors.

\begin{proposition}[Series model: stationarity and range]\label{prop:tsseries}
(1) Every AR polynomial the sampler uses is stationary: partial autocorrelations in $(-1,1)$ map one to one onto
stationary AR coefficients. (2) The seasonal filter $Y_t=\Phi Y_{t-P}+\sqrt{1-\Phi^2}\,V_t$ with $|\Phi|<1$ and
unit-variance stationary $V$ keeps unit variance and has lag-$P$ correlation $\Phi$ plus the lag-$P$ correlation of
$V$ times $(1-\Phi^2)$-weighted terms; in particular it is $\Phi$ for white $V$. (3) Sampled parameter vectors lie in
the Mahalanobis ball of probability $0.99$ of their mixture component and in the real range widened by $5\%$, and
$\sigma\ge0.01$, so no sampled series has a level, amplitude or noise outside what the real series show.
\end{proposition}
\begin{proof}
(1) The Durbin--Levinson recursion is a bijection between $(-1,1)^p$ and the stationary region (Barndorff-Nielsen
and Schou 1973; Monahan 1984). (2) Stationarity of the first-order recursion in $Y_{t-P}$ with $|\Phi|<1$; the
variance solves $v=\Phi^2v+(1-\Phi^2)$, so $v=1$; the lag-$P$ covariance is $\Phi v+\sqrt{1-\Phi^2}\,
\operatorname{Cov}(Y_{t-P},V_t)$, which is $\Phi$ when $V$ is white. C1205's companion unit test checks both. (3) By
construction (\texttt{\_Mixture.sample} rejects draws outside the $\chi^2_{d,0.99}$ ball;
\texttt{SeriesModel.\_fit\_params} clips).
\end{proof}

\paragraph{Incremental runs.} \texttt{relsynth extend} adds the events of the days after a dataset. For entity $i$
with base count $n_i$, exposure $U_i$ and segment rate law $\mathrm{Gamma}(a_k,b_k)$, the rate has posterior
$r_i\sim\mathrm{Gamma}(a_k+n_i,b_k+U_i)$, the gap shape $\alpha_i$ is drawn from the segment's random effect, and new
gaps are $\mathrm{Gamma}(\alpha_i,\alpha_ir_i)$ (mean $1/r_i$) in operational time. Every uniform comes from a
counter-based hash of (seed, entity key, event number, purpose) and is turned into a variate by its inverse CDF.

\begin{proposition}[Incremental runs equal one long run]\label{prop:tsincr}
(1) The first new gap of an entity whose last event was $c$ units before the start is drawn from
$\mathrm{Gamma}(\alpha,\beta)$ conditioned on exceeding $c$; then the new events have the law of the renewal process
continued past the start given its history. (2) The events of $[s,e)$ do not depend on $e$, so runs of $[s,e_1)$,
$[e_1,e_2)$, $\dots$ with their outputs concatenated equal one run of $[s,e_K)$, for any step (day, week, month).
\end{proposition}
\begin{proof}
(1) Given the last event at $u_{\mathrm{last}}$ and no event in $(u_{\mathrm{last}},s)$, the current gap $G$ has the
law of $G\mid G>c$, $c=s-u_{\mathrm{last}}$, by the renewal property; it is drawn by inverse CDF:
$G=Q_\Gamma\big(1-U\,\bar F_\Gamma(c)\big)$, with $Q_\Gamma$ and $\bar F_\Gamma$ the regularized incomplete gamma
inverse and complement. Later gaps are independent of the past. (2) Event $j$ of an entity is
$u_{\mathrm{last}}+G_1+\dots+G_j$, where each $G$ is a deterministic function of the hash of
(seed, key, $j$, purpose) and of $(\alpha_i,r_i)$, which are functions of the hash and of the base dataset only. The
run stops at the first event at or after $e$, so the events before $e_1$ are the same for every $e\ge e_1$. The
\texttt{Extender} uses the base dataset (not earlier deltas) for every state and continues key numbering from the
previous output, so step-wise and one-shot runs give equal tables (unit test
\texttt{test\_extension\_is\_step\_invariant}); C1206 checks (1) against a simulated renewal and (2) for prefixes.
\end{proof}
\setcounter{proposition}{\value{savedpropts}}

%=====================================================================
\section{Partitioned generation (release 3.5)}\label{sec:runtime}
\texttt{synthesis.runtime} generates the fact tables of a dataset in partitions (\texttt{relsynth.runtime}, docs/runtime.md). The driver draws the parameters, every entity, and per fact $f$ the \emph{state} $S_f$: the anchor rows' segments $z^1$, key columns and degree weights $w^1$ (after the eligibility mask), and the partner rows' $z^2,w^2$ and extra-dimension weights. The anchor rows $\{0,\dots,n_1-1\}$ are cut into consecutive blocks $A_0,\dots,A_{P-1}$. Partition $p$ draws the rows of $A_p$ with the generator $G_{f,p}=\mathrm{Philox}(\mathrm{SHA256}(s,d,f,p))$ ($s$ the seed, $d$ the dataset index). Rows are joined in partition order, surrogate keys numbered over the join and generated codes drawn over it from $\mathrm{Philox}(s,d,f,\texttt{codes},c)$; aggregates and the external form follow on the driver from $\mathrm{Philox}(s,d,\texttt{tail})$. A fact is split in one of three modes chosen by \texttt{partition\_rule}: \emph{partition} (all per partition), \emph{structure} (joins on the driver, Prop.~\ref{prop:rtstructure}) or \emph{whole} (one partition). Code: \texttt{runtime/partitioned.py}, \texttt{runtime/streams.py}, \texttt{Synthesizer.\_fact\_state}, \texttt{\_fact\_structure}, \texttt{\_fact\_attributes}.

\begin{proposition}[Worker independence; Prop.~900]\label{prop:rtdet}
Let the blocks be a function of $(S_f, \texttt{partition\_rows})$, let every draw of partition $p$ come from $G_{f,p}$ or from a cache whose entries are drawn from streams keyed by their names, and let the executor return the partitions' results in partition order. Then the synthetic dataset is a function of $(s, d, \texttt{partition\_rows})$ and the fitted model only: it is the same for every executor, every number of workers and every order in which partitions run.
\end{proposition}
\begin{proof}
The entity phase and $S_f$ are computed on the driver from streams keyed by $(s,d)$ and $(s,d,f,\texttt{state})$, the same in every run. The block boundaries depend on $S_f$ and \texttt{partition\_rows} only (\texttt{blocks\_for} on expected rows; \texttt{blocks\_by\_rows} on the driver's joins). A task reads $S_f$, the entities and its block, and draws only from $G_{f,p}$; the cached Bayesian-bootstrap weights are drawn, whenever first needed, from $\mathrm{Philox}(s,d,\texttt{bb},\text{name})$ (\texttt{KeyedRunCache}), so their value does not depend on which task needs them first. Hence each task's result is a function of $(S_f, p)$. The join in partition order, the numbering and the tail are deterministic functions of the ordered results and keyed streams. Distinct label tuples give distinct Philox keys except with probability $2^{-128}$ per pair (SHA-256 truncated to 128 bits), and Philox streams with distinct keys are the independent streams of Salmon et al.\ (2011).
\end{proof}

\begin{proposition}[Partition mode keeps the law; Prop.~901]\label{prop:rtlaw}
Given $S_f$, suppose the rows of $f$ are generated as: (single) $N_i$ rows for anchor $i$, independent over $i$ with law $L_{z_i}(w_i)$ (Poisson, or the DP count law with its bootstrap weights fixed for the dataset); or (pairs, no degree law) for each block $(k,l)$, $N_{kl}\sim\mathrm{Poisson}(\eta^*_{kl}W^1_kW^2_l)$ rows, each with anchor drawn $\propto w^1$ in segment $k$ and partner $\propto w^2$ in segment $l$; and, given its anchor (and partner, block), each row's attributes, rule redraws and finishing are drawn independently of other anchors' rows. If every unique key, group effect, running balance and sequence of $f$ contains the anchor key, then the rows drawn per block $A_p$ and joined have the same law as the rows drawn for all anchors at once.
\end{proposition}
\begin{proof}
(single) The $N_i$ are independent, so drawing them block by block changes nothing. (pairs) Mark each of the $N_{kl}$ rows by the block of its anchor. By the marking (thinning) theorem for Poisson processes, the counts $N^p_{kl}$ of rows marked $p$ are independent $\mathrm{Poisson}(\eta^*_{kl}W^{1,p}_kW^2_l)$ with $W^{1,p}_k=\sum_{i\in A_p,z_i=k}w^1_i$, and given $N^p_{kl}$ the anchors are i.i.d.\ $\propto w^1$ on $A_p\cap\{z=k\}$ and the partners i.i.d.\ $\propto w^2$ on segment $l$. This is exactly what partition $p$ draws (\texttt{\_sample\_poisson\_pairs} on the restricted state). Conversely, independent Poisson counts per block superpose to $\mathrm{Poisson}(\eta^*_{kl}W^1_kW^2_l)$ and the union of the per-block anchor draws is the global multinomial allocation. Rows given their anchors are then drawn by the same per-row kernels (\texttt{\_fact\_rows}: ``rows are independent given $(u,v,\mathrm{blk})$''; \texttt{\_refill}: rejection per row from $Q(\text{row}\mid\text{anchor},\text{valid})$). A unique key containing the anchor key compares only rows of one anchor, which lie in one partition; keeping the first of an exchangeable sequence of rows has the same law in any order. Group effects, running balances and sequences within the anchor key act on one anchor's rows. The weights $w^1,w^2$ are normalised once on the driver (in $S_f$), not per block.
\end{proof}

\begin{proposition}[Structure mode keeps the law; Prop.~902]\label{prop:rtstructure}
For an occupancy fact or a pair fact with degree laws, let the joins $(u,v,\mathrm{blk})$ be drawn on the driver from $S_f$ by the single-process samplers (\texttt{\_fact\_structure}), and let partition $p$ draw the attributes of the rows with $u\in A_p$. Under the key conditions of Prop.~\ref{prop:rtlaw}, the joined rows have the law of the single-process sampler.
\end{proposition}
\begin{proof}
The joins are drawn by the same function of $S_f$ (Props.~\ref{prop:linkdeg} and the occupancy model), so their law is unchanged. Given the joins the remaining steps are the per-row kernels of Prop.~\ref{prop:rtlaw}, and the occupancy key (anchor, partner) contains the anchor key.
\end{proof}

\begin{proposition}[Keys and links across partitions; Prop.~903]\label{prop:rtkeys}
In a partitioned dataset every primary key is unique, every foreign key resolves, and every unique key and check rule of a fact that the single-process sampler guarantees is guaranteed.
\end{proposition}
\begin{proof}
Entities are drawn on the driver as before. A fact's anchor, partner and extra keys are taken from the key arrays of the driver's entities (or of the joined parent fact for nested facts), so they resolve. Surrogate keys generated in finishing are replaced by $1..n$ over the join, and generated codes by one draw of \texttt{\_unique\_codes} over the join, so both are unique. A unique key containing the anchor key is enforced within each partition and anchors are disjoint across partitions; a fact with any other unique key is \emph{whole}. Check rules are row rules, enforced per row by the same redraws.
\end{proof}

\paragraph{What is not partitioned.} A fact is \emph{whole} when a unique key, a group effect or a running balance does not contain the anchor key, when it has link affinity (a share of the whole table), a sequence numbered over the table, or is nested on a fact other than its anchor. Entities are always drawn on the driver: their relational step moves links across rows and enforces keys over the whole table.

\paragraph{Checks.} \texttt{tests/unit/test\_runtime.py}: streams depend only on their labels and the keyed cache on the key, not the order; blocks are fixed by the model; local, 2, 8 and 64 worker processes give identical tables on two scenarios (one with a \emph{structure} fact) and another dataset index differs (Prop.~\ref{prop:rtdet}); primary keys unique, foreign keys resolving and the declared unique key holding across partitions (Prop.~\ref{prop:rtkeys}); the scorecard of a partitioned dataset matches the single-process one (Props.~\ref{prop:rtlaw}, \ref{prop:rtstructure}); Spark gives the same tables (slow). On Chinook (\texttt{partition\_rows} 2{,}000; both facts in \emph{structure} mode) the partitioned datasets score 140/140 and are identical with the local executor, 8 and 64 processes and Spark (\texttt{local[2]}); on Contoso (\texttt{OrderRows} \emph{whole}: a group effect crosses anchors) they score 246/246 with the local executor and 4 processes, identical; no primary or foreign key violation in either (\texttt{roadmap/reports/T10/scripts/runtime\_realdb.py}).

\subsection{Partition-aware statistics}

\begin{proposition}[Drift tests across partitions; Prop.~904]\label{prop:rtdrift}
Let a table be split into $k$ partitions by its partition fields, and for each of $m$ columns let $p_j$ be the p-value of Welch's heteroscedastic ANOVA (numbers) or of Pearson's chi-square test of homogeneity after lumping levels until every expected count is at least 5 (categories), computed from the per-partition sufficient statistics (counts, means and sample variances; level counts). Flag column $j$ when its Holm-adjusted p-value is below $\alpha$ and its effect size ($\eta^2$, or Cramér's $V$) is at least a threshold $e_j\ge 0$. If every column has the same law in every partition, the probability that any column is flagged is at most $\alpha$ asymptotically (as the partition sizes grow).
\end{proposition}
\begin{proof}
Welch's statistic is a function of $(n_i,\bar x_i,s_i^2)$ only, and under equal means it is asymptotically $F(k-1,\nu)$ with Welch's approximate $\nu$, whatever the variances (Welch, Biometrika 1951); a common law implies equal means. Pearson's statistic is a function of the counts and is asymptotically $\chi^2_{(k-1)(c-1)}$ under homogeneity; lumping levels by their pooled totals is a function of the margins and preserves homogeneity of the lumped table. So each $p_j$ is asymptotically valid: $\Pr(p_j\le u)\le u$. Holm's step-down procedure controls the family-wise error rate at $\alpha$ for any dependence among valid p-values (Holm 1979). The effect threshold only removes flags, so the rate cannot grow.
\end{proof}

\paragraph{Check.} \texttt{tests/unit/test\_drift.py}: Welch's $F$ and the chi-square agree with direct and scipy computations; over 300 homogeneous tables of 1{,}200 rows in 12 monthly partitions (two numbers, two categories, $e_j=0$) the family-wise rejection rate is at most $\alpha+3\,\mathrm{SE}$; a shifted mean and a changed category share in the second half-year are flagged and the unchanged columns are not; the same on an Iceberg table partitioned by \texttt{month(day)} and an identity field, read through DuckDB.

%=====================================================================
\section{Synthetic text (v5.0)}\label{sec:v50text}
% numbered from 1500 (wave D range of T14)
\newcounter{savedproptx}\setcounter{savedproptx}{\value{proposition}}
\setcounter{proposition}{1499}
Version 5.0 adds row-conditioned free text (\texttt{relsynth/textgen}, guide \texttt{docs/text.md}, design record
\texttt{docs/adr/0010-synthetic-text.md}). A text column chooses a method: the 1.0 n-gram model (Props.~12, 13), the
DP n-gram release, a local language model with optional DP fine-tuning, or Private Evolution. A PII layer replaces
detected personal data by typed placeholders before any model sees a text, and a leakage audit measures what the
output still reveals. The checks are \texttt{tests/unit/test\_textgen.py} and C1500--C1507 in
\texttt{tests/science/test\_text\_checks.py}; numbers on data are in the T14 report.

\paragraph{Setting.} A \emph{person} $u$ (the value of \texttt{text.dp.unit\_column}) holds rows $(x_r,c_r)$, a text
and its conditioning values. Neighbouring datasets $D\simeq D'$ differ by all rows of one person (add/remove).
$\rho$-zCDP (Bun and Steinke 2016) and $(\varepsilon,\delta)$-DP are as in Section~\ref{sec:v25dp}; the conversion
from $\rho$ to $\varepsilon(\delta)$ is the one of Canonne, Kamath and Steinke (2020) used by the ledger.

\begin{proposition}[Canary exposure under DP]\label{prop:txexposure}
Let a canary person $u$ hold copies of $p\,s$, where the secret $s$ is uniform on a candidate set $R$ and independent
of everything else, and let $M$ be $(\varepsilon,\delta)$-DP for $u$. Let $\mathrm{rank}(s)$ be the position of $s$
among all candidates ordered by any score computed from $M(D)$ and public data, ties broken by independent uniform
randomness, and $\mathrm{exposure}=\log_2|R|-\log_2\mathrm{rank}(s)$. Then for every $t\ge0$
\[P(\mathrm{exposure}\ge t)\le e^{\varepsilon}\,2^{-t}+\delta.\]
For $\rho$-zCDP plus a partition-selection $\delta_0$, the bound holds with
$\min_{\delta'}\{e^{\varepsilon(\rho,\delta')}2^{-t}+\delta'\}+\delta_0$.
\end{proposition}
\begin{proof}
Let $D^-$ be the data without $u$. The event $E_t=\{\mathrm{rank}(s)\le|R|2^{-t}\}=\{\mathrm{exposure}\ge t\}$ is a
post-processing of the output (the score, the public candidate list and the tie-breaking coins), so
$P(E_t\mid D)\le e^\varepsilon P(E_t\mid D^-)+\delta$ for every value of $s$, hence also after averaging over $s$. Under
$D^-$ the output is independent of $s$, and with random tie-breaking the rank of a uniform $s$ is uniform on
$\{1,\dots,|R|\}$ for any fixed scores, so $P(E_t\mid D^-)=\lfloor|R|2^{-t}\rfloor/|R|\le2^{-t}$. For zCDP every pair
$(\varepsilon(\rho,\delta'),\delta')$ is valid, and the partition-selection mechanisms add their $\delta_0$ by
composition.
\end{proof}

\begin{proposition}[DP n-gram release]\label{prop:txngram}
The release of \texttt{DPNgramModel.fit\_dp} with levels $n=0,\dots,N-1$, budgets $\rho_n$, $\rho_\emptyset$ and
selection deltas $\delta_n$ is $(\rho_\emptyset+\sum_n\rho_n)$-zCDP up to the event of probability at most
$\sum_n\delta_n$ (user level), and every text it generates is post-processing.
\end{proposition}
\begin{proof}
After bounding, person $u$ keeps at most $k$ texts of at most $L$ tokens. At level $n$, $u$ contributes the indicator
of a set $K_u$ of at most $m$ keys, so the count vector over keys changes by a $0/1$ vector with at most $m$ ones:
$L_2$ sensitivity $\sqrt m$, and at most $1$ on any key. The discrete Gaussian with scale
$\sigma_n=\lceil\sqrt m\rceil/\sqrt{2\rho_n}$ is $\rho_n$-zCDP on the keys present in both datasets
(Canonne et al.\ 2020, Thm.~4; Prop.~900). Keys held by $u$ alone have true count at most $1$ and are released only if
$1+X>\tau_n$; with $\tau_n$ from Prop.~903 for count bound $1$ and $m$ keys per person, the probability that any of
them is released is at most $\delta_n$, which is the partition-selection argument of Prop.~903 (Gopi et al.\ 2020;
Kim et al.\ 2021, Sec.~3). The candidate set at level $n$ depends on the data only through the release at level
$n-1$ (the shorter context and the context must be released) and is otherwise public. The share of empty texts is a
2-vector with sensitivity at most $k$ in $L_2$, released with $\rho_\emptyset$. Adaptive composition of zCDP adds the
$\rho$'s (Bun and Steinke 2016, Lemma~2.3). The back-off model, the sampler and the placeholder filler use only the
released counts, the public start/end tokens and fresh randomness.
\end{proof}

\begin{proposition}[User-level DP-GD]\label{prop:txdpgd}
\texttt{LMTextModel} with \texttt{dp} runs $T$ steps $\theta_{t+1}=\mathrm{opt}\big(\theta_t,(\sum_u
\mathrm{clip}_C(g_u(\theta_t))+Z_t)/\tilde n\big)$ with $Z_t\sim N(0,\sigma^2C^2I)$ and $\tilde n$ a noisy count of
people. With $\sigma=\sqrt{T/(2\rho_{\mathrm{steps}})}$ and $\tilde n$ released with $\rho_{\mathrm{n}}$, the final
weights are $(\rho_{\mathrm{steps}}+\rho_{\mathrm n})$-zCDP for one person, and so is every text generated from them.
\end{proposition}
\begin{proof}
$g_u$ is the gradient of the mean loss over the at most $k$ kept texts of $u$; after clipping,
$\|\mathrm{clip}_C(g_u)\|_2\le C$, so the sum over people has $L_2$ sensitivity $C$ under add/remove of $u$, given
$\theta_t$. The Gaussian mechanism with variance $\sigma^2C^2$ is then $1/(2\sigma^2)$-zCDP (Bun and Steinke 2016,
Prop.~1.6); every person is used in every step (no sampling), so adaptive composition over $T$ steps gives
$T/(2\sigma^2)=\rho_{\mathrm{steps}}$. The count has sensitivity $1$. The optimiser state, the division and the model
are post-processing. The noise is floating-point (as in Opacus); the guarantee is stated for exact Gaussian noise
(Mironov 2012 for the caveat).
\end{proof}

\begin{proposition}[Private Evolution votes]\label{prop:txpe}
In \texttt{PETextModel.fit\_dp} with $T$ iterations, each vote histogram has $L_2$ sensitivity at most
$S=\texttt{SCALE}$ for one person, the iterations cost $\sum_t\rho_t$, the class levels cost $\rho_c$ and $\delta$, and
the final pool and weights are post-processing; the guarantee is user-level.
\end{proposition}
\begin{proof}
The embedding is a function of the text and a public constant only (hash and projection seeded by 1503), and the
proposals are produced by a public API from the previous noisy release, so the only data access per iteration is the
histogram. Person $u$ with votes $v_u\in\mathbb N^{|S|}$ contributes $\lfloor Sv_u/\|v_u\|_2\rfloor$ (elementwise), whose
$L_2$ norm is at most $S$; removing $u$ removes that vector. The discrete Gaussian with scale
$S/\sqrt{2\rho_t}$ is $\rho_t$-zCDP (Prop.~900), and adaptive composition adds the costs. The class levels are a
partition selection with $k+1$ keys per person (Prop.~903). Resampling, the API's variations and the final draws use
the noisy histograms and public randomness only (Lin et al.\ 2024, Thm.~1; Xie et al.\ 2024).
\end{proof}

\begin{proposition}[Expected exposure]\label{prop:txexpect}
Under the assumptions of Prop.~\ref{prop:txexposure},
$\mathbb E[\mathrm{exposure}]\le(\varepsilon+1)/\ln2+\delta\log_2|R|$.
\end{proposition}
\begin{proof}
Exposure is nonnegative and at most $\log_2|R|$, so $\mathbb E[\mathrm{exposure}]=\int_0^{\log_2|R|}P(\mathrm{exposure}
\ge t)\,dt\le\int_0^{\infty}\min(1,e^\varepsilon2^{-t})\,dt+\delta\log_2|R|$. With $t_0=\varepsilon/\ln2$ the integral is
$t_0+\int_{t_0}^\infty e^\varepsilon2^{-t}dt=\varepsilon/\ln2+1/\ln2$.
\end{proof}

\begin{proposition}[Near-duplicate search]\label{prop:txminhash}
For shingle sets $A,B$ with Jaccard similarity $s$, one MinHash coordinate agrees with probability $s$ (Broder 1997)
when the hash family is min-wise independent; with $b$ bands of $r$ coordinates the pair becomes a candidate with
probability $1-(1-s^r)^b$. Every reported similarity is the exact Jaccard similarity of a candidate pair, so the
reported counts of near duplicates are never too high, and a pair with similarity $s$ is missed with probability at
most $(1-s^r)^b$ ($2\cdot10^{-4}$ at $s=0.7$, $b=32$, $r=4$).
\end{proposition}
\begin{proof}
For a random permutation $\pi$, $\min\pi(A)=\min\pi(B)$ iff the minimiser of $\pi$ on $A\cup B$ lies in $A\cap B$,
which has probability $|A\cap B|/|A\cup B|$. Bands use disjoint coordinates, so they agree independently; a band
agrees with probability $s^r$ and the pair is a candidate if some band agrees. The implementation hashes with
$x\mapsto ax+b\bmod 2^{64}$ over 64-bit BLAKE2 hashes, which is only approximately min-wise independent; check C1505
measures the candidate rate against the formula. The exact verification step makes the counts conservative.
\end{proof}

\begin{proposition}[PII placeholders]\label{prop:txpii}
(a) No character of a span found by the recognizers reaches a text model: the model is fitted to the redacted texts.
(b) Surrogates are drawn from fresh randomness and reserved or fictitious ranges (RFC 2606 domains, 555-01xx numbers,
RFC 5737 addresses, Luhn-valid test cards, mod-97-valid IBANs of a test bank code), independent of the source values,
and each is format-valid. (c) With $x$ of $n$ labelled entities caught, the Wilson interval covers the recall with
probability close to $0.95$ for $n\ge 30$.
\end{proposition}
\begin{proof}
(a) \texttt{redact} rebuilds the text from the gaps between the merged spans and the placeholders, and
\texttt{TextColumn.fit} passes only the redacted list on. Overlapping spans are merged, so a partly overlapping detection
removes the union. (b) By construction of \texttt{surrogate}: the check digits are computed (Luhn: the last digit is
chosen to make the sum $0\bmod 10$; IBAN: ISO 7064 mod 97-10). (c) Wilson (1927); Brown, Cai and DasGupta (2001) for
its coverage. The recall is a property of the recognizers on the labelled distribution only, which is why the audit
runs on the output as well.
\end{proof}

\begin{proposition}[Consistency with the row]\label{prop:txconsist}
Let $f$ be a classifier fitted on real training rows that predicts a conditioning value $c$ from a text. If the
synthetic pairs $(\tilde x,\tilde c)$ have the law of the real pairs, then
$\mathbb E[\mathrm{acc}_{\mathrm{syn}}]=\mathbb E[\mathrm{acc}_{\mathrm{held\text{-}out}}]$, so the consistency ratio is $1$
up to sampling error; if the text is independent of the row, $\mathbb E[\mathrm{acc}_{\mathrm{syn}}]=\sum_c
q(c)\,P(f(\tilde x)=c)$, the chance agreement.
\end{proposition}
\begin{proof}
Given $f$ (fitted on data disjoint from both samples), $\mathrm{acc}$ is the mean of i.i.d.\ indicators
$1\{f(x)=c\}$ whose expectation depends only on the joint law of $(x,c)$. Under independence,
$P(f(\tilde x)=\tilde c)=\sum_cP(\tilde c=c)P(f(\tilde x)=c)$.
\end{proof}
\setcounter{proposition}{\value{savedproptx}}


%=====================================================================
\section{Evaluation thresholds}\label{sec:thresholds}
A synthetic sample of size $m$ drawn from a model fitted on $n$ real rows deviates from the real sample by the real sample's own noise plus the fit's noise plus the synthetic sampling noise, so the effective size for two-sample tests is $n_\text{eff}=(2/n+1/m)^{-1}$. The scorecard uses: the two-sample KS critical value $\sqrt{-\tfrac12\log(\alpha/2)}\,n_\text{eff}^{-1/2}$ at $\alpha=1\%$ (at least 0.08); for categorical columns $2.5$ times the expected TVD, $\tfrac12\sum_l\sqrt{2p_l(1-p_l)/(\pi n_\text{eff})}$ (at least 0.08); for association differences $2.5\,n_\text{eff}^{-1/2}$ (at least 0.15); for joins, the size of the ancestor table rather than the number of joined rows, since a join repeats each ancestor row. Without this, a table of 12 departments fails on noise alone, as the first runs showed. Since v0.3, rows whose level occurs in fewer than \texttt{suppress\_small\_counts} rows are left out of association pairs (a single record's coincidence is not a dependence the model should copy, and copying it would be a disclosure), the correlation ratio between a categorical and a numeric column is computed on ranks, and rows per partner are compared after rescaling by the realised size of the grand anchor (tree and grain tables have random sizes). Grain metrics report the duplicate rate of each unique key, the share of repeated pairs, and the share of synthetic rows whose combination of two reference keys never occurs in the real data where the real data covers its combinations densely. Since v0.4 every limit is at least the sampling limit of the quantity it tests: null-rate differences use $2.5\sqrt{p(1-p)/n_\text{eff}}$; an association pair uses the rows that carry information on it (non-null rows, and for a categorical the rows away from its mode: a flag set on 5 rows has 5 such rows, however long the table); joined columns use Kish's effective size $(\sum w)^2/\sum w^2$ of the ancestor rows weighted by how many fact rows each carries; tree fan-out and attribute-by-depth use the KS and TVD limits at the tree's size. Check C41 measures the false-alarm rate of the null-rate and association limits on independent samples from one law (0 of 400 at $n=59$).

%=====================================================================
\section{Checks without ground truth (real data)}\label{sec:real}
On real data the segments and parameters are unknown, so recovery checks are replaced by checks that need none. \texttt{local/replicate\_checks.py} runs them where the data lives and writes only aggregates:
\begin{description}
\item[R1] the scorecard of Section~\ref{sec:thresholds} (marginals, dependence, discriminator AUC, joins, cardinalities, trees, time, text, privacy, rules);
\item[R2] copula convergence: split $\hat R$ of every correlation and of $\log\det C$, and the largest change of a correlation between chain halves (and the Metropolis acceptance in \texttt{bayes} mode);
\item[R3] posterior predictive checks (Proposition~\ref{prop:ppc}, a parameter redraw before each copy): for rows per parent, share of parents without rows, dispersion, tree depth and leaves, busiest-month share and fact means, the real value must lie within $3$ standard deviations of $D$ synthetic copies (or within 3\%);
\item[R4] the scale rule (Proposition~\ref{prop:scale}) on the user's own schema;
\item[R5] structural validity: FK integrity and acyclic hierarchies;
\item[R6] stability: agreement (ARI) of segmentations from two seeds, the real-data stand-in for recovery;
\item[R7] reproducibility.
\end{description}
\paragraph{Privacy.} For entity tables with at least 3 modelled columns and 50 rows: exact-copy rate (at most $\max(1\%,\,2\times$ the real table's own duplicate rate$)$) and the distance-to-closest-record ratio, the 5th percentile of synthetic-to-real distances over the 5th percentile of real-to-real leave-one-out distances, in rank-normal and one-hot space (at least 0.5; on coarse tables split as in Proposition~\ref{prop:dcr}). Direct identifiers never pass through the model (\texttt{pii} columns are regenerated from formats). \emph{Check C04}: no copies, ratio 0.98, no shared emails.

\paragraph{Rehearsal.} \texttt{local/make\_sqlite\_demo.py} loads the enterprise scenario into SQLite with declared keys; introspection then drafts the spec from the catalog alone (Section~\ref{sec:doctor}).

%=====================================================================
\section{Diagnostics loop}\label{sec:doctor}
The bundle holds the spec (credentials masked), a column profile without values or level names, the model summary, the scorecard, the replication checks and any error traceback. \texttt{relsynth doctor} applies rules D01--D22 and R2--R6, each mapping a failed metric plus profile evidence to a cause and a spec patch (for example: a categorical column that is 95\% e-mail addresses $\to$ \texttt{pii}; all segments occupied and checks on that table failing $\to$ double the cap; rare text copied $\to$ raise $k$, shorten $L$; rows per parent too dispersed and a small parent table $\to$ frailty or copy mode). Version 0.3 adds: rows per parent \emph{less} dispersed than Poisson $\to$ a grain, with the unique key proposed from the profile's key candidates (D03, D21); a broken unique key (D18); repeated pairs (D19); combinations of reference keys the real data never has (D20, stratify); a stored computed column that disagrees with its rule (D22); and slow copula mixing $\to$ rank mode, longer chains, then a higher rare-column threshold (D07). Reference tables in copy mode are exempt from the leak rules. Issues whose cause is in the code (failed steps, broken foreign keys) are marked \emph{engine} and come back to the project with their traceback. \texttt{relsynth tune} loops run $\to$ doctor $\to$ apply.

\noindent\emph{Check C19}: four planted spec errors in the enterprise schema (email declared categorical, segment cap 1, text guards off, hierarchy depth capped at 3) are all identified from the bundle alone; tuning removes all four and raises the scorecard from 117/129 to 122/122 within three rounds.

%=====================================================================
\section{Changes to layer 1 found while building and testing}\label{sec:changes}
\begin{enumerate}
\item \textbf{Measures conditional on parents.} Segment-level means (layer~1's $m_{kl}$) lose the link between a sale amount and the product's own price; parent attributes are now conditioning columns of the fact copula (Section~\ref{sec:snow}).
\item \textbf{Degree correction.} Pure Poisson blocks under-disperse rows per parent; Gamma frailties, rescaled within segments (Proposition~\ref{prop:split}(c)), fix this without breaking the scale rule. Estimating a frailty needs $\ge30$ rows in the segment: an estimate from 6 executives once killed most synthetic hierarchies.
\item \textbf{Segment marginals} use a Dirichlet-process posterior instead of falling back to the table marginal below a row threshold, and a weak prior ($a=0.25$): with $a=2$, 7\% of synthetic rows in a 6-row segment took the wrong level.
\item \textbf{Small tables} draw segment labels by largest-remainder allocation of $n^*\pi$ rather than multinomially, and reference tables can be copied.
\item \textbf{Strongly dependent categoricals} are modelled jointly (Remark~\ref{rem:nominal}); nominal levels are ordered by their most associated column.
\item \textbf{Columns constant within a segment} are excluded from that segment's copula, whose correlations they cannot inform.
\item \textbf{Trees} use posterior predictive checks and a depth limit from the data.
\item \textbf{Text guards} count rows, not occurrences, so common phrases are allowed and rare details are not.
\item \textbf{(v0.3) Grain.} A unique key of two foreign keys makes a table an occupancy model of its parents' product, not a Poisson count (Propositions~\ref{prop:repeat}--\ref{prop:unique}).
\item \textbf{(v0.3) Structural zeros} in every prior that gives a combination positive mass: occupancy, activity, null rates, category levels and extra keys (Propositions~\ref{prop:zoib}--\ref{prop:sparsedir}), pooled within strata.
\item \textbf{(v0.3) Computed columns} are computed, never sampled, on both sides (Proposition~\ref{prop:computed}).
\item \textbf{(v0.3) Nullable keys} are a level, not a reason to drop the row (Proposition~\ref{prop:nullfk}).
\item \textbf{(v0.3) Copula}: the rank likelihood with parameter expansion replaces Metropolis on vine partial correlations as the default, with chain extension by $\hat R$, and rare columns are kept out (Propositions~\ref{prop:rank}, \ref{prop:rare}).
\item \textbf{(v0.4) Rows per parent} are chosen between Poisson/NB and a DP by Bayes factor, for facts and for links between entities, and DP counts of entity links are drawn given the segment totals (Propositions~\ref{prop:countbf}, \ref{prop:condcounts}); parents that always have children use $1+$Poisson (Proposition~\ref{prop:mandatory}).
\item \textbf{(v0.4) Categorical hierarchies} are joint categoricals over their observed combinations (Proposition~\ref{prop:fd}).
\item \textbf{(v0.4) Tiny tables} are stratified by an own column and generated with their stratum counts; hierarchies are conditioned on their size (Propositions~\ref{prop:attrstrata}, \ref{prop:treesize}).
\item \textbf{(v0.4) Values}: few-valued numbers are categoricals on their values, tail levels never leak a placeholder, aggregates keep the stored precision (Proposition~\ref{prop:tail}).
\end{enumerate}

%=====================================================================
\section{Limitations and open items}\label{sec:limits}
\begin{itemize}
\item \textbf{Copula mixing} is solved by the parameter-expanded rank sampler (Proposition~\ref{prop:rank}); \texttt{bayes} mode (vine Metropolis) remains available and still mixes slowly for near-deterministic pairs.
\item \textbf{Structure the data cannot show.} A rule that holds in every real row but is not declared (e.g.\ ``a filed form always has its header items'') is reproduced statistically, not exactly (S8: 99.8\%). Declaring it as a check rule makes it exact (Proposition~\ref{prop:rules}).
\item \textbf{Within-segment dependence is Gaussian.} Non-monotone relations inside a segment are captured only through more segments. ClavaDDPM and TabDDPM remain the baselines to compare against on real data (layer~1 item~9).
\item \textbf{Text is $n$-gram based}: fluent for templated enterprise text (tickets, notes), weak for long free prose. The foundation-model route needs a model endpoint inside the data perimeter.
\item \textbf{The prototype is in memory.} Layer~1's in-database statistics (GROUP BY pushdown, sketches) are the layer~3 path to $10^8$-row facts; every quantity the model uses is already an aggregate of that form.
\item \textbf{Formal DP} (layer~1 Algorithm~6) is not wired into this prototype; the privacy checks here are empirical.
\item \textbf{Four real databases so far} (the reporting model, Chinook, AdventureWorks DW, Contoso V2). Their generic lessons are built in and tested on S8--S11 and C57--C68; other real schemas will test other situations.
\item \textbf{Unknown members on facts} are restored under a missing-at-random assumption (Proposition~\ref{prop:unknown}); when unknown references depend on other columns the member should be modelled as an ordinary row.
\item \textbf{The frailty shape} is an empirical-Bayes (moment) estimate, not a posterior draw; the participation hurdle (Proposition~\ref{prop:hurdle}) corrects the zero share it can miss.
\item \textbf{Stated column distributions} are drawn independently of the other columns.
\item \textbf{Hierarchies of a handful of rows} (8 employees) are learned from one example of each relation. Even with strata and size conditioning, about one draw in five places a manager under the wrong person; the reproducible fix is to declare the hierarchy as reference data if it is not confidential.
\item \textbf{Regularities the source generator imposed} (each track bought at most twice, cells spread evenly over submissions) are under-dispersion on the partner side; the engine models counts per anchor, not per partner, and reports these.
\end{itemize}

%=====================================================================
\section{Results}\label{sec:results}
All 79 of 79 checks pass (507 of 507 individual assertions). Each row names the property it tests; the per-assertion values are in \texttt{checks/results.md}.

\begin{longtable}{@{}l p{8.2cm} p{3.6cm} c r@{}}
\caption{Layer 2 checks on fake data with known ground truth.}\label{tab:results}\\
\toprule
Id & Check & Property & Items & Pass \\ \midrule \endfirsthead
\toprule Id & Check & Property & Items & Pass \\ \midrule \endhead
C00 & Layer 1 algorithm checks still pass & Layer 1 paper, Table 4 & 1/1 & yes \\
C01 & Copula synthesis preserves every marginal (mixed types, nulls) & Prop. 1, 2 & 8/8 & yes \\
C02 & Posterior of the latent correlation covers the true copula & Prop. 3 (Alg. 2 of layer 1) & 4/4 & yes \\
C03 & Missingness that depends on other columns is reproduced & Prop. 2 (null as lowest latent bin) & 4/4 & yes \\
C04 & Single-table privacy: no copies, distance to closest record & Sec. 9 & 3/3 & yes \\
C05 & Star schema: latent segments of both dimensions are recovered & Prop. 4 (identifiability up to relabelling) & 2/2 & yes \\
C06 & Star schema: block rates eta\_kl inside their posterior intervals & Prop. 5 & 2/2 & yes \\
C07 & Rows per customer follow the fitted Poisson mixture (Poisson splitting) & Prop. 6 & 3/3 & yes \\
C08 & Scale rule: resizing the partner keeps rows per anchor; resizing the anchor scales the fact & Prop. 7 (layer 1 scale rule) & 4/4 & yes \\
C09 & Snowflake: segments recovered at every level and carried down the chain & Prop. 8 & 7/7 & yes \\
C10 & Self-referencing tree: multi-type Galton-Watson fit and finiteness; level-conditioned generator & Props. 9, 10, 54 & 5/5 & yes \\
C11 & Self-referencing tree: synthetic hierarchy is valid and matches the source shape & Prop. 10, Sec. 6 & 5/5 & yes \\
C12 & Temporal fact: segment-specific seasonality recovered and reproduced & Prop. 11 & 4/4 & yes \\
C13 & Text: proper predictive distribution, no rare content leaks, class signal kept & Prop. 12, 13 & 6/6 & yes \\
C14 & Enterprise schema end to end: integrity, rules, full scorecard & Sec. 7, Prop. 14 & 4/4 & yes \\
C15 & Rejection sampling for business rules gives the exact conditional law & Prop. 14 & 2/2 & yes \\
C16 & Gamma frailty gives the negative-binomial variance mu + mu\^{}2/alpha & Prop. 6 (degree correction) & 2/2 & yes \\
C17 & Dirichlet-process marginal posterior: mean is (a F0 + n Fn)/(a + n) & Prop. 2 & 1/1 & yes \\
C18 & Same seed gives an identical synthetic database & Sec. 7 (reproducibility) & 1/1 & yes \\
C19 & Doctor: finds planted spec errors from the bundle alone and the fixes improve the result & Sec. 8 (diagnostics loop) & 6/6 & yes \\
C20 & Rule language: three-valued logic, SQL null semantics and string functions & Prop. 15 (three-valued rules) & 9/9 & yes \\
C21 & Structural (spike-and-slab) priors: samplers match the exact posterior & Props. 17-19 (spike-and-slab priors) & 13/13 & yes \\
C22 & Empirical-Bayes slabs: Beta-binomial concentration recovered, shrinkage lowers error & Props. 20-21 (empirical-Bayes slabs) & 4/4 & yes \\
C23 & Dropping duplicate rows of a unique key is exact conditioning on presence & Prop. 24 (unique keys) & 2/2 & yes \\
C24 & Pair model selection: occupancy for grain data, Poisson for count data & Prop. 22 (repeat test) & 4/4 & yes \\
C25 & Rank copula: the fit depends on the order of the data only & Prop. 28 (rank likelihood) & 3/3 & yes \\
C26 & Reporting scenario end to end: grain, keys, aggregates, rules, scorecard & Props. 16, 23, 25, 26 combined (S8); Prop. 52 (structural zeros) & 8/8 & yes \\
C27 & Reporting scenario: omitted and partial filings reproduced & Props. 18, 27 (grouped zeros, nested refinement) & 3/3 & yes \\
C28 & Reporting scenario: open rows repeat, closed items never do & Props. 21, 24 (multiplicity, unique keys) & 3/3 & yes \\
C29 & Grain entity scales with its anchor and stays unique & Prop. 23 (grain entities) & 4/4 & yes \\
C30 & Introspect drafts keys, links and grains from PascalCase warehouse names & Sec. 12 (introspect) & 6/6 & yes \\
C31 & Doctor proposes the missing grain from the bundle alone & Sec. 12 (doctor D03, D21) & 2/2 & yes \\
C32 & Computed columns are filled on the source and stored ones are audited & Prop. 16 (computed columns) & 4/4 & yes \\
C33 & Posterior predictive draws: redraw gives spread that covers the source & Prop. 30 (posterior predictive) & 3/3 & yes \\
C34 & Rare columns are kept out of the copula and the chain converges & Prop. 29 (rare columns) & 3/3 & yes \\
C35 & Count-model choice by Bayes factor: DP only for counts no Poisson/NB mixture explains & Prop. 32 (count-model evidence) & 3/3 & yes \\
C36 & DP counts given a total: the pairwise Gibbs sampler draws the exact conditional & Prop. 33 (conditioned counts) & 3/3 & yes \\
C37 & Functional dependencies: city -> state -> country kept exactly, no extra disclosure & Prop. 36 (dependency groups) & 5/5 & yes \\
C38 & Tiny parent: strata by its own column keep who receives children; hierarchy keeps its size & Props. 37, 38 (attribute strata, conditioned tree) & 4/4 & yes \\
C39 & Small shop end to end: invoices per customer, lines per invoice, lookups, aggregates, precision & Props. 32-35, 39, 40 with Prop. 16 (S9) & 10/10 & yes \\
C40 & Introspect on a CSV shop: lookups, aggregate with a product, anchors, grains, party tables, questions & Sec. 12 (introspect) & 8/8 & yes \\
C41 & Evaluator limits hold their nominal level on small and clustered samples; no placeholder leaks & Sec. 10 (evaluation), Prop. 40 (tail levels) & 5/5 & yes \\
C42 & Doctor proposes a grain only when independence would have produced repeats & Prop. 41 (grain evidence) & 5/5 & yes \\
C43 & Sales warehouse end to end: one header per order, lines 1..n, date gaps, identities, stock, who buys what & Props. 43-45, 48, 49, 53 (headers, offsets, refill, tied segments, running balances, mandatory lines) & 18/18 & yes \\
C44 & Introspect on the sales warehouse (SQLite, declared keys): calendar, headers, offsets, lookups, identities, running balance, first purchase; same spec twice & Prop. 42 (significance of discovered rules) & 10/10 & yes \\
C45 & Same seed gives the same data in a new process; per-row shares keep count dispersion & Prop. 46 (posterior-draw shares of copied rows) & 3/3 & yes \\
C46 & Warehouse catalog regressions: binary columns, composite keys, partly declared links, order header, validation, money TVD, pooled level after reordering & Findings F01-F08, F22-F24 of the AdventureWorks log & 10/10 & yes \\
C47 & Introspect drafts reference tables, keyed and mapped lookups, views, line numbers from 0, date gaps, implied dimensions, ratio columns, cross-table rules and key order & Props. 55-58; findings K01-K16 of the Contoso run & 16/16 & yes \\
C48 & Synthesis keeps views, keyed and mapped lookups, implied dimensions, ratio factors, date gaps, line numbers, key order and cross-table rules exactly, without biasing the marginals & Props. 55-59 & 26/26 & yes \\
C49 & Building blocks: date functions, keyed lookups across date storage, ratio factors at the coarsest precision, computed-column cycles, file formats & Props. 55-58; findings K03, K09, K14, K15 & 8/8 & yes \\
C50 & Introspect on Parquet files without declared keys and with renamed columns: line keys, the view with its renamed columns, the dimension code not asked as a foreign key; the view rebuilt exactly & Props. 55, 60; findings K02, K04, K15 & 7/7 & yes \\
C51 & Links under a rule across two parents: exact draws per row keep the rule on every row, the source's cell margin, and the source's law of each parent within its cell & Props. 59, rake; finding K10 & 4/4 & yes \\
C52 & Conditional inclusion on a fact's anchor: facts that only reference some parent rows (quotas only for sales staff) never reference the excluded rows, and keep their row counts & Prop. anchorincl; AW finding (quotas linked to non-sales staff) & 15/15 & yes \\
C53 & Labels nested in a coarser code (cities in states): drafted as whole-value text within the code; synthetic cities sit in their source states, with no label held by fewer than 3 source rows & Prop. nested; finding K17 & 4/4 & yes \\
C54 & DP row counts under mandatory participation: the base measure 1 + Poisson keeps every parent with a child, with the same mean, in unconditional and total-conditioned draws & Prop. mandatory DP; finding K18 & 5/5 & yes \\
C55 & A small dimension with no label to anchor on anchors itself: the code stays modelled (keeps its dependence on the other columns) and the columns it determines are read by its value & Prop. dimension (self-anchored); finding R6 & 3/3 & yes \\
C56 & Conditional inclusion on an occupancy fact's partner: order lines never reference products that are never sold (null price), even when one segment mixes them with sold products, and keep lines per order & Prop. anchorincl (partner side); AW regression (one reseller line on a product without a price) & 6/6 & yes \\
C57 & One posterior draw per dataset: a dataset is reproducible from its index; across datasets the spread of a share and of rows per parent is posterior plus sampling variance (about twice the sampling variance), and about the sampling variance alone in posterior-mean mode & Prop. postpred (v0.7) & 6/6 & yes \\
C58 & A rule that holds on most rows (an age on a reference date) is found by the interval sweep and synthesised as a mixture: the share that keeps the rule and the offsets of the rest match the source & Prop. soft (v0.7) & 5/5 & yes \\
C59 & A value shared by all rows of a group (a daily price factor per product category) is found by introspect and drawn once per group: every synthetic group has one factor, and the factor law matches & Prop. groupfx (v0.7) & 4/4 & yes \\
C60 & Two links that point at parents in the same area more often than chance (customer and store) are found by introspect and synthesised as a diagonal-inflated mixture: the same-area share matches and the country rule still holds & Prop. affinity (v0.7) & 5/5 & yes \\
C61 & Participation of parents in an entity's link is a hurdle: the share of customers without an order matches the source (half of them, and none when every customer has one) & Prop. hurdle (v0.7) & 9/9 & yes \\
C62 & A latent segment never mixes parent rows that may have children with rows that never do (only sales staff have quotas): the eligibility column stratifies the segments, so the children per eligible parent match the source in every dataset, even with one segment allowed & Prop. strata (v0.7) & 4/4 & yes \\
C63 & A type 2 slowly changing dimension (versions of a product code with valid-from/to dates and a current flag) is found by introspect and synthesised per natural key: version counts match, stable attributes repeat across versions, each version ends the day before the next starts and only the latest is open & Prop. scd2 (v0.7) & 10/10 & yes \\
C64 & An unknown member (key -1, 'Unknown') that facts point at is found by introspect, kept out of the model, copied, and referenced with the source's share & Prop. unknown (v0.7) & 3/3 & yes \\
C65 & A bridge with allocation weights (accounts of a customer, weights summing to 1) is found by introspect; synthetic weights sum to 1 in every group and keep their spread & Prop. bridge (v0.7) & 3/3 & yes \\
C66 & Recipes of the specification: copy every dimension and synthesise only the fact; draw a column from a stated distribution; name the anchor of a role-playing fact by its key column & spec-reference.md (v0.7) & 7/7 & yes \\
C67 & The specification is strict: a misspelt key anywhere is an error with the closest valid key, and every spec drafted or written for the earlier tests still validates & schema.py (v0.7) & 2/2 & yes \\
C68 & The annotated example of the specification reference (every v0.7 option in one spec) validates, runs, and keeps each option's property & spec-reference.md (v0.7) & 8/8 & yes \\
C69 & Optimised kernels compute exactly what they replace & Prop. 'Exact kernels': same IEEE operations in the same order and the same random stream & 8/8 & yes \\
C70 & Keys and links identified from the data, declared keys hidden & Props. 'Inclusion by chance' and 'Randomness bound': no false link; every missed link is asked & 6/6 & yes \\
C71 & Identification statistics: chance inclusion, randomness, forest test & Props. 'Inclusion by chance', 'Randomness bound', 'Chance of a forest' & 11/11 & yes \\
C72 & Exported model: verified, strict JSON, draws are the engine's datasets & Props. 'Exported draw m is the parameter of dataset m' and 'Verification conditions' & 12/12 & yes \\
C73 & Standalone replica: keys, sizes and fidelity & Prop. 'The replica samples the core generative process' & 10/10 & yes \\
C74 & Posterior blocks agree with Stan & Props. 'Block-rate posterior', 'Spike-and-slab Bernoulli', 'Rank likelihood with parameter expansion' & 1/1 & yes \\
C75 & v0.9 engine: pooled weights, exact conditioning, coverage, warnings & Section 'v0.9' propositions & 14/14 & yes \\
C76 & v0.9 scorecard calibration and fidelity statistics & Props. 'ppnull', 'sidak' & 14/14 & yes \\
C77 & v0.9 identification: permutation forest test, chance inclusion & Props. 'forestperm', 'stringchance' & 14/14 & yes \\
C78 & v0.9 input, output and security hardening & Spec reference & 28/28 & yes \\
\bottomrule
\end{longtable}


\begin{thebibliography}{99}\small
\bibitem{nist800226} J.~P.~Near, D.~Darais, N.~Lefkovitz, G.~S.~Howarth (2025). Guidelines for evaluating differential privacy guarantees. NIST SP 800-226. doi:10.6028/NIST.SP.800-226.
\bibitem{ponomareva2023} N.~Ponomareva et al. How to DP-fy ML: a practical guide to machine learning with differential privacy. \emph{J.\ Artif.\ Intell.\ Res.} 77 (2023).
\bibitem{bun2016} M.~Bun, T.~Steinke (2016). Concentrated differential privacy: simplifications, extensions, and lower bounds. \emph{TCC}.
\bibitem{cks2020} C.~L.~Canonne, G.~Kamath, T.~Steinke (2020). The discrete Gaussian for differential privacy. \emph{NeurIPS}.
\bibitem{mironov2017} I.~Mironov (2017). R\'enyi differential privacy. \emph{IEEE CSF}.
\bibitem{dworkroth2014} C.~Dwork, A.~Roth (2014). The algorithmic foundations of differential privacy. \emph{Found.\ Trends Theor.\ Comput.\ Sci.} 9.
\bibitem{ballewang2018} B.~Balle, Y.-X.~Wang (2018). Improving the Gaussian mechanism for differential privacy: analytical calibration and optimal denoising. \emph{ICML}.
\bibitem{dong2022} J.~Dong, A.~Roth, W.~J.~Su (2022). Gaussian differential privacy. \emph{J.\ R.\ Statist.\ Soc.\ B} 84.
\bibitem{rfc8785} A.~Rundgren, B.~Jordan, S.~Erdtman (2020). JSON Canonicalization Scheme (JCS). RFC 8785.
\bibitem{rfc8032} S.~Josefsson, I.~Liusvaara (2017). Edwards-curve digital signature algorithm (EdDSA). RFC 8032.
\bibitem{giomi2023} M.~Giomi, F.~Boenisch, C.~Wehmeyer, B.~Tasn\'adi (2023). A unified framework for quantifying privacy risk in synthetic data. \emph{Proc.\ Privacy Enhancing Technologies} 2023(2).
\bibitem{wilson1927} E.~B.~Wilson (1927). Probable inference, the law of succession, and statistical inference. \emph{J.\ Amer.\ Statist.\ Assoc.} 22.
\bibitem{clopper1934} C.~J.~Clopper, E.~S.~Pearson (1934). The use of confidence or fiducial limits illustrated in the case of the binomial. \emph{Biometrika} 26.
\bibitem{kairouz2015} P.~Kairouz, S.~Oh, P.~Viswanath (2015). The composition theorem for differential privacy. \emph{ICML}.
\bibitem{houssiau2022} F.~Houssiau, J.~Jordon, S.~N.~Cohen, O.~Daniel, A.~Elliott, J.~Geddes, C.~Mole, C.~Rangel-Smith, L.~Szpruch (2022). TAPAS: a toolbox for adversarial privacy auditing of synthetic data. \emph{NeurIPS Workshop on Synthetic Data for Empowering ML Research}.
\bibitem{meeus2023} M.~Meeus, F.~Gu\'epin, A.-M.~Cre\c{t}u, Y.-A.~de Montjoye (2023). Achilles' heels: vulnerable record identification in synthetic data publishing. \emph{ESORICS}.
\bibitem{stadler2022} T.~Stadler, B.~Oprisanu, C.~Troncoso (2022). Synthetic data -- anonymisation groundhog day. \emph{USENIX Security}.
\bibitem{vanbreugel2023} B.~van Breugel, H.~Sun, Z.~Qian, M.~van der Schaar (2023). Membership inference attacks against synthetic data through overfitting detection. \emph{AISTATS}.
\bibitem{carlini2022} N.~Carlini, S.~Chien, M.~Nasr, S.~Song, A.~Terzis, F.~Tram\`er (2022). Membership inference attacks from first principles. \emph{IEEE S\&P}.
\bibitem{edpb2026} European Data Protection Board (2026). Guidelines 02/2026 on anonymisation, version 1.0 for public consultation, adopted 7 July 2026.
\bibitem{ferguson} T.~Ferguson (1973). A Bayesian analysis of some nonparametric problems. \emph{Ann.\ Statist.} 1.
\bibitem{rubin} D.~Rubin (1981). The Bayesian bootstrap. \emph{Ann.\ Statist.} 9.
\bibitem{lkj} D.~Lewandowski, D.~Kurowicka, H.~Joe (2009). Generating random correlation matrices based on vines and extended onion method. \emph{J.\ Multivar.\ Anal.} 100.
\bibitem{hoff} P.~Hoff (2007). Extending the rank likelihood for semiparametric copula estimation. \emph{Ann.\ Appl.\ Stat.} 1.
\bibitem{agg} R.~Arratia, L.~Goldstein, L.~Gordon (1990). Poisson approximation and the Chen--Stein method. \emph{Statist.\ Sci.} 5.
\bibitem{fk} D.~Feldman, D.~Kowal (2025). Bayesian Gaussian copula with auxiliary marginal quantiles. \emph{Bayesian Analysis}.
\bibitem{haario} H.~Haario, E.~Saksman, J.~Tamminen (2001). An adaptive Metropolis algorithm. \emph{Bernoulli} 7.
\bibitem{kemp} C.~Kemp et al.\ (2006). Learning systems of concepts with an infinite relational model. \emph{AAAI}.
\bibitem{kn} B.~Karrer, M.~Newman (2011). Stochastic blockmodels and community structure in networks. \emph{Phys.\ Rev.\ E} 83.
\bibitem{athreya} K.~Athreya, P.~Ney (1972). \emph{Branching Processes}. Springer.
\bibitem{gms} A.~Gelman, X.-L.~Meng, H.~Stern (1996). Posterior predictive assessment of model fitness via realized discrepancies. \emph{Statistica Sinica} 6.
\bibitem{mp} D.~MacKay, L.~Peto (1995). A hierarchical Dirichlet language model. \emph{Natural Language Engineering} 1.
\bibitem{xie} C.~Xie et al.\ (2024). Differentially private synthetic data via foundation model APIs 2: text. \emph{ICML}.
\bibitem{liuwu} J.~Liu, Y.~Wu (1999). Parameter expansion for data augmentation. \emph{JASA} 94.
\bibitem{snoke} J.~Snoke et al.\ (2018). General and specific utility measures for synthetic data. \emph{JRSS A} 181.
\bibitem{kleene} S.~C.~Kleene (1952). \emph{Introduction to Metamathematics}. North-Holland.
\bibitem{kleinman} J.~Kleinman (1973). Proportions with extraneous variance: single and independent samples. \emph{JASA} 68.
\bibitem{bmm} J.~Barnard, R.~McCulloch, X.-L.~Meng (2000). Modeling covariance matrices in terms of standard deviations and correlations. \emph{Statistica Sinica} 10.
\bibitem{crr} M.~Cowles, G.~Roberts, J.~Rosenthal (1999). Possible biases induced by MCMC convergence diagnostics. \emph{J.\ Statist.\ Comput.\ Simul.} 64.
\bibitem{mb} T.~Mitchell, J.~Beauchamp (1988). Bayesian variable selection in linear regression. \emph{JASA} 83.
\bibitem{ospina} R.~Ospina, S.~Ferrari (2010). Inflated beta distributions. \emph{Statistical Papers} 51.
\bibitem{rrr} T.~Raghunathan, J.~Reiter, D.~Rubin (2003). Multiple imputation for statistical disclosure limitation. \emph{J.\ Official Statist.} 19.
\bibitem{rk} J.~Reiter, S.~Kinney (2012). Inferentially valid, partially synthetic data: generating from posterior predictive distributions not necessary. \emph{J.\ Official Statist.} 28.
\bibitem{hrw} J.~Hu, J.~Reiter, Q.~Wang (2018). Dirichlet process mixture models for modeling and generating synthetic versions of nested categorical data. \emph{Bayesian Analysis} 13.
\bibitem{goodman} L.~Goodman (1968). The analysis of cross-classified data: independence, quasi-independence, and interactions in contingency tables with or without missing entries. \emph{JASA} 63.
\bibitem{mullahy} J.~Mullahy (1986). Specification and testing of some modified count data models. \emph{J.\ Econometrics} 33.
\bibitem{shental} N.~Shental, A.~Bar-Hillel, T.~Hertz, D.~Weinshall (2003). Computing Gaussian mixture models with EM using equivalence constraints. \emph{NIPS}.
\bibitem{parciak} M.~Parciak et al.\ (2024). Measuring approximate functional dependencies: a comparative study. \emph{ICDE}.
\bibitem{kimball} R.~Kimball, M.~Ross (2013). \emph{The Data Warehouse Toolkit}, 3rd ed. Wiley.
\bibitem{reiter05} J.~P.~Reiter (2005). Using CART to generate partially synthetic public use microdata. \emph{J.\ Official Statist.} 21.
\bibitem{synthpop} B.~Nowok, G.~M.~Raab, C.~Dibben (2016). synthpop: Bespoke creation of synthetic data in R. \emph{J.\ Statist.\ Software} 74.
\bibitem{sdv} N.~Patki, R.~Wedge, K.~Veeramachaneni (2016). The Synthetic Data Vault. \emph{IEEE DSAA}.
\bibitem{clava} J.~Pang et al.\ (2024). ClavaDDPM: Multi-relational data synthesis with cluster-guided diffusion models. \emph{NeurIPS}.
\bibitem{goodman74} L.~A.~Goodman (1974). Exploratory latent structure analysis using both identifiable and unidentifiable models. \emph{Biometrika} 61.
\bibitem{dunsonxing} D.~B.~Dunson, C.~Xing (2009). Nonparametric Bayes modeling of multivariate categorical data. \emph{J.\ Amer.\ Statist.\ Assoc.} 104.
\bibitem{vonneumann} J.~von Neumann (1951). Various techniques used in connection with random digits. \emph{NBS Appl.\ Math.\ Series} 12.
\bibitem{alaa} A.~Alaa, B.~van Breugel, E.~Saveliev, M.~van der Schaar (2022). How faithful is your synthetic data? Sample-level metrics for evaluating and auditing generative models. \emph{ICML}.
\bibitem{esteban} C.~Esteban, S.~Hyland, G.~R\"atsch (2017). Real-valued (medical) time series generation with recurrent conditional GANs. arXiv:1706.02633.
\bibitem{pang} W.~Pang, M.~Shafieinejad, L.~Liu, S.~Hazlewood, X.~He (2024). ClavaDDPM: multi-relational data synthesis with cluster-guided diffusion models. \emph{NeurIPS}.
\bibitem{syntherela} V.~Hudovernik, M.~Jurkovi\v{c}, E.~\v{S}trumbelj (2026). SyntheRela: a benchmark for synthetic relational database generation. \emph{TMLR}; arXiv:2410.03411.
\bibitem{synthcity} Z.~Qian, B.-C.~Cebere, M.~van der Schaar (2023). Synthcity: a benchmark framework for diverse use cases of tabular synthetic data. \emph{NeurIPS Datasets and Benchmarks}.
\bibitem{vdvw} A.~van der Vaart, J.~Wellner (1996). \emph{Weak Convergence and Empirical Processes}. Springer.
\bibitem{vdv} A.~van der Vaart (1998). \emph{Asymptotic Statistics}. Cambridge University Press.
\bibitem{delbarrio} E.~del Barrio, E.~Gin\'e, C.~Matr\'an (1999). Central limit theorems for the Wasserstein distance between the empirical and the true distributions. \emph{Ann.\ Probab.} 27.
\bibitem{hahn} G.~Hahn, W.~Meeker, L.~Escobar (2017). \emph{Statistical Intervals}, 2nd ed. Wiley.
\bibitem{ruschendorf} L.~R\"uschendorf (2009). On the distributional transform, Sklar's theorem, and the empirical copula process. \emph{J.\ Statist.\ Plann.\ Inference} 139.
\bibitem{lopezpaz} D.~Lopez-Paz, M.~Oquab (2017). Revisiting classifier two-sample tests. \emph{ICLR}.
\bibitem{mironov2012} I.~Mironov. On significance of the least significant bits for differential privacy. \emph{ACM CCS}, 2012.
\bibitem{wilson2020} R.~J.~Wilson, C.~Y.~Zhang, W.~Lam, D.~Desfontaines, D.~Simmons-Marengo and B.~Gipson. Differentially private SQL with bounded user contribution. \emph{PoPETs}, 2020(2), 2020.
\bibitem{desfontaines2022} D.~Desfontaines, J.~Voss, B.~Gipson and C.~Mandayam. Differentially private partition selection. \emph{PoPETs}, 2022(1), 2022.
\end{thebibliography}
\end{document}
